% Options for packages loaded elsewhere
\PassOptionsToPackage{unicode}{hyperref}
\PassOptionsToPackage{hyphens}{url}
\documentclass[
]{article}
\usepackage{xcolor}
\usepackage{amsmath,amssymb}
\setcounter{secnumdepth}{-\maxdimen} % remove section numbering
\usepackage{iftex}
\ifPDFTeX
  \usepackage[T1]{fontenc}
  \usepackage[utf8]{inputenc}
  \usepackage{textcomp} % provide euro and other symbols
\else % if luatex or xetex
  \usepackage{unicode-math} % this also loads fontspec
  \defaultfontfeatures{Scale=MatchLowercase}
  \defaultfontfeatures[\rmfamily]{Ligatures=TeX,Scale=1}
\fi
\usepackage{lmodern}
\ifPDFTeX\else
  % xetex/luatex font selection
\fi
% Use upquote if available, for straight quotes in verbatim environments
\IfFileExists{upquote.sty}{\usepackage{upquote}}{}
\IfFileExists{microtype.sty}{% use microtype if available
  \usepackage[]{microtype}
  \UseMicrotypeSet[protrusion]{basicmath} % disable protrusion for tt fonts
}{}
\makeatletter
\@ifundefined{KOMAClassName}{% if non-KOMA class
  \IfFileExists{parskip.sty}{%
    \usepackage{parskip}
  }{% else
    \setlength{\parindent}{0pt}
    \setlength{\parskip}{6pt plus 2pt minus 1pt}}
}{% if KOMA class
  \KOMAoptions{parskip=half}}
\makeatother
\usepackage{longtable,booktabs,array}
\newcounter{none} % for unnumbered tables
\usepackage{multirow}
\usepackage{calc} % for calculating minipage widths
% Correct order of tables after \paragraph or \subparagraph
\usepackage{etoolbox}
\makeatletter
\patchcmd\longtable{\par}{\if@noskipsec\mbox{}\fi\par}{}{}
\makeatother
% Allow footnotes in longtable head/foot
\IfFileExists{footnotehyper.sty}{\usepackage{footnotehyper}}{\usepackage{footnote}}
\makesavenoteenv{longtable}
\setlength{\emergencystretch}{3em} % prevent overfull lines
\providecommand{\tightlist}{%
  \setlength{\itemsep}{0pt}\setlength{\parskip}{0pt}}
\usepackage{bookmark}
\IfFileExists{xurl.sty}{\usepackage{xurl}}{} % add URL line breaks if available
\urlstyle{same}
\hypersetup{
  pdftitle={OS-9 Operating System},
  hidelinks,
  pdfcreator={LaTeX via pandoc}}

\title{OS-9 Operating System}
\usepackage{etoolbox}
\makeatletter
\providecommand{\subtitle}[1]{% add subtitle to \maketitle
  \apptocmd{\@title}{\par {\large #1 \par}}{}{}
}
\makeatother
\subtitle{System Programmer\textquotesingle s Manual}
\author{}
\date{}

\begin{document}
\maketitle

\section{Introduction}

OS-9 Level One is a versatile multiprogramming/multitasking operating
system for computers utilizing the Motorola 6809 microprocessor,. It is
well-suited for a wide range of applications on 6809 computers of almost
any size or complexity. Its main features are:

\begin{itemize}
\item
  Comprehensive management of all system resources: memory, input/output
  and CPU time.
\item
  A powerful user interface that is easy to learn and use.
\item
  True multiprogramming operation.
\item
  Efficient operation in typical microcomputer configuratjons.
\item
  Expandable, device-independent unified I/O system.
\item
  Full support for modular ROMed software.
\item
  Upward and downward compatibility with OS-9 Level Two.
\end{itemize}

This manual is intended to provide the information necessary to install,
maintain, expand, or write assembly-language software for OS-9 systems.
It assumes that the reader is familiar with the 6809 architecture,
instruction set, and assembly language.

\subsection{History And Design Philosophy}

OS-9 Level One is one of the products of the BASIC09 Advanced 6809
Programming Language development effort undertaken by Microware and
Motorola from 1978 to 1980. During the course of the project it became
evident that a fairly sophisticated operating system would be required
to support BASIC09 and similar high-performance 6809 software.

OS-9\textquotesingle s design was modeled after Bell Telephone
Laboratories\textquotesingle{} UNIX operating system, which is becoming
widely recognized as a standard for mini and micro multiprogramming
operating systems because of its versatility and relatively simple, yet
elegant structure. Even though a ``clone'' of UNIX for the 6809 is
relatively easy to implement, there are a number of problems with this
approach. UNIX was designed for fairly large-scale minicomputers (such
as large PDP-11s) that have high CPU throughput, large fast disk storage
devices and a static I/O environment. Also, UNIX is not particularly
time or disk-storage efficient, especially when used with low-cost disk
drives.

For these reasons, OS-9 was designed to retain the overall concept and
user interface of UNIX, but its implementation is considerably
different. OS-9\textquotesingle s design is tailored to typical
microcomputer performance ranges and operational environments. As an
example, OS-9, unlike UNIX, does not dynamically swap running programs
on and off disk This is because floppy disks and many lower-cost
Winchester-type hard disks are simply too slow to do this efficiently.
Instead, OS-9 always keeps running programs in memory and emphasizes
more efficient use of available ROM or RAM.

OS-9 also introduces some important new features that are intended to
make the most of the capabilities of third-generation microprocessors,
such as support of reentrant, position-independent software that can be
shared by several users simultaneously to reduce overall memory
requirements.

Perhaps the most innovative part of OS-9 is its ``memory module''
management system, which provides extensive support for modular
software, particularly ROMed software. This will play an increasingly
important role in the future as a method of reducing software costs. The
``memory module'' and LINK capabilities of OS-9 permit modules to be
automatically identified, linked together, shared, updated or repaired.
Individual modules in ROM which are defective may be repaired (without
reprogramming the ROM) by placing a ``fixed'' module, with the same
name, but a higher revision number into memory. Memory modules have many
other advantages, for example, OS-9 can allow several programs to share
a common math subroutine module. The same module could automatically be
replaced with a module containing drivers for a hardware arithmetic
processor without any change to the programs which call the module.

Users experienced with UNIX should have little difficulty adapting to
OS-9. Here are some of the main differences between the two systems:

\begin{enumerate}
\def\labelenumi{\arabic{enumi}.}
\item
  OS-9 is written in 6809 assembly language, not C. This improves
  program size and speed characteristics.
\item
  OS-9 was designed for a mixed RAM/ROM microcomputer memory environment
  and more effectively supports reentrant, position-independent code.
\item
  OS-9 introduces the ``memory module'' concept for organizing object
  code with built-in dynamic inter-module linkage.
\item
  OS-9 supports multiple file managers, which are modules that interface
  a class of devices to the file system.
\item
  ``Fork'' and ``Execute'' calls are faster and more memory efficient
  than the UNIX equivalents.
\end{enumerate}

\subsection{System Hardware Requirements}

The OS-9 Operating system consists of building blocks called memory
modules, which are automatically located and linked together when the
system starts up. This makes it extremely easy to reconfigure the
system. For example, reconfiguring the system to handle additional
devices is simply a matter of placing the corresponding modules into
memory. Because OS-9 is so flexible, the minimum hardware requirements
are difficult to define. A bare-bones LEVEL I system requires 4K of ROM
and 2K of RAM, which may be expanded to 56K RAM.

Shown below are the requirements for a typical OS-9 software development
system. Actual hardware requirements may vary depending upon the
particular application.

\begin{itemize}
\item
  6809 MPU
\item
  24K Bytes RAM Memory for Assembly Language Development. 40K Bytes RAM
  Memory for High Level Languages such as BASIC09 (RAM Must Be
  Contiguous From Address Zero Upward)
\item
  4K Bytes of ROM: 2K must be addressed at \$F800 - \$FFFF, the other 2K
  is position-independent and self-locating. Some disk systems may
  require three 2K ROMs.
\item
  Console terminal and interface using serial, parallel, or memory
  mapped video.
\item
  Optional printer using serial or parallel interface.
\item
  Optional real-time clock hardware.
\end{itemize}

I/O device controller addresses can be located anywhere in the memory
space, however it is good practice to place them as high as possible to
maximize RAM expansion capability. Standard OS-9 packages for computers
made by popular manufacturers usually conform to the
system\textquotesingle s customary memory map.

\section{Basic System Organization}

OS-9 is composed of a group of modules, each of which provides specific
functions. When OS-9 is configured for a specific system various modules
are selected to provide a given level of functionality. For example, a
small control computer without a disk does not need the disk-related
OS-9 modules. Most examples in this manual describe a fully-configured
OS-9 system.

\begin{figure}
\centering
\begin{verbatim}
                 +-----------------------+
+----------+     !                       !     +----------+
!          !     !                       !     !          !
!   INIT   ! - - !     OS-9 KERNEL       ! - - !  Clock   !
!          !     !        (ROM)          !     !          !
+----------+     !                       !     +----------+
                 +-----------------------+
                             !
                             !
                 +-----------------------+
                 !                       !
                 ! Input/Output Manager  !
                 !       (IOMAN)         !
                 !                       !
                 +-----------------------+
                    !                 !
                    !                 !
   +--------------------+       +--------------------+
   !                    !       !                    !
   ! Disk File Manager  !       ! Char. File Manager !    More
   !      (RBFMAN)      !       !      (SCFMAN)      ! -> opt.
   !                    !       !                    !
   +--------------------+       +--------------------+
        !          !                !           !
        !          !                !           !
   +--------+  +--------+       +--------+  +--------+
   !        !  !        !       !        !  !        !
   !  Disk  !  !  Disk  !       !  ACIA  !  !  PIA   !    More
   ! Driver !  ! Driver !       ! Driver !  ! Driver ! -> opt.
   !        !  !        !       !        !  !        !
   +--------+  +--------+       +--------+  +--------+
     !     !     !     !          !     !     !     !
     !     !     !     !          !     !     !     !
   +---+ +---+  +---+ +---+     +---+ +---+  +---+ +---+
   !D0 ! !D1 !  !D2 ! !D3 !     !T1 ! !T2 !  !P1 ! !P2 !-> More
   +---+ +---+  +---+ +---+     +---+ +---+  +---+ +---+   opt.
    RBF Device Descriptors        SCF Device Descriptors
\end{verbatim}
\caption{OS-9 Component Module Organization}
\end{figure}

Notice that the diagram on the previous page indicates a multilevel
organization.

The first level is the KERNEL and the CLOCK MODULE. The kernel provide
basic system services such as multitasking, memory management, and links
all other system modules. The CLOCK module is a software handler for the
specific real-time-clock hardware. INIT is an initialization table used
by the kernel during system startup. It specifies initial table sizes,
initial system device names, etc.

The second level is the Input/Output Manager. It provides common
processing all I/O operations. It is required if any OS-supported I/O is
to be performed.

The third level is the File Manager level. File managers perform I/O
request processing for similar classes of I/O devices. The Random Block
File Manager (RBFMAN) processes all disk-type device functions, and the
Sequential Character File Manager (SCFMAN) handles all non-mass storage
devices that basically operate a character at a time, such as terminals
and printers. The user can add additional File Managers to handle
classes of devices not covered by SCFMAN or RBFMAN.

The fourth level is the Device Driver Level. Device drivers handle basic
physical I/O functions for specific I/O controller hardware. Standard
OS-9 systems are typically supplied with a disk driver, a ACIA driver
for terminals and serial printers, and a PIA driver for parallel
printers. Many users add customized drivers of their own design or
purchased from a hardware vendor.

The fifth level is the Device Descriptor Level. These modules are small
tables that are associate specific I/O ports with their logical names,
and the port\textquotesingle s device driver and file manager. They also
contain the physical address of the port and initialization data. By use
of device descriptors, only one copy of each driver is required for each
specific type of I/O controller regardless of how many controllers the
system uses.

One important component not shown is the shell, which is the command
interpreter. It is technically a program and not part of the operating
system itself, and is described fully in the OS-9 Users Manual.

Even though all modules can be resident in ROM, generally only the
KERNEL and INIT modules are ROMed in disk-based systems. All other
modules are loaded into RAM during system startup by a disk bootstrap
module (not shown on diagram) which is also resident in ROM.

\section{Basic Functions of the Kernel}

The nucleus of OS-9 is the ``kernel'', which serves as the system
administrator, supervisor, and resource manager. It is about 3K bytes
long and normally resides in two 2K byte ROMs: ``P1'' residing at
addresses \$F800 - \$FFFF, and ``P2'', which is position-independent. P2
only occupies about half (1K) of the ROM, the other space in the ROM is
reserved for the disk bootstrap module.

The kernel\textquotesingle s main functions are:

\begin{enumerate}
\def\labelenumi{\arabic{enumi}.}
\item
  System initialization after restart.
\item
  Service request processing.
\item
  Memory management.
\item
  MPU management (multiprogramming).
\item
  Basic interrupt processing.
\end{enumerate}

Notice that input/output functions were not included in the list above;
this is because the kernel does not directly process them. The kernel
passes I/O service requests directly to another the Input/Output Manager
(IOMAN) module for processing.

After a hardware reset, the kernel will initialize the system which
involves: locating ROMs in memory, determining the amount of RAM
available, loading any required modules not already in ROM from the
bootstrap device, and running the system startup task (``SYSGO''). The
INIT module is a table used during startup to specify initial table
sizes and system device names.

\subsection{Kernel Service Request Processing}

Service requests (system calls) are used to communicate between OS-9 and
assembly-language-level programs for such things as allocating memory,
creating new processes, etc. System calls use the SWI2 instruction
followed by a constant byte representing the code. Parameters for system
calls are usually passed in MPU registers. In addition to I/O and memory
management functions, there are other service request functions
including interprocess control and timekeeping.

A system-wide assembly languaqe equate file called \texttt{OS9Defs}
defines symbolic names for all service requests. This file is included
when assembling hand-written or compiler-generated code. The OS-9
Assembler has a built-in macro to generate system calls, for example:

\begin{verbatim}
OS9 I$Read
\end{verbatim}

is recongnized and assembled as the equivalent to:

\begin{verbatim}
SWI2
FCB  I$Read
\end{verbatim}

Service requests are divided into two categories:

I/O REQUESTS perform various input/output functions. Requests of this
type are passed by the kernel to IOMAN for processing. The symbolic
names for this category have a ``I\$'' prefix, for example, the ``read''
service request is called \hyperref[i.read]{I\$Read}.

FUNCTION REQUESTS perform memory management, multiprogramming, and
miscellaneous functions. Most are processed by the kernel. The symbolic
names for this category begins with ``F\$''.

\subsection{Kernel Memory Management Functions}

Memory management is an important operating system function. OS-9
manages both the physical assignment of memory to programs \emph{and}
the logical contents of memory, by using entities called ``memory
modules''. All programs are loaded in memory module format, allowing
OS-9 to maintain a directory which contains the name, address, and other
related information about each module in memory. These structures are
the foundation of OS-9\textquotesingle s modular software environment.
Some of its advantages are: automatic run-time ``linking'' of programs
to libraries of utility modules; automatic ``sharing'' of reentrant
programs; replacement of small sections of large programs for update or
correction (even when in ROM); etc.

\subsection{Memory Utilization}

All usable RAM memory must be contiguous from address 0 upward. During
the OS-9 start-up sequence the upper bound of RAM is detemined by an
automatic search, or from the configuration module. Some RAM is reserved
by OS-9 for its own data structures at the top and bottom of memory. The
exact amount depends on the sizes of system tables that are specified in
the configuration module.

All other RAM memory is pooled into a ``free memory'' space. Memory
space is dynamically taken from and returned to this pool as it is
allocated or deallocated for various purposes. The basic unit of memory
allocation is the 256-byte page . Memory is always allocated in whole
numbers of pages.

The data structure used to keep track of memory allocation is a 32-byte
bit-map located at addresses \$0100 - \$011F. Each bit in this table is
associated with a specific page of memory. Bits are cleared to indicate
that the page is free and available for assignment, or set to indicate
that the page is in use or that no RAM memory is present at that
address.

Automatic memory allocation occurs when:

\begin{enumerate}
\def\labelenumi{\arabic{enumi}.}
\item
  Program modules are loaded into RAM.
\item
  Processes are created.
\item
  Processes request additional RAM.
\item
  OS-9 needs I/O buffers, larger tables, etc.
\end{enumerate}

All of the above usually have inverse functions that cause previously
allocated memory to be deallocated and returned to the free memory pool.

In general, memory is allocated for program modules and buffers from
high addresses downward, and for process data areas from lower addresses
upward.

\begin{figure}
\centering
\begin{verbatim}
   TYPICAL MEMORY MAP


+-----------------------+  <- $FFFF
|                       |
|     OS-9 ROMS (4K)    |
|                       |
+-----------------------+  <- $F000
|                       | 
|  I/O DEVICE ADDRESSES |
|                       | 
+-----------------------+  <- $E000
|                       | 
|    SPACE FOR MORE     |
|    OPTIONAL ROMS      |
|                       | 
+-----------------------+  <- END OF RAM MEMORY
|                       | 
|    FILE MANAGERS      |
|  DEVICE DRIVERS, ETC. |
|   (APPROXIMATELY 6K)  |
|                       |
+-----------------------+
|                       |
|      SHELL (1K)       |
|                       |
+-----------------------+
|                       |
| OS-9 DATA STRUCTURES  |
|  (APPROXIMATELY 1K)   |
|                       |
+-----------------------+
|                       |
|    FREE MEMORY FOR    |
|     GENERAL USE       |
|                       |
+-----------------------+ <- $0400
|                       |
| OS-9 DATA STRUCTURES  |
|   AND DIRECT PAGE     |
|                       |
+-----------------------+ <- $0000 BEGINNING OF RAM MEMORY
\end{verbatim}
\caption{}
\end{figure}

The map above is for a ``typical'' system. Actual memory sizes and
addresses may vary depending on the exact system configuration.

\subsection{Overview of Multiprogramming}

OS-9 is a multiprogramming operating system, which allows several
independent programs called ``processes'' can be executed
simultaneously. Each process can have access to any system resource by
issuing appropriate service requests to OS-9. Multiprogramming functions
use a hardware real-time clock that generates interrupts at a regular
rate of about 10 times per second. MPU time is therefore divided into
periods typically 100 milliseconds in duration. This basic time unit is
called a tick . Processes that are ``active'' (meaning not waiting for
some event) are run for a specific system-assigned period called a
``time slice''. The duration of the time slice depends on a
process\textquotesingle s priority value relative to the priority of all
other active processes. Many OS-9 service requests are available to
create, terminate, and control processes.

\subsection{Process Creation}

New processes are created when an existing process executes a
\hyperref[f.fork]{F\$Fork} service request. Its main argument is the
name of the program module (called the ``primary module'') that the new
process is to initially execute. OS-9 first attempts to find the module
in the ``module directory'', which includes the names of all program
modules already present in memory. If the module cannot be found there.
OS-9 usually attempts to load into memory a mass-storage file using the
requested module name as a file name.

Once the module has been located, a data structure called a ``process
descriptor'' is assigned to the new process. The process descriptor is a
64-byte package that contains information about the process, its state,
memory allocations, priority, queue pointers, etc. The process
descriptor is automatically initialized and maintained by OS-9. The
process itself has no need, and is not permitted to access the
descriptor.

The next step in the creation of a new process is allocation of data
storage (RAM) memory for the process. The primary
module\textquotesingle s header contains a storage size value that is
used unless the \hyperref[f.fork]{F\$Fork} system call requested an
optionally larger size. OS-9 then attempts to allocate a CONTIGUOUS
memory area of this size from the free memory space.

If any of the previous steps cannot be performed, creation of the new
process is aborted, and the process that originated the
\hyperref[f.fork]{F\$Fork} is informed of the error. Otherwise, the new
process is added to the active process queue for execution scheduling.

The new process is also assigned a unique number called a ``process ID''
which is used as its identifier. Other processes can commnunciate with
it by referring to its ID in various system calls. The process also has
associated with it a ``user ID'' which is used to identify all processes
and files belonging to a particular user. The user ID is inherited from
the parent process.

Processes terminate when they execute an \hyperref[f.exit]{F\$Exit}
system service request, or when they receive fatal signals. The process
termination closes any open paths, deallocates its memory, and unlinks
its primary module.

\subsection{Process States}

At any instant, a process can be in one of three states:

ACTIVE - The process is active and ready for execution.

WAITING - The process is suspended until a child process terminates or a
signal is received.

SLEEPING - The process is suspended for a specific period of time or
until a signal is received.

There is a queue for each process state. The queue is a linked list of
the ``process descriptors'' of processes in the corresponding state.
State changes are performed by moving a process descriptor to another
queue.

\subsubsection{The Active State}

This state includes all ``runnable'' processes, which are given time
slices for execution according to their relative priority with respect
to all other active processes. The scheduler uses a pseudo-round-robin
scheme that gives all active processes some CPU time, even if they have
a very low relative priority.

\subsubsection{The Wait State}

This state is entered when a process executes a
\hyperref[f.wait]{F\$Wait} system service request. The process remains
suspended until the death of any of its descendant processes, or, until
it receives a signal.

\subsubsection{The Sleeping State}

This state is entered when a process executes a
\hyperref[f.sleep]{F\$Sleep} service request, which specifies a time
interval. (a specific number of ticks) for which the process is to
remain suspended. The process remains asleep until the specified time
has elapsed, or until a signal is received.

\subsection{Execution Scheduling}

The kernel contains a scheduler that is responsible for allocation of
CPU time to active processes. OS-9 uses a scheduling algorithm that
ensures all processes get some execution time.

All active processes are members of the active process queue, which is
kept sorted by process ``age''. Age is a count of how many process
switches have occurred since the process\textquotesingle{} last time
slice. When a process is moved to the active process queue from another
queue, its ``age'' is initialized by setting it to the
process\textquotesingle{} assigned priority, i.e., processes having
relatively higher priority are placed in the queue with an artificially
higher age. Also, whenever a new process is activated, the ages of all
other processes are incremented.

Upon conclusion of the currently executing process\textquotesingle{}
time slice, the scheduler selects the process having the highest age to
be executed next. Because the queue is kept sorted by age, this process
will be at the bead of the queue. At this time the ages of all other
active processes are incremented (ages are never incremented beyond
255).

An exception is newly-active processes that were previously deactivated
while they were in the system state. These processes are noted and given
higher priority than others because they are usually executing critical
routines that affect shared system resources and therefore could be
blocking other unrelated processes.

When there are no active processes, the kernel will set itself up to
handle the next interrupt and then execute a CWAI instruction, which
decreases interrupt latency time.

\subsection{Signals}

``Signals'' are an asynchronous control mechanism used for interprocess
communication and control. A signal behaves like a software interrupt in
that it can cause a process to suspend a program, execute a specific
routine, and afterward return to the interrupted program. Signals can be
sent from one process to another process (by means of the SEND service
request), or they can be sent from OS-9 system routines to a process.

Status information can be conveyed by the signal in the form of a
one-byte numeric value. Some of the signal ``codes'' (values) have
predefined meanings, but all the rest are user-defined. The defined
signal codes are:

0 = KILL (non-interceptable process abort)

1 = WAKEUP - wake up sleeping process

2 = KEYBOARD ABORT

3 = KEYBOARD INTERRUPT

4 - 255 USER DEFINED

When a signal is sent to a process, the signal is noted and saved in the
process descriptor. If the process is in the sleeping or waiting state,
it is changed to the active state. It then becomes eligible for
execution according to the usual MPU scheduler criteria. When it gets
its next time slice, the signal is processed.

What happens next depends on whether or not the process had previously
set up a ``signal trap'' (signal service routine) by executing an
\hyperref[f.icpt]{F\$ICPT} service request. If it had not, the process
is immediately aborted. It is also aborted if the signal code is zero.
The abort will be deferred if the process is in system mode: the process
dies upon its return to user state.

If a signal intercept trap has been set up, the process resumes
execution at the address given in the \hyperref[f.icpt]{F\$ICPT} service
request. The signal code is passed to this routine, which should
terminate with an RTI instruction to resume normal execution of the
process.

NOTE: ``Wakeup'' signals activate a sleeping process: they \emph{do not}
vector through the intercept routine.

If a process has a signal pending (usually because it has not been
assigned a time slice since the signal was received), and some other
process attempts to send it another signal, the new signal is aborted
and the ``send'' service request will return an error status. The sender
should then execute a sleep service request for a few ticks before
attempting to resend the signal, so the destination process has an
opportunity to process the previously pending signal.

\subsection{Interrupt Processing}

Interrupt processing is another important function of the kernel. All
hardware interrupts are vectored to specific processing routines. IRQ
interrupts are handled by a prioritized polling system (actually part of
IOMAN) which automatically identifies the source of the interrupt and
dispatches to the associated user or system defined service routine. The
real-time clock will generate IRQ interrupts. SWI, SWI2, and SWI3
interrupts are vectored to user-definable addresses which are ``local''
to each procedure, except that SWI2 is normally used for OS-9 service
requests calls. The NMI and FIRQ interrupts are not normally used and
are vectored through a RAM address to an RTI instruction.

\subsubsection{Physical Interrupt Processing}

The OS-9 kernel. ROMs contain the hardware vectors required by the 6809
MPU at addresses \$FFF0 through \$FFFF. These vectors each point to
jump-extended-indirect instruction which vector the MPU to the actual
interrupt service routine. A RAM vector table in page zero of memory
contains the target addresses of the jump instructions as follows:

{\def\LTcaptype{none} % do not increment counter
\begin{longtable}[]{@{}
  >{\raggedright\arraybackslash}p{(\linewidth - 2\tabcolsep) * \real{0.5000}}
  >{\raggedright\arraybackslash}p{(\linewidth - 2\tabcolsep) * \real{0.5000}}@{}}
\toprule\noalign{}
\begin{minipage}[b]{\linewidth}\raggedright
INTERRUPT
\end{minipage} & \begin{minipage}[b]{\linewidth}\raggedright
ADDRESS
\end{minipage} \\
\midrule\noalign{}
\endhead
\bottomrule\noalign{}
\endlastfoot
SWI3 & \$002C \\
SWI2 & \$002E \\
FIRQ & \$0030 \\
IRQ & \$0032 \\
SWI & \$0034 \\
NMI & \$0036 \\
\end{longtable}
}

OS-9 initializes each of these locations after reset to point to a
specific service routine in the kernel. The SWI, SWI2, and SWI3 vectors
point to specific routines which in turn read the corresponding pseudo
vector from the process\textquotesingle{} process descriptor and
dispatch to it. This is why the \hyperref[f.sswi]{F\$SSWI} service
request to be local to a process since it only changes a pseudo vector
in the process descriptor. The IRQ routine points directly to the IRQ
polling system, or to it indirectly via the real-time clock device
service routine. The FIRQ and NMI vectors are not normally used by OS-9
and point to RTI instructions.

A secondary vector table located at \$FFE0 contains the addresses of the
routines that the RAM vectors are initialized to. They may be used when
it is necessary to restore the original service routines after altering
the RAM vectors. On the next page are the definitions of both the actual
hardware interrupt vector table, and the secondary vector table:

{\def\LTcaptype{none} % do not increment counter
\begin{longtable}[]{@{}
  >{\raggedright\arraybackslash}p{(\linewidth - 4\tabcolsep) * \real{0.2200}}
  >{\raggedright\arraybackslash}p{(\linewidth - 4\tabcolsep) * \real{0.2200}}
  >{\raggedright\arraybackslash}p{(\linewidth - 4\tabcolsep) * \real{0.5600}}@{}}
\toprule\noalign{}
\begin{minipage}[b]{\linewidth}\raggedright
VECTOR
\end{minipage} & \begin{minipage}[b]{\linewidth}\raggedright
ADDRESS
\end{minipage} & \begin{minipage}[b]{\linewidth}\raggedright
\end{minipage} \\
\midrule\noalign{}
\endhead
\bottomrule\noalign{}
\endlastfoot
Secondary Vector Table & & \\
TICK & \$FFE0 & Clock Tick Service Routine \\
SWI3 & \$FFE2 & \\
SWI2 & \$FFE4 & \\
FIRQ & \$FFE6 & \\
IRQ & \$FFE8 & \\
SWI & \$FFEA & \\
NMI & \$FFEC & \\
WARM & \$FFEE & Reserved for warm-start \\
Hardware Vector Table & & \\
SWI3 & \$FFF2 & \\
SWI2 & \$FFF4 & \\
FIRQ & \$FFF6 & \\
IRQ & \$FFF8 & \\
SWI & \$FFFA & \\
NMI & \$FFFC & \\
RESTART & \$FFFE & \\
\end{longtable}
}

If it is necessary to alter the RAM vectors use the secondary vector
table to exit the substitute routine. The technique of altering the IRQ
pointer is usually used by the clock service routines to reduce latency
time of this frequent interrupt source.

\subsubsection{Logical Interrupt Polling System}

In OS-9 systems, most I/O devices use IRQ-type interrupts, so OS-9
includes a sophisticated polling system that automatically identifies
the source of the interrupt and dispatches to its associated
user-defined service routine. The information required for IRQ polling
is maintained in a data structure called the ``IRQ polling table''. The
table has a 9-byte entry for each possible IRQ-generatinq device. The
table size is static and defined by an initialization constant in the
System Configuration Module.

The polling system is prioritized so devices having a relatively greater
importance (i.e., interrupt frequency) are polled before those of lesser
priority. This is accomplished by keeping the entries sorted by
priority, which is a number between 0 (lowest) and 255 (highest). Each
entry in the table has 6 variables:

\begin{enumerate}
\def\labelenumi{\arabic{enumi}.}
\item
  POLLING ADDRESS: The address of the device\textquotesingle s status
  register, which must have a bit or bits that indicate it is the source
  of an interrupt.
\item
  MASK BYTE; This byte selects one or more bits within the device status
  register that are interrupt request flag(s). A set bit identifies the
  active bit(s).
\item
  FLIP BYTE: This byte selects whether the bits in the device status
  register are true when set or true when cleared. Cleared bits indicate
  active when set.
\item
  SERVICE ROUTINE ADDRESS: The user-supplied address of the
  device\textquotesingle s interrupt service routine.
\item
  STATIC STORAGE ADDRESS: a user-supplied ter to the permanent storage
  required by the device service routine.
\item
  PRIORITY; The device priority number: 0 to 255. This value determines
  the order in which the devices in the polling table will be polled.
  Note: this is not the same as a process priority which is used by the
  execution scheduler to decide which process gets the next time slice
  for MPU execution.
\end{enumerate}

When an IRQ interrupt occurs, the polling system is entered via the
corresponding RAM interrupt vector. It starts polling the devices, using
the entries in the polling table in priority order. For each entry, the
status register address is loaded into accumulator A using the device
address from the table. An exclusive-or operation using the flip-byte is
executed, followed by a logical-and operation using the mask byte. If
the result is non-zero, the device is assumed to be the cause of the
interrupt.

The device\textquotesingle s static storage address and service routine
address is read from the table and executed.

The interrupt service routine should terminate with an an \emph{RTS},
not an RTI instruction.

Entries can be made to the IRQ polling table by use of a special OS-9
service request called \hyperref[f.irq]{F\$IRQ}. This is a priviledged
service request that can be executed only when OS-9 is in System Mode
(which is the case when device drivers are executed).

The actual code for the interrupt polling system is located in the IOMAN
module. The kernel P1 and P2 modules contain the physical interrupt
processing routines.

\section{Memory Modules}

Any object to be loaded into the memory of an OS-9 system must use the
memory module format and conventions. The memory module concept allows
OS-9 to manage the logical contents as well as the physical contents of
memory. The basic idea is that all programs are individual, named
objects.

The operating system keeps track of modules which are in memory at all
times by use of a ``module directory'' . It contains the addresses and a
count of how many processes are using each module. When modules are
loaded into memory, they are added to the directory. When they are no
longer needed, their memory is deallocated and their name removed from
the directory (except ROMs, which are discussed later). In many
respects, modules and memory in general, are managed just like a disk.
In fact, the disk and memory management sections of OS-9 share many
subroutines.

Each module has three parts; a module header, module body and a
cyclic-redundancy-check (CRC) value. The header contains information
that describes the module and its use. This information includes: the
modules size, its type (machine language. BASIC09 compiled code, etc);
attributes (executable, reentrant, etc), data storage memory
requirements, execution starting address, etc. The CRC value is used to
verify the integrity of a module.

There are several different kinds of modules, each type having a
different usage and function. Modules do not have to be complete
programs, or even 6809 machine language. They may contain BASIC09
``I-code'', constants, single subroutines, subroutine packages, etc. The
main requirements are that modules do not modify themselves arid that
they be position-independent so OS-9 can load or relocate them wherever
memory space is available. In this respect, the module format is the
OS-9 equivalent of ``load records'' used in older-style operating
systems.

\subsection{Memory Module Structure}

At the beginning (lowest address) of the module is the module header,
which can have several forms depending on the module\textquotesingle s
usage. OS-9 family software such as BASIC09, Pascal, C, the assembler,
and many utility programs automatically generate modules and headers.
Following the header is the program/constant section which is usually
pure code. The module name string is included somewhere in this area.
The last three bytes of the module are a three-byte Cyclic Redundancy
Check (CRC) value used to verify the integrity of the module.

\begin{longtable}[]{@{}
  >{\raggedright\arraybackslash}p{(\linewidth - 0\tabcolsep) * \real{1.0000}}@{}}
\caption{Module Format}\tabularnewline
\toprule\noalign{}
\endfirsthead
\endhead
\bottomrule\noalign{}
\endlastfoot
MODULE HEADER \\
PROGRAM OR CONSTANTS \\
CRC \\
\end{longtable}

The 24-bit CRC is performed over the entire module from the first byte
of the module header to the byte just before the CRC itself. The CRC
polynomial used is \$800063. (See \hyperref[f.crc]{F\$CRC})

Because most OS-9 family software (such as the assembler) automatically
generate the module header and CRC values, the programner usually does
not have to be concerned with writing routines to generate them.

\subsection{Module Header Definitions}

The first nine bytes of all module headers are identical:

{\def\LTcaptype{none} % do not increment counter
\begin{longtable}[]{@{}
  >{\raggedleft\arraybackslash}p{(\linewidth - 2\tabcolsep) * \real{0.1837}}
  >{\raggedright\arraybackslash}p{(\linewidth - 2\tabcolsep) * \real{0.8163}}@{}}
\toprule\noalign{}
\begin{minipage}[b]{\linewidth}\raggedleft
MODULE OFFSET
\end{minipage} & \begin{minipage}[b]{\linewidth}\raggedright
DESCRIPTION
\end{minipage} \\
\midrule\noalign{}
\endhead
\bottomrule\noalign{}
\endlastfoot
\$0,\$1 = & Sync Bytes (\$87,\$CD). These two constant bytes are used to
locate modules. \\
\$2,\$3 = & Module Size. The overall size of the module in bytes
(includes CRC). \\
\$4,\$5 = & Offset to Module Name. The address of the module name string
relative to the start (first sync byte) of the module. The name string
can be located anywhere in the module and consists of a string of ASCII
characters having the sign bit set on the last character. \\
\$6 = & Module Type/Language Type. See text. \\
\$7 = & Attributes/Revision Level. See text. \\
\$8 = & Header Check. The one\textquotesingle s compliment of (the
vertical parity (exclusive OR) of) the previous eight bytes \\
\end{longtable}
}

\subsubsection{Type/Language Byte}

The module type is coded into the tour most significant bits of byte 6
of the module header. Eight types are pre-defined by convention, some of
which are for OS-9\textquotesingle s internal use only. The type codes
are:

{\def\LTcaptype{none} % do not increment counter
\begin{longtable}[]{@{}
  >{\raggedleft\arraybackslash}p{(\linewidth - 2\tabcolsep) * \real{0.2308}}
  >{\raggedright\arraybackslash}p{(\linewidth - 2\tabcolsep) * \real{0.7692}}@{}}
\toprule\noalign{}
\endhead
\bottomrule\noalign{}
\endlastfoot
\$1 & Program module \\
\$2 & Subroutine module \\
\$3 & Multi-module (for future use) \\
\$4 & Data module \\
\$5-\$B & User-definable \\
\$C & OS-9 System module \\
\$D & OS-9 File Manager module \\
\$E & OS-9 Device Driver module \\
\$F & OS-9 Device Descriptor module \\
\end{longtable}
}

NOTE: 0 is not a legal type code.

``user-defined'' types having type codes of 0 through 9. They have six
more bytes in their headers defined as follows:

{\def\LTcaptype{none} % do not increment counter
\begin{longtable}[]{@{}
  >{\raggedleft\arraybackslash}p{(\linewidth - 2\tabcolsep) * \real{0.1837}}
  >{\raggedright\arraybackslash}p{(\linewidth - 2\tabcolsep) * \real{0.8163}}@{}}
\toprule\noalign{}
\begin{minipage}[b]{\linewidth}\raggedleft
MODULE OFFSET
\end{minipage} & \begin{minipage}[b]{\linewidth}\raggedright
DESCRIPTION
\end{minipage} \\
\midrule\noalign{}
\endhead
\bottomrule\noalign{}
\endlastfoot
\$9,\$A = & Execution Offset. The program or
subroutine\textquotesingle s starting address, relative to the first
byte of the sync code. Modules having multiple entry points (cold start,
warm start, etc.) may have a branch table starting at this address. \\
\$B,\$C = & Permanent Storage Requirement. This is the minimum number of
bytes of data storage required to run. This is the number used by
\hyperref[f.fork]{F\$Fork} and \hyperref[f.chain]{F\$Chain} to allocate
a process\textquotesingle{} data area.

If the module will not be directly executed by a
\hyperref[f.chain]{F\$Chain} or \hyperref[f.fork]{F\$Fork} service
request (for instance a subroutine package), this entry is not used by
OS-9. It is commonly used to specify the maximum stack size required by
reentrant subroutine modules. The calling program can check this value
to determine if the subroutine has enough stack space. \\
\end{longtable}
}

\subsection{Executable Memory Module Format}

\begin{figure}
\centering
\begin{verbatim}
Relative            Usage                    Check Range
Address

        +------------------------------+  ---+--------+---
 $00    |                              |     |        |
        +--    Sync Bytes ($87CD)    --+     |        |
 $01    |                              |     |        |
        +------------------------------+     |        |
 $02    |                              |     |        |
        +--   Module Size (bytes)    --+     |        |
 $03    |                              |     |        |
        +------------------------------+     |        |
 $04    |                              |     |        |
        +--   Module Name Offset     --+   header     |
 $05    |                              |   parity     |
        +------------------------------+     |        |
 $06    |     Type     |   Language    |     |        |
        +------------------------------+     |        |
 $07    |  Attributes  |   Revision    |     |        |
        +------------------------------+  ---+--    module
 $08    |     Header Parity Check      |             CRC
        +------------------------------+              |
 $09    |                              |              |
        +--     Execution Offset     --+              |
 $0A    |                              |              |
        +------------------------------+              |
 $0B    |                              |              |
        +--  Permanent Storage Size  --+              |
 $0C    |                              |              |
        +------------------------------+              |
 $0D    |                              |              |
        |  (Add'l optional header      |              |
        |   extensions located here    |              |
        |                              |              |
        |  .  .  .  .  .  .  .  .  .   |              |
        |                              |              |
        |                              |              |
        |       Module Body            |              |
        | object code, constants, etc. |              |
        |                              |              |
        |                              |              |
        +------------------------------+              |
        |                              |              |
        +--                          --+              |
        |       CRC Check Value        |              |
        +--                          --+              |
        |                              |              |
        +------------------------------+  ------------+---
\end{verbatim}
\caption{}
\end{figure}

\subsection{ROMed Memory Modules}

When OS-9 starts after a system reset, it searches the entire memory
space for ROMed modules. It detects them by looking for the module
header sync code (\$87,\$CD) which are unused 6809 opcodes. When this
byte pattern is detected, the header check is performed to verify a
correct header. If this test succeeds, the module size is obtained from
the header and a 24-bit CRC is performed over the entire module. If the
CRC matches correctly, the module is considered valid, and it is entered
into the module directory. The chances of detecting a ``false module''
are virtually nil.

In this manner all ROMed modules present in the system at startup are
automatically included in the system module directory. Some of the
modules found initially are various parts of OS-9: file managers, device
driver, the configuration module, etc.

After the module search OS-9 links to whichever of its component modules
that it found. This is the secret of OS-9\textquotesingle s
extraordinary adaptability to almost any 6809 computer; it automatically
locates its required and optional component modules, wherever they are,
and rebuilds the system each time that it is started.

ROMs containing non-system modules are also searched so any
user-supplied software is located during the start-up process and
entered into the module directory.

\section{The OS-9 Unified Input/Output System}

OS-9 has a unified I/O system that provides system-wide
hardware-independent I/O services for user programs and OS-9 itself. All
I/O service requests (system call) are received by the kernel and passed
to the Input/Output Manager (IOMAN) module for processing IOMAN performs
some processing (such as allocating data structures for the I/O path)
and calls the file managers and device drivers to do much of the actual
work. File manager, device driver, and device descriptor modules are
standard memory modules that can be loaded into memory from files and
used while the system is running.

The structural organization of I/O-related modules in an OS-9 system is
hierarchical, as illustrated below:

\begin{figure}
\centering
\begin{verbatim}
               +-----------------------+
               !                       !
               ! Input/Output Manager  !
               !       (IOMAN)         !
               !                       !
               +-----------------------+
                  !                 !
                  !                 !
 +--------------------+       +--------------------+
 !                    !       !                    !
 ! Disk File Manager  !       ! Char. File Manager !    More
 !      (RBFMAN)      !       !      (SCFMAN)      ! -> opt.
 !                    !       !                    !
 +--------------------+       +--------------------+
      !          !                !           !
      !          !                !           !
 +--------+  +--------+       +--------+  +--------+
 !        !  !        !       !        !  !        !
 !  Disk  !  !  Disk  !       !  ACIA  !  !  PIA   !    More
 ! Driver !  ! Driver !       ! Driver !  ! Driver ! -> opt.
 !        !  !        !       !        !  !        !
 +--------+  +--------+       +--------+  +--------+
   !     !     !     !          !     !     !     !
   !     !     !     !          !     !     !     !
+---+ +---+  +---+ +---+     +---+ +---+  +---+ +---+
!D0 ! !D1 !  !D2 ! !D3 !     !T1 ! !T2 !  !P1 ! !P2 !-> More
+---+ +---+  +---+ +---+     +---+ +---+  +---+ +---+   opt.
 RBF Device Descriptors        SCF Device Descriptors
\end{verbatim}
\caption{}
\end{figure}

\subsection{The Input/Output Manager (IOMAN)}

The Input/output Manager (IOMAN) module provides the first level of
service for I/O system calls by routing data on I/O paths from processes
to/from the appropriate file managers and device drivers. It maintains
two important internal OS-9 data structures: the device table and the
path table. This module is used in all OS-9 Level One systems and should
never be modified.

When a path is opened, IOMAN attempts to link to a memory module having
the device name given (or implied) in the pathlist. This module is the
device\textquotesingle s descriptor, which contains the names of the
device driver and file manager for the device. This information is saved
by IOMAN so subsequent system calls can be routed to these modules.

\subsection{File Managers}

OS-9 systems can have any number of File Manager modules. The function
of a file manager is to process the raw data stream to or from device
drivers for a similar class of devices to conform to the OS-9 standard
I/O and file structure, removing as many unique device operational
characteristics as possible from I/O operations. They are also
responsible for mass storage allocation and directory processing if
applicable to the class of devices they service.

File managers usually buffer the data stream and issue requests to the
kernel for dynamic allocation of buffer memory. They may also monitor
and process the data stream, for example, adding line feed characters
after carriage return cbaracters.

The file managers are reentrant and one file manager may be used for an
entire class of devices having similar operational ctiaracteristics. The
two standard OS-9 file managers are:

RBFMAN: The Random Block File Manager which operates random-access,
block-structured devices such as disk systems, bubble memories, etc.

SCFMAN: Sequential Character File Manager which is used with
single-character-oriented devices such as CRT or hardcopy terminals,
printers, modems etc.

\subsection{Device Driver Modules}

The device driver modules are subroutine packages that perform basic,
low-level I/O transfers to or from a specific type of I/O device
hardware controller. These modules are reentrant so one copy of the
module can simultaneously run several different devices which use
identical I/O ccntrollers. For example the device driver for 6850 serial
interfaces is called ``ACIA'' and can communicate to any number of
serial terminals.

Device driver modules use a standard module header and are given a
module type of ``device driver'' (code \$E).The execution offset address
in the module header points to a branch table that has a minimum of six
(three-byte) entries.Each entry is typically a LBRA to the corresponding
subroutine. The File Managers call specific routines in the device
driver through this table, passing a pointer to a path decriptor and the
hardware control register address in the MPU registers. The branch table
looks like:

+0~=~Device~Initialization~Routine\\
+3~=~Read~From~Device\\
+6~=~Write~to~Device\\
+9~=~Get~Device~Status\\
+\$C~=~Set~Device~Status\\
+\$F~=~Device~Termination~Routine\\

For a complete description of the parameters passed to these subroutines
see the file manager descriptions. Also see the appendicies on writing
device drivers.

\subsection{Device Descriptor Modules}

Device descriptor modules are small, non-executable modules that provide
information that associates a specific I/O device with its logical name,
hardware controller address(es), device driver name, file manager name,
and initialization paramaters.

Recall that device drivers and file managers both operate on general
classes of devices, not specific I/O ports. The device descriptor
modules tailor their functions to a specific I/O device. One device
descriptor module must exist for each I/O device in the system.

The name of the module is the name the device is known by to the system
and user (i.e. it is the device name given in pathlists\}. Its format
consists of a standard module header that has a type ``device
descriptor'' (code \$F). The rest of the device descriptor header
consists of:

{\def\LTcaptype{none} % do not increment counter
\begin{longtable}[]{@{}
  >{\raggedleft\arraybackslash}p{(\linewidth - 2\tabcolsep) * \real{0.1935}}
  >{\raggedright\arraybackslash}p{(\linewidth - 2\tabcolsep) * \real{0.8065}}@{}}
\toprule\noalign{}
\endhead
\bottomrule\noalign{}
\endlastfoot
\$9,\$A = & File manager name string relative address. \\
\$B,\$C = & Device driver name string relative address \\
\$D = & Mode/Capabilities. (D S PE PW PR E W R) \\
\$E,\$F,\$10 = & Device controller absolute physical (24-bit) address \\
\$11 = & Number of bytes ( ``n'' bytes in intialization table) \\
\$12,\$12+n = & Initialization table \\
\end{longtable}
}

The initialization table is copied into the ``option section'' of the
path descriptor when a path to the device is opened. The values in this
table may be used to define the operating parameters that are changeable
by the OS9 \hyperref[i.getstt]{I\$GetStt} and
\hyperref[i.setstt]{I\$SetStt} service requests. For example, a
terminal\textquotesingle s initialization parameters define which
control characters are used for backspace, delete, etc. The maximum size
of initialization table which may be used is 32 bytes. If the table is
less than 32 bytes long, the remaining values in the path descriptor
will be set to zero.

You may wish to add additional devices to your system. If a similar
device controller already exists, all you need to do is add the new
hardware and load another device descriptor. Device descriptors can be
in ROM or loaded into RAM from mass-storage files while the system is
running.

The diagram on the next page illustrates the device descriptor module
format.

\begin{figure}
\centering
\begin{verbatim}
MODULE     DEVICE DESCRIPTOR MODULE FORMAT
OFFSET

           +-----------------------------+  ---+--------+---
  $0       |                             |     |        |
           +--   Sync Bytes ($87CD)    --+     |        |
  $1       |                             |     |        |
           +-----------------------------+     |        |
  $2       |                             |     |        |
           +--   Module Size (bytes)   --+     |        |
  $3       |                             |     |        |
           +-----------------------------+     |        |
  $4       |                             |     |        |
           +-- Offset to Module Name   --+   header     |
  $5       |                             |   parity     |
           +-----------------------------+     |        |
  $6       | $F (TYPE)   |  $1 (LANG)    |     |        |
           +-----------------------------+     |        |
  $7       | Attributes  |   Revision    |     |        |
           +-----------------------------+  ---+--    module
  $8       |  Header Parity Check        |             CRC
           +-----------------------------+              |
  $9       |                             |              |
           +--  Offset to File Manager --+              |
  $A       |         Name String         |              |
           +-----------------------------+              |
  $B       |                             |              |
           +-- Offset to Device Driver --+              |
  $C       |         Name String         |              |
           +-----------------------------+              |
  $D       |        Mode Byte            |              |
           +-----------------------------+              |
  $E       |                             |              |
           +--    Device Controller    --+              |
  $F       | Absolute Physical Address   |              |
           +--       (24 bit)          --+              |
 $10       |                             |              |
           +-----------------------------+              |
 $11       |  Initialization Table Size  |              |
           +-----------------------------+              |
$12,$12+N  |                             |              |
           |  (Initialization Table)     |              |
           |                             |              |
           +-----------------------------+              |
           |    (Name Strings etc)       |              |
           +-----------------------------+              |
           |      CRC Check Value        |              |
           +-----------------------------+  ------------+---
\end{verbatim}
\caption{}
\end{figure}

\subsection{Path Descriptors}

Every open path is represented by a data structure called a path
descriptor (``PD''). It contains the information required by the file
managers and device drivers to perform I/O functions. Path descriptors
are exactly 64 bytes long and are dynamically allocated and deallocated
by IOMAN as paths are opened and closed.

PDs are INTERNAL data structures that are not normally referenced from
user or applications programs. In fact, it is almost impossible to
locate a path\textquotesingle s PD when OS-9 is in user mode. The
description of PDs is mostly of interest to, and presented here for
those programmers who need to write custom file managers, device
drivers, or other extensions to OS-9.

PDs have three sections: the first 10-byte section is defined
universally for all file managers and device drivers, as shown below.

\begin{longtable}[]{@{}
  >{\raggedright\arraybackslash}p{(\linewidth - 6\tabcolsep) * \real{0.1250}}
  >{\raggedright\arraybackslash}p{(\linewidth - 6\tabcolsep) * \real{0.1250}}
  >{\raggedleft\arraybackslash}p{(\linewidth - 6\tabcolsep) * \real{0.1250}}
  >{\raggedright\arraybackslash}p{(\linewidth - 6\tabcolsep) * \real{0.6250}}@{}}
\caption{Universal Path Descriptor Definitions}\tabularnewline
\toprule\noalign{}
\begin{minipage}[b]{\linewidth}\raggedright
Name
\end{minipage} & \begin{minipage}[b]{\linewidth}\raggedright
Addr
\end{minipage} & \begin{minipage}[b]{\linewidth}\raggedleft
Size
\end{minipage} & \begin{minipage}[b]{\linewidth}\raggedright
Description
\end{minipage} \\
\midrule\noalign{}
\endfirsthead
\toprule\noalign{}
\begin{minipage}[b]{\linewidth}\raggedright
Name
\end{minipage} & \begin{minipage}[b]{\linewidth}\raggedright
Addr
\end{minipage} & \begin{minipage}[b]{\linewidth}\raggedleft
Size
\end{minipage} & \begin{minipage}[b]{\linewidth}\raggedright
Description
\end{minipage} \\
\midrule\noalign{}
\endhead
\bottomrule\noalign{}
\endlastfoot
PD.PD & \$00 & 1 & Path number \\
PD.MOD & \$01 & 1 & Access mode: 1=read 2=write 3=update \\
PD.CNT & \$02 & 1 & Number of paths using this PD \\
PD.DEV & \$03 & 2 & Address of associated device table entry \\
PD.CPR & \$05 & 1 & Requester\textquotesingle s process ID \\
PD.RGS & \$06 & 2 & Caller\textquotesingle s MPU register stack
address \\
PD.BUF & \$08 & 2 & Address of 236-byte data buffer (if used) \\
PD.FST & \$0A & 22 & Defined by file manager \\
PD.OPT & \$20 & 32 & Reserved for GETSTAT/SETSTAT options \\
\end{longtable}

The 22-byte section called ``PD.FST'' is reserved for and defined by
each type of file manager for file pointers, permanent variables, etc.

The 32-byte section called ``PD.OPT'' is used as an ``option'' area for
dynamically-alterable operating parameters for the file or device. These
variables are initialized at the time the path is opened by copying the
initialization table contained in the device descriptor module, and can
be altered later by user programs by means of the
\hyperref[i.getstt]{I\$GetStt} and \hyperref[i.setstt]{I\$SetStt} system
calls.

These two sections are defined each file manager\textquotesingle s in
the assembly language equate file (\texttt{SCFDefs} for SCFMAN and
\texttt{RBFDefs} for RBFMAN).

\section{Random Block File Manager}

The Random Block File Manager (RBFMAN) is a file manager module that
supports random access block-oriented mass storage devices such as disk
systems, bubble memory systems, and high-performance tape systems.
RBFMAN can handle any number or type of such systems simultaneously. It
is a reentrant subroutine package called by IOMAN for I/O service
requests to random-access devices. It is responsible for maintaining the
logical and physical file structures.

In the course of normal operation, RBFMAN requests allocation and
deallocation of 256-byte data buffers; usually one is required for each
open file. When physical I/O functions are necessary, RBFMAN directly
calls the subroutines in the associated device drivers. All data
transfers are performed using 256-byte data blocks. RBFMAN does not
directly deal with physical addresses such as tracks, cylinders, etc.
Instead, it passes to device driver modules address parameters using a
standard address called a ``logical sector number'', or ``LSN''. LSNs
are integers in the range of 0 to n-1, where n is the maximum number of
sectors on the media. The driver is responsible for translating the
logical sector number to actual cylinder/track/sector values.

Because RBFMAN is designed to support a wide range of devices having
different performance and storage capacity, it is highly
parameter-driven. The physical parameters it uses are stored on the
media itself. On disk systems, this information is written on the first
few sectors of track number zero. The device drivers also use this
information, particularly the physical parameters stored on sector 0.
These parameters are written by the ``format'' program that initializes
and tests the media.

\subsection{Logical and Physical Disk Organization}

All mass storage volumes (disk media) used by OS-9 utilize the first few
sectors of the volume to store basic identification structure, and
storage allocation information.

Logical sector zero (LSN 0) is called the \emph{Identification Sector}
which contains description of the physical and logical format of the
volume.

Logical sector one (LSN 1) contains an allocation map which indicates
which disk sectors are free and available for use in new or expanded
files.

The volume\textquotesingle s root directory usually starts at logical
sector two.

\subsubsection{Identification Sector}\label{ident-sector}

Logical sector number zero contains a description of the physical and
logical characteristics of the volume. These are established by the
\texttt{format} command program when the media is initialized, the table
below gives the OS-9 mnemonic name, byte address, size, and description
of each value stored in this sector.

{\def\LTcaptype{none} % do not increment counter
\begin{longtable}[]{@{}
  >{\raggedright\arraybackslash}p{(\linewidth - 6\tabcolsep) * \real{0.2400}}
  >{\raggedright\arraybackslash}p{(\linewidth - 6\tabcolsep) * \real{0.1200}}
  >{\raggedleft\arraybackslash}p{(\linewidth - 6\tabcolsep) * \real{0.1200}}
  >{\raggedright\arraybackslash}p{(\linewidth - 6\tabcolsep) * \real{0.5200}}@{}}
\toprule\noalign{}
\begin{minipage}[b]{\linewidth}\raggedright
Name
\end{minipage} & \begin{minipage}[b]{\linewidth}\raggedright
Addr
\end{minipage} & \begin{minipage}[b]{\linewidth}\raggedleft
Size
\end{minipage} & \begin{minipage}[b]{\linewidth}\raggedright
Description
\end{minipage} \\
\midrule\noalign{}
\endhead
\bottomrule\noalign{}
\endlastfoot
DD.TOT & \$00 & 3 & Total number of sectors on media \\
DD.TKS & \$03 & 1 & Number of sectors per track \\
DD.MAP & \$04 & 2 & Number of bytes in allocation map \\
DD.BIT & \$06 & 2 & Number of sectors per cluster \\
DD.DIR & \$08 & 3 & Starting sector of root directory \\
DD.OWN & \$0B & 2 & Owner\textquotesingle s user number \\
DD.ATT & \$0D & 1 & Disk attributes \\
DD.DSK & \$05 & 2 & Disk identification (for internal use) \\
DD.FMT & \$10 & 1 & Disk format: density, number of sides \\
DD.SPT & \$11 & 2 & Number of sectors per track \\
DD.RES & \$13 & 2 & Reserved for future use \\
DD.BT & \$15 & 3 & Starting sector of bootstrap file \\
DD.BSZ & \$18 & 2 & Size of bootstrap file (in bytes) \\
DD.DAT & \$1A & 5 & Time of creation: Y:M:D:H:M \\
DD.NAM & \$1F & 32 & Volume name: last char has sign bit set \\
DD.OPT & \$3F & 32 & Option area \\
\end{longtable}
}

\subsubsection{Disk Allocation Map Sector}\label{lsn1}

One sector (usually LSN 1) of the disk is used for the ``disk allocation
map'' that specifies which clusters on the disk are available for
allocation of file storage space The address of this sector is always
assigned logical sector 1 by the format program DD.MAP specifies the
number of bytes in this sector which are actually used in the map.

Each bit in the map corresponds to a cluster of sectors on the disk. The
number of sectors per cluster is specified by the ``DD.BIT'' variable in
the identification sector, and is always an integral power of two, i,e.,
1, 2, 4, 8, 16, etc. There are a maximum of 4096 bits in the map, so
media such as double-density double-sided floppy disks and hard disks
will use a cluster size of two or more sectors. Each bit is cleared if
the corresponding cluster is available for allocation, or set if the
sector is already allocated, non-existent, or physically defective. The
bitmap is initially created by the \texttt{format} utility program.

\subsubsection{File Descriptor Sectors}

The first sector of every file is called a ``file descriptor'', which
contains the logical and physical description of the file.. The table
below describes the contents of the descriptor.

{\def\LTcaptype{none} % do not increment counter
\begin{longtable}[]{@{}
  >{\raggedright\arraybackslash}p{(\linewidth - 6\tabcolsep) * \real{0.1600}}
  >{\raggedright\arraybackslash}p{(\linewidth - 6\tabcolsep) * \real{0.1200}}
  >{\raggedleft\arraybackslash}p{(\linewidth - 6\tabcolsep) * \real{0.1200}}
  >{\raggedright\arraybackslash}p{(\linewidth - 6\tabcolsep) * \real{0.6000}}@{}}
\toprule\noalign{}
\begin{minipage}[b]{\linewidth}\raggedright
Name
\end{minipage} & \begin{minipage}[b]{\linewidth}\raggedright
Addr
\end{minipage} & \begin{minipage}[b]{\linewidth}\raggedleft
Size
\end{minipage} & \begin{minipage}[b]{\linewidth}\raggedright
Description
\end{minipage} \\
\midrule\noalign{}
\endhead
\bottomrule\noalign{}
\endlastfoot
FD.ATT & \$0 & 1 & File Attributes: D S PE PW PR E W R \\
FD.OWN & \$1 & 2 & Owner\textquotesingle s User ID \\
FD.DAT & \$3 & 5 & Date Last Modified; Y M D H M \\
FD.LNK & \$8 & 1 & Link Count \\
FD.SIZ & \$9 & 4 & File Size (number of bytes) \\
FD.DCR & \$D & 3 & Date Created: Y M D \\
FD.SEG & \$10 & 240 & Segment List: see below \\
\end{longtable}
}

The attribute byte contains the file permission bits. Bit 7 is set to
indicate a directory file, bit 6 indicates a ``sharable'' file, bit 5 is
public execute, bit 4 is public write, etc.

The segment list consists of up to 48 five-byte entries that have the
size and address of each block of storage that comprise the file in
logical order. Each entry has a three-byte logical sector number of the
block, and a two-byte block size (in sectors). The entry following the
last segment will be zero.

When a file is created, it initially has no data segments allocated to
it. Write operations past the current end-of-file (the first write is
always past the end-of-file) cause additional sectors to be allocated to
the file. If the file has no segments, it is given an initial segment
having the number of sectors specified by the minimum allocation entry
in the device descriptor, or the number of sectors requested if greater
than the minimum. Subsequent expansions of the file are also generally
made in minimum allocation increments. An attempt is made to expand the
last segment wherever possible rather than adding a new segment. When
the file is closed, unused sectors in the last segment are truncated.

A note about disk allocation: OS-9 attempts to minimize the number of
storage segments used in a file. In fact, many files will only have one
segment in which case no extra read operations are needed to randomly
access any byte on the file. Files can have multiple segments if the
free space of the disk becomes very fragmented, or if a file is
repeatedly closed, then opened and expanded at some later time. This can
be avoided by writing a byte at the highest address to be used on a file
before writing any other data.

\subsubsection{Directory Files}

Disk directories are files that have the ``D'' attribute set. Directory
files contain an integral number of directory entries each of which can
bold the name and LSN of a single regular or directory file.

Each directory entry is 32 bytes long, consisting of 29 bytes for the
file name followed by a three byte logical sector number of the
file\textquotesingle s descriptor sector. The file name is
left-justified in the field with the sign bit of the last character set.
Unused entries have a zero byte in the first file name character
position.

Every mass-storage media must have a master directory called the ``root
directory''. The beginning logical sector number of this directory is
stored in the identification sector, as previously described.

\subsection{RBFMAN Definitions of the Path Descriptor.}

The table below describes the usage of the file-manager-reserved section
of path descriptors used by RBFMAN.

{\def\LTcaptype{none} % do not increment counter
\begin{longtable}[]{@{}
  >{\raggedright\arraybackslash}p{(\linewidth - 6\tabcolsep) * \real{0.2000}}
  >{\raggedright\arraybackslash}p{(\linewidth - 6\tabcolsep) * \real{0.1200}}
  >{\raggedleft\arraybackslash}p{(\linewidth - 6\tabcolsep) * \real{0.1200}}
  >{\raggedright\arraybackslash}p{(\linewidth - 6\tabcolsep) * \real{0.5600}}@{}}
\toprule\noalign{}
\begin{minipage}[b]{\linewidth}\raggedright
Name
\end{minipage} & \begin{minipage}[b]{\linewidth}\raggedright
Addr
\end{minipage} & \begin{minipage}[b]{\linewidth}\raggedleft
Size
\end{minipage} & \begin{minipage}[b]{\linewidth}\raggedright
Description
\end{minipage} \\
\midrule\noalign{}
\endhead
\bottomrule\noalign{}
\endlastfoot
Universal Section (same for all file managers) & & & \\
PD.PD & \$00 & 1 & Path number \\
PD.MOD & \$01 & 1 & Mode (read/write/update) \\
PD.CNT & \$02 & 1 & Number of open images \\
PD.DEV & \$03 & 2 & Address of device table entry \\
PD.CPR & \$05 & 1 & Current process ID \\
PD.RGS & \$06 & 2 & Address of callers register stack \\
PD.BUF & \$08 & 2 & Buffer address \\
RBFMAN Path Descriptor Definitions & & & \\
PD.SMF & \$0A & 1 & State flags (see next page) \\
PD.CP & \$0B & 4 & Current logical file position (byte addr) \\
PD.SIZ & \$0F & 4 & File size \\
PD.SBL & \$13 & 3 & Segment beginning logical sector number \\
PD.SBP & \$16 & 3 & Segment beginning physical sector number \\
PD.SSZ & \$19 & 2 & Segment size \\
PD.DSK & \$15 & 2 & Disk ID (for internal use only) \\
PD.DTB & \$lD & 2 & Address of drive table \\
RBFMAN Option Section Definitions (Copied from device descriptor) & &
& \\
& \$20 & 1 & Device class 0= SCF 1=NSF 2=PIPE 3=SBF \\
PD.DRV & \$21 & 1 & Drive number (0..N) \\
PD.STP & \$22 & 1 & Step rate \\
PD.TYV & \$23 & 1 & Device type \\
PD.UNS & \$24 & 1 & Density capability \\
PD.CYL & \$25 & 2 & Number of cylinders (tracks) \\
PD.SID & \$27 & 1 & Number of sides (surfaces) \\
PD.VFY & \$28 & 1 & 0 = verify disk writes \\
PD.SCT & \$29 & 2 & Default number of sectors/track \\
PD.T0S & \$2B & 2 & Default number of sectors/track (track 0) \\
PD.ILV & \$2D & 1 & Sector interleave factor \\
PD.SAS & \$2E & 1 & Segment allocation size \\
(the following values are \emph{not} copied from the device descriptor)
& & & \\
PD.ATT & \$33 & 1 & File attributes (D S PE PW PR E W R) \\
PD.FD & \$34 & 3 & File descriptor PSN (physical sector \#) \\
PD.DFD & \$37 & 3 & Directory file descriptor PSN \\
PD.DCP & \$3A & 4 & File\textquotesingle s directory entry pointer \\
PD.DVT & \$3E & 2 & Address of device table entry \\
\end{longtable}
}

State Flag (PD.SMF): the bits of this byte are defined as:

bit~0~=~set~if~current~buffer~has~been~altered\\
bit~1~=~set~if~current~sector~is~in~buffer\\
bit~2~=~set~if~descriptor~sector~in~buffer\\

The first section of the path descriptor is universal for all file
managers, the second and third sections are defined by RBFMAN and
RBFMAN-type device drivers. The option section of the path descriptor
contains many device operating parameters which may be read and/or
written by the OS9 \hyperref[i.getstt]{I\$GetStt} and
\hyperref[i.setstt]{I\$SetStt} service requests. This section is
initialized by IOMAN which copies the initialization table of the device
descriptor into the option section of the path descriptor when a path to
a device is opened. Any values not determined by this table will default
to zero.

\subsection{RBF Device Descriptor Modules}

This section describes the definitions and use of the initialization
table contained in device descriptor modules for RBF-type devices.

{\def\LTcaptype{none} % do not increment counter
\begin{longtable}[]{@{}
  >{\raggedright\arraybackslash}p{(\linewidth - 6\tabcolsep) * \real{0.1562}}
  >{\raggedright\arraybackslash}p{(\linewidth - 6\tabcolsep) * \real{0.1094}}
  >{\raggedright\arraybackslash}p{(\linewidth - 6\tabcolsep) * \real{0.1094}}
  >{\raggedright\arraybackslash}p{(\linewidth - 6\tabcolsep) * \real{0.6250}}@{}}
\toprule\noalign{}
\begin{minipage}[b]{\linewidth}\raggedright
Module Offset
\end{minipage} & \begin{minipage}[b]{\linewidth}\raggedright
\end{minipage} & \begin{minipage}[b]{\linewidth}\raggedright
\end{minipage} & \begin{minipage}[b]{\linewidth}\raggedright
\end{minipage} \\
\midrule\noalign{}
\endhead
\bottomrule\noalign{}
\endlastfoot
0-\$11 & & & Standard Device Descriptor Module Header \\
\$12 & IT.DTP & RMB 1 & DEVICE TYPE (0=SCF 1=RBF 2=PIPE 3=SBF) \\
\$13 & IT.DRV & RMB 1 & DRIVE NUMBER \\
\$14 & IT.STP & RMB 1 & STEP RATE \\
\$15 & IT.TYP & RMB 1 & DEVICE TYPE (See RBFMAN path descriptor) \\
\$16 & IT.DNS & RMB 1 & MEDIA DENSITY (0 - SINGLE, 1-DOUBLE) \\
\$17 & IT.CYL & RMB 2 & NUMBER OF CYLINDERS (TRACKS) \\
\$19 & IT.SID & RMB 1 & NUMBER OF SURFACES (SIDES) \\
\$1A & IT.VFY & RMB 1 & 0 = VERIFY DISK WRITES \\
\$1B & IT.SCT & RMB 2 & Default Sectors/Track \\
\$1D & IT.T0S & RMB 2 & Default Sectors/Track (Track 0) \\
\$1F & IT.ILV & RMB 1 & SECTOR INTERLEAVE FACTOR \\
\$20 & IT.SAS & RMB 1 & SEGMENT ALLOCATION SIZE \\
\end{longtable}
}

IT.DRV - This location is used to associate a one byte integer with each
drive that a controller will handle. The drives for each controller
should be numbered 0 to n-1, where n is the maximum number of drives the
controller can handle.

IT.STP - (Floppy disks) This location sets the head stepping rate that
will be used with a drive. The step rate should be set to the fastest
value that the drive is capable of to reduce access time. The actual
values stored depended on the specific disk controller and disk driver
module used. Below are the values which are used by the popular Western
Digital floppy disk controller IC:

{\def\LTcaptype{none} % do not increment counter
\begin{longtable}[]{@{}
  >{\centering\arraybackslash}p{(\linewidth - 8\tabcolsep) * \real{0.1724}}
  >{\raggedleft\arraybackslash}p{(\linewidth - 8\tabcolsep) * \real{0.2069}}
  >{\raggedleft\arraybackslash}p{(\linewidth - 8\tabcolsep) * \real{0.2069}}
  >{\raggedleft\arraybackslash}p{(\linewidth - 8\tabcolsep) * \real{0.2069}}
  >{\raggedleft\arraybackslash}p{(\linewidth - 8\tabcolsep) * \real{0.2069}}@{}}
\toprule\noalign{}
\begin{minipage}[b]{\linewidth}\centering
Step Code
\end{minipage} & \begin{minipage}[b]{\linewidth}\centering
FD1771
\end{minipage} & \begin{minipage}[b]{\linewidth}\centering
FD179X Family
\end{minipage} & \begin{minipage}[b]{\linewidth}\raggedleft
\end{minipage} & \begin{minipage}[b]{\linewidth}\raggedleft
\end{minipage} \\
\midrule\noalign{}
\endhead
\bottomrule\noalign{}
\endlastfoot
& 5" & 8" & 5" & 8" \\
0 & 40ms & 20ms & 30ms & 15ms \\
1 & 20ms & 10ms & 20ms & 10ms \\
2 & 12ms & 6ms & 12ms & 6ms \\
3 & 12ms & 6ms & 6ms & 3ms \\
\end{longtable}
}

IT.TYP~-~Device~type~(All~types)\\
\strut ~~~~bit~0~-\/-~0~=~5"~floppy~disk\\
\strut ~~~~~~~~~~~~~1~=~8"~floppy~disk\\
\strut ~~~~bit~6~-\/-~0~=~Standard~OS-9~format\\
\strut ~~~~~~~~~~~~~1~=~Non-standard~format\\
\strut ~~~~bit~7~-\/-~0~=~Floppy~disk\\
\strut ~~~~~~~~~~~~~1~=~Hard~disk\\
\strut ~\\
IT.DNS~-~Density~capabilities~(Floppy~disk~only)\\
\strut ~~~~bit~0~-\/-~0~=~Single~bit~density~(FM)\\
\strut ~~~~~~~~~~~~~1~=~Double~bit~density~(MFM)\\
\strut ~\\
\strut ~~~~bit~1~-\/-~0~=~Single~track~density~(5",~48~TPI)\\
\strut ~~~~~~~~~~~~~1~=~Double~track~density~(5",~96~TPI)\\

IT.SAS - This value specifies the minimum number of sectors to be
allocated at any one time.

\subsection{RBF-type Device Drivers}

An RBF type device driver module contains a package of subroutines that
perform sector oriented I/O to or from a specific hardware controller.
These modules are usually reentrant so that one copy of the module can
simultaneously run several different devices that use identical I/O
controllers. IOMAN will allocate a static storage area for each device
(which may control several drives). The size of the storage area is
given in the device driver module header. Some of this storage area will
be used by IOMAN and RBFMAN, the device driver is free to use the
remainder in any manner. This static storage is used as follows:

\begin{longtable}[]{@{}
  >{\raggedright\arraybackslash}p{(\linewidth - 6\tabcolsep) * \real{0.1200}}
  >{\raggedright\arraybackslash}p{(\linewidth - 6\tabcolsep) * \real{0.1400}}
  >{\raggedright\arraybackslash}p{(\linewidth - 6\tabcolsep) * \real{0.2400}}
  >{\raggedright\arraybackslash}p{(\linewidth - 6\tabcolsep) * \real{0.5000}}@{}}
\caption{Static Storage Definitions}\tabularnewline
\toprule\noalign{}
\begin{minipage}[b]{\linewidth}\raggedright
Offset
\end{minipage} & \begin{minipage}[b]{\linewidth}\raggedright
ORG 0
\end{minipage} & \begin{minipage}[b]{\linewidth}\raggedright
\end{minipage} & \begin{minipage}[b]{\linewidth}\raggedright
\end{minipage} \\
\midrule\noalign{}
\endfirsthead
\toprule\noalign{}
\begin{minipage}[b]{\linewidth}\raggedright
Offset
\end{minipage} & \begin{minipage}[b]{\linewidth}\raggedright
ORG 0
\end{minipage} & \begin{minipage}[b]{\linewidth}\raggedright
\end{minipage} & \begin{minipage}[b]{\linewidth}\raggedright
\end{minipage} \\
\midrule\noalign{}
\endhead
\bottomrule\noalign{}
\endlastfoot
0 & V.PAGE & RMB 1 & PORT EXTENDED ADDRESS (A20 - A16) \\
1 & V.PORT & RMB 2 & DEVICE BASE ADDRESS \\
3 & V.LPRC & RMB 1 & LAST ACTIVE PROCESS ID \\
4 & V.BUSY & RMB 1 & ACTIVE PROCESS ID (0 = NOT BUSY) \\
5 & V.WAKE & RMB 1 & PROCESS ID TO REAWAKEN \\
& V.USER & EQU . & END OF OS9 DEFINITIONS \\
6 & V.NDRV & RMB 1 & NUMBER OF DRIVES \\
& DRVBEG & EQU . & BEGINNING OF DRIVE TABLES \\
7 & TABLES & RMB DRVMEM*N & RESERVE N DRIVE TABLES \\
& & RMB 1 & \\
& FREE & EQU & FREE FOR DRIVER TO USE \\
\end{longtable}

NOTE: V.PAGE through V.USER are predefined in the \texttt{OS9Defs} file.
V.NDRV, DRVBEG, DRVMEM are predefined in the \texttt{RBFDefs} file.

V.PAGE, V.PORT These three bytes are defined by IOMAN as the 24-bit
device address.

V.LPRC This location contains the process ID of the last process to use
the device. Not used by RBF-type device drivers.

V.BUSY This location contains the process ID of the process currently
using the device. Defined by RBFMAN.

V.WAKE This location contains the process-ID of any process that is
waiting for the device to complete I/O (0 = NO PROCESS WAITING). Defined
by device driver.

V.NDRV This location contains the number of drives that the controller
can use. Defined by the device driver as the maximum number of drives
that the controller can work with. RBFMAN will assume that there is a
drive table for each drive. Also see the driver INIT routine in this
section.

TABLES This area contains one table for each drive that the controller
will handle (RBFMAN will assume that there are as many tables as
indicated by V.NDRV). Some time after the driver INIT routine has been
called, RBFMAN will issue a request for the driver to read the
identification sector (logical sector zero) from a drive. At this time
the driver will initialize the corresponding drive table by copying the
first part of the identification sector (up to DD.SIZ) into it, Also see
the ``\hyperref[ident-sector]{Identification Sector}'' section of this
manual. The format of each drive table is as given below:

{\def\LTcaptype{none} % do not increment counter
\begin{longtable}[]{@{}
  >{\raggedright\arraybackslash}p{(\linewidth - 6\tabcolsep) * \real{0.1200}}
  >{\raggedright\arraybackslash}p{(\linewidth - 6\tabcolsep) * \real{0.1400}}
  >{\raggedright\arraybackslash}p{(\linewidth - 6\tabcolsep) * \real{0.2400}}
  >{\raggedright\arraybackslash}p{(\linewidth - 6\tabcolsep) * \real{0.5000}}@{}}
\toprule\noalign{}
\begin{minipage}[b]{\linewidth}\raggedright
Offset
\end{minipage} & \begin{minipage}[b]{\linewidth}\raggedright
ORG 0
\end{minipage} & \begin{minipage}[b]{\linewidth}\raggedright
\end{minipage} & \begin{minipage}[b]{\linewidth}\raggedright
\end{minipage} \\
\midrule\noalign{}
\endhead
\bottomrule\noalign{}
\endlastfoot
\$00 & DD.TOT & RMB 3 & Total number of sectors on media \\
\$03 & DD.TKS & RMB 1 & Number of sectors per track \\
\$04 & DD.MAP & RMB 2 & Number of bytes in allocation map \\
\$06 & DD.BIT & RMB 2 & Number of sectors per cluster \\
\$08 & DD.DIR & RMB 3 & Starting sector of root directory \\
\$0B & DD.OWN & RMB 2 & Owner\textquotesingle s user number \\
\$0D & DD.ATT & RMB 1 & Disk attributes \\
\$05 & DD.DSK & RMB 2 & Disk identification \\
\$10 & DD.FMT & RMB 1 & Disk format: density, number of sides \\
\$11 & DD.SPT & RMB 2 & Number of sectors per track \\
\$13 & DD.RES & RMB 2 & Reserved for future use \\
& DD.SIZ & EQU . & \\
\$15 & V.TRAK & RMB 2 & Current Track Number \\
\$17 & V.BMB & RMB 1 & Bit-map Use Flag \\
\$18 & DRVMEM & EQU . & Size of Each Drive Table \\
\end{longtable}
}

DD.TOT This location contains the total number of sectors contained on
the disk.

DD.TKS This location contains the track size (in sectors).

DD.MAP This location contains the number of bytes in the disk allocation
bit map.

DD.BIT This location contains the number of sectors that each bit
represents in the disk allocation bit map.

DD.DIR This location contains the logical sector number of the disk root
directory.

DD.OWN This location contains the disk owner\textquotesingle s user
number.

DD.APT This location contains the disk access permission attributes as
defined below:

BIT~7~-~D~(DIRECTORY~IF~SET)\\
BIT~6~-~S~(SHARABLE~IF~SET)\\
BIT~5~-~PX~(PUBLIC~EXECUTE~IF~SET)\\
BIT~4~-~PW~(PUBLIC~WRITE~IF~SET)\\
BIT~3~-~PR~(PUBLIC~READ~IF~SET)\\
BIT~2~-~X~(EXECUTE~IF~SET)\\
BIT~1~-~W~(WRITE~IF~SET).\\
BIT~0~-~R~(READ~IF~SET)\\

DD.DSK This location contains a pseudo random number which is used to
identify a disk so that OS-9 may detect when a disk is removed from the
drive and another inserted in its place.

DD.FMT DISK FORMAT:

BIT~B0~-~SIDE\\
\strut ~~~~0~=~SINGLE~SIDED\\
\strut ~~~~1~=~DOUBLE~SIDED\\
\strut ~\\
BIT~B1~-~DENSITY\\
\strut ~~~~0~=~SINGLE~DENSITY\\
\strut ~~~~1~=~DOUBLE~DENSITY\\
\strut ~\\
BIT~B2~-~TRACK~DENSITY\\
\strut ~~~~0~=~SINGLE~(48~TPI)\\
\strut ~~~~1=~DOUBLE~(96~TPI)\\

DD.SPT Number of sectors per track (track zero may use a different
value, specified by IT.T0S in the device descriptor).

DD.RES RESERVED FOR FUTURE USE

V.TRAK This location contains the current track which the head is on and
is updated by the driver.

V.BMB This location is used by RBFMAN to indicate whether or not the
disk allocation bit map is currently in use (0 = not in use). The disk
driver routines must not alter this location.

\subsection{RBFMAN Device Drivers}

As with all device drivers, RBFMAN-type device drivers use a standard
executable memory module format with a module type of ``device driver''
(CODE \$E). The execution offset address in the module header points to
a branch table that has six three byte entries. Each entry is typically
a LBRA to the corresponding subroutine. The branch table is defined as
follows:

{\def\LTcaptype{none} % do not increment counter
\begin{longtable}[]{@{}
  >{\raggedright\arraybackslash}p{(\linewidth - 6\tabcolsep) * \real{0.2000}}
  >{\raggedright\arraybackslash}p{(\linewidth - 6\tabcolsep) * \real{0.2000}}
  >{\raggedright\arraybackslash}p{(\linewidth - 6\tabcolsep) * \real{0.2000}}
  >{\raggedright\arraybackslash}p{(\linewidth - 6\tabcolsep) * \real{0.4000}}@{}}
\toprule\noalign{}
\endhead
\bottomrule\noalign{}
\endlastfoot
\multirow{6}{=}{ENTRY} & LBRA & INIT & INITIALIZE DRIVE \\
& LBRA & READ & READ SECTOR \\
& LBRA & WRITE & WRITE SECTOR \\
& LBRA & GETSTA & GET STATUS \\
& LBRA & SETSTA & SET STATUS \\
& LBRA & TERM & TERMINATE DEVICE \\
\end{longtable}
}

Each subroutine should exit with the condition code register C bit
cleared if no error occurred. Otherwise the C bit should be set and an
appropriate error code returned in the B register. Below is a
description of each subroutine, its input parameters, and its output
parameters.

\subsubsection{NAME: INIT}

{\def\LTcaptype{none} % do not increment counter
\begin{longtable}[]{@{}
  >{\raggedright\arraybackslash}p{(\linewidth - 2\tabcolsep) * \real{0.3182}}
  >{\raggedright\arraybackslash}p{(\linewidth - 2\tabcolsep) * \real{0.6818}}@{}}
\toprule\noalign{}
\endhead
\bottomrule\noalign{}
\endlastfoot
NAME: & INIT \\
INPUT: & \begin{minipage}[t]{\linewidth}\raggedright
(U)~=~ADDRESS~OF~DEVICE~STATIC~STORAGE\\
(Y)~=~ADDRESS~OF~THE~DEVICE~DESCRIPTOR~MODULE
\end{minipage} \\
OUTPUT: & NONE \\
ERROR OUTPUT: & \begin{minipage}[t]{\linewidth}\raggedright
(CC)~=~C~bit~set.\\
(B)~=~Appropriate~error~code.
\end{minipage} \\
FUNCTION: & INITIALIZE DEVICE AND ITS STATIC STORAGE AREA \\
\end{longtable}
}

\begin{enumerate}
\def\labelenumi{\arabic{enumi}.}
\item
  If disk writes are verified, use the \hyperref[f.srqmem]{F\$SRqMem}
  service request to allocate a 256 byte buffer area where a sector may
  be read back and verified after a write.
\item
  Initialize the device permanent storage. For floppy disk controller
  typically this consists of initializing V.NDRV to the number of drives
  that the controller will work with, initializing DD.TOT in the drive
  table to a non-zero value so that sector zero may be read or written
  to, and initializing V.TRAK to \$FF so that the first seek will find
  track zero.
\item
  Place the IRQ service routine on the IRQ polling list by using the OS9
  \hyperref[f.irq]{F\$IRQ} service request.
\item
  Initialize the device control registers (enable interrupts if
  necessary).
\end{enumerate}

NOTE: Prior to being called, the device permanent storage will be
cleared (set to zero) except for V.PAGE and V.PORT which will contain
the 24 bit device address. The driver should initialize each drive table
appropriately for the type of disk the driver expects to be used on the
corresponding drive.

\subsubsection{NAME: READ}

{\def\LTcaptype{none} % do not increment counter
\begin{longtable}[]{@{}
  >{\raggedright\arraybackslash}p{(\linewidth - 2\tabcolsep) * \real{0.3182}}
  >{\raggedright\arraybackslash}p{(\linewidth - 2\tabcolsep) * \real{0.6818}}@{}}
\toprule\noalign{}
\endhead
\bottomrule\noalign{}
\endlastfoot
NAME: & READ \\
INPUT: & \begin{minipage}[t]{\linewidth}\raggedright
(U)~=~ADDRESS~OF~THE~DEVICE~STATIC~STORAGE\\
(Y)~=~ADDRESS~OF~THE~PATH~DESCRIPTOR\\
(B)~=~NSB~OF~DISK~LOGICAL~SECTOR~NUMBER\\
(X)~=~LSB\textquotesingle s~OF~DISK~LOGICAL~SECTOR~NUMBER
\end{minipage} \\
OUTPUT: & SECTOR IS RETURNED IN THE SECTOR BUFFER \\
ERROR OUTPUT: & \begin{minipage}[t]{\linewidth}\raggedright
(CC)~=~C~bit~set.\\
(B)~=~Appropriate~error~code.
\end{minipage} \\
FUNCTION: & READ A 256 BYTE SECTOR \\
\end{longtable}
}

Read a sector from the disk and place it in the sector buffer (256
byte). Below are the things that the disk driver must do:

1. Get the sector buffer address from PD.BUF in the path descriptor.

2. Get the drive number from PD.,DRV in the path descriptor.

3. Compute the physical disk address from the logical sector.

number.

4. Initiate the read operation.

5. Copy V.BUSY to V.WAKE, then go to sleep and wait for the I/O to
complete (the IRQ service routine is responsible for sending a wake up
signal). After awakening, test V.WAKE to see if it is clear, if not, go
back to sleep.

If the disk controller can not be interrupt driven it will be necessary
to perform programmed I/O.

NOTE 1: Whenever logical sector zero is read, the first part of this
sector must be copied into the proper drive table (get the drive number
from PD.DRV in the path descriptor). The number of bytes to copy is
DD.SIZ.

NOTE 2: The drive number (PD.DRv) should be used to compute the offset
to the corresponding drive table as follows:

\begin{verbatim}
LDA PD.DRV.Y Get drive number
LDB #DRVMEM Get size of a drive table
MUL
LEAX DRVBEG,U Get address of first table
LEAX D,X           Compute address of table N
\end{verbatim}

\subsubsection{NAME: WRITE}

{\def\LTcaptype{none} % do not increment counter
\begin{longtable}[]{@{}
  >{\raggedright\arraybackslash}p{(\linewidth - 2\tabcolsep) * \real{0.3182}}
  >{\raggedright\arraybackslash}p{(\linewidth - 2\tabcolsep) * \real{0.6818}}@{}}
\toprule\noalign{}
\endhead
\bottomrule\noalign{}
\endlastfoot
NAME: & WRITE \\
INPUT: & \begin{minipage}[t]{\linewidth}\raggedright
(U)~=~ADDRESS~OF~THE~DEVICE~STATIC~STORAGE\\
(Y)~=~ADDRESS~OF~THE~PATH~DESCRIPTOR\\
(B)~=~MSB~OF~DISK~LOGICAL~SECTOR~NUMBER\\
(X)~=~LSB\textquotesingle s~OF~DISK~LOGICAL~SECTOR~NUMBER
\end{minipage} \\
OUTPUT: & THE SECTOR BUFFER IS WRITTEN OUT TO DISK \\
ERROR OUTPUT: & \begin{minipage}[t]{\linewidth}\raggedright
(CC)~=~C~bit~set.\\
(B)~=~Appropriate~error~code.
\end{minipage} \\
FUNCTION: & WRITE A SECTOR \\
\end{longtable}
}

Write the sector buffer (256 bytes) to the disk. Below are the things
that a disk driver must do:

1. Get the sector buffer address from PD.BUF in the path descriptor.

2. Get the drive number from PD.DRV in the path descriptor.

3. Compute the physical disk address from the logical sector number.

4. Initiate the write operation.

5. Copy V.BUSY to V.WAKE, then go to sleep and wait for the I/O to
complete (the IRQ service routine is responsible for sending the wakeup
signal). After awakening, test V.WAKE to see if it is clear, if it is
not, then go back to sleep. If the disk controller can not be
interrupt-driven, it will be necessary to perform a programmed I/O
transfer.

6. If PD.VFY in the path descriptor is equal to zero, read the sector
back in and verify that it was written correctly. This usually does not
involve a compare of the data.

NOTE 1: If disk writes are to be verified, the INIT routine must request
the buffer where the sector may be placed when it is read back in. Do
not copy sector zero into the drive table when it is read back to be
verified.

NOTE 2: Use the drive number (PD.DRV) to compute the offset to the
corresponding drive table as shown for the READ routine.

\subsubsection{NAME: GETSTA PUTSTA}

{\def\LTcaptype{none} % do not increment counter
\begin{longtable}[]{@{}
  >{\raggedright\arraybackslash}p{(\linewidth - 2\tabcolsep) * \real{0.3182}}
  >{\raggedright\arraybackslash}p{(\linewidth - 2\tabcolsep) * \real{0.6818}}@{}}
\toprule\noalign{}
\endhead
\bottomrule\noalign{}
\endlastfoot
NAME: & GETSTA/PUTSTA \\
INPUT: & \begin{minipage}[t]{\linewidth}\raggedright
(U)~=~ADDRESS~OF~THE~DEVICE~STATIC~STORAGE~AREA\\
(Y)~=~ADDRESS~OF~THE~PATH~DESCRIPTOR\\
(A)~=~STATUS~CODE
\end{minipage} \\
OUTPUT: & (DEPENDS UPON THE FUNCTION CODE) \\
ERROR OUTPUT: & \begin{minipage}[t]{\linewidth}\raggedright
(CC)~=~C~bit~set.\\
(B)~=~Appropriate~error~code.
\end{minipage} \\
FUNCTION: & GET/SET DEVICE STATUS \\
\end{longtable}
}

These routines are wild card calls used to get (set) the
device\textquotesingle s operating parameters as specified for the OS9
\hyperref[i.getstt]{I\$GetStt} and \hyperref[i.setstt]{I\$SetStt}
service requests.

It may be necessary to examine or change the register stack which
contains the values of MPU registers at the time of the
\hyperref[i.getstt]{I\$GetStt} or \hyperref[i.setstt]{I\$SetStt} service
request. The address of the register stack may be found in PD.RGS, which
is located in the path descriptor, . The following offsets may be used
to access any particular value in the register stack:

{\def\LTcaptype{none} % do not increment counter
\begin{longtable}[]{@{}
  >{\raggedright\arraybackslash}p{(\linewidth - 8\tabcolsep) * \real{0.1277}}
  >{\raggedright\arraybackslash}p{(\linewidth - 8\tabcolsep) * \real{0.1702}}
  >{\raggedright\arraybackslash}p{(\linewidth - 8\tabcolsep) * \real{0.1064}}
  >{\raggedright\arraybackslash}p{(\linewidth - 8\tabcolsep) * \real{0.0638}}
  >{\raggedright\arraybackslash}p{(\linewidth - 8\tabcolsep) * \real{0.5319}}@{}}
\toprule\noalign{}
\begin{minipage}[b]{\linewidth}\raggedright
OFFSET
\end{minipage} & \begin{minipage}[b]{\linewidth}\raggedright
MNEMONIC
\end{minipage} & \begin{minipage}[b]{\linewidth}\raggedright
MPU REGISTER
\end{minipage} & \begin{minipage}[b]{\linewidth}\raggedright
\end{minipage} & \begin{minipage}[b]{\linewidth}\raggedright
\end{minipage} \\
\midrule\noalign{}
\endhead
\bottomrule\noalign{}
\endlastfoot
\$0 & R\$CC & RMB & 1 & CONDITIONS CODE REGISTER \\
\$1 & R\$D & EQU & . & D REGISTER \\
\$1 & R\$A & RMB & 1 & A REGISTER \\
\$2 & R\$B & RMB & 1 & B REGISTER \\
\$3 & R\$DP & RMB & 1 & DP REGISTER \\
\$4 & R\$X & RMB & 2 & X REGISTER \\
\$6 & R\$Y & RMB & 2 & Y REGISTER \\
\$8 & R\$U & RMB & 2 & U REGISTER \\
\$A & R\$PC & RMB & 2 & PROGRAM COUNTER \\
\end{longtable}
}

\subsubsection{NAME: TERM}

{\def\LTcaptype{none} % do not increment counter
\begin{longtable}[]{@{}
  >{\raggedright\arraybackslash}p{(\linewidth - 2\tabcolsep) * \real{0.3182}}
  >{\raggedright\arraybackslash}p{(\linewidth - 2\tabcolsep) * \real{0.6818}}@{}}
\toprule\noalign{}
\endhead
\bottomrule\noalign{}
\endlastfoot
NAME: & TERM \\
INPUT: & (U) = ADDRESS OF DEVICE STATIC STORAGE AREA \\
OUTPUT: & NONE \\
ERROR OUTPUT: & \begin{minipage}[t]{\linewidth}\raggedright
(CC)~=~C~bit~set.\\
(B)~=~Appropriate~error~code.
\end{minipage} \\
FUNCTION: & TERMINATE DEVICE \\
\end{longtable}
}

This routine is called when a device is no longer in use in the system,
which is defined to be when the link count of its device descriptor
module becomes zero). The TERM routine must:

1. Wait until any pending I/O has completed.

2. Disable the device interrupts.

3. Remove the device from the IRQ polling list.

4. If the INIT routine reserved a 256 byte buffer for verifying disk
writes, return the memory with the \hyperref[f.mem]{F\$Mem} service
request.

\subsubsection{NAME: IRQ SERVICE ROUTINE}

{\def\LTcaptype{none} % do not increment counter
\begin{longtable}[]{@{}
  >{\raggedright\arraybackslash}p{(\linewidth - 2\tabcolsep) * \real{0.3182}}
  >{\raggedright\arraybackslash}p{(\linewidth - 2\tabcolsep) * \real{0.6818}}@{}}
\toprule\noalign{}
\endhead
\bottomrule\noalign{}
\endlastfoot
NAME: & IRQ SERVICE ROUTINE \\
FUNCTION: & SERVICE DEVICE INTERRUPTS \\
\end{longtable}
}

Although this routine is not included in the device driver module branch
table and is not called directly by RBFMAN, it is an key routine in
interrupt-driven device drivers. Its function is to:

1. Service device interrupts.

2. When the I/O is complete, the IRQ service routine should send a wake
up signal to the process whose process ID is in V.WAKE

Also clear V.WAKE as a flag to the mainline program that the IRQ has
indeed occurred.

NOTE: When the IRQ service routine finishes servicing an interrupt it
must clear the array and exit with an RTS instruction.

\subsubsection{NAME: BOOT (Bootstrap Module)}

{\def\LTcaptype{none} % do not increment counter
\begin{longtable}[]{@{}
  >{\raggedright\arraybackslash}p{(\linewidth - 2\tabcolsep) * \real{0.3182}}
  >{\raggedright\arraybackslash}p{(\linewidth - 2\tabcolsep) * \real{0.6818}}@{}}
\toprule\noalign{}
\endhead
\bottomrule\noalign{}
\endlastfoot
NAME: & TERM \\
INPUT: & None. \\
OUTPUT: & \begin{minipage}[t]{\linewidth}\raggedright
(U)~=~SIZE~OF~THE~BOOT~FILE~(in~bytes)\\
(X)~=~ADDRESS~OF~WHERE~THE~BOOT~FILE~WAS~LOADED~IN~MEMORY
\end{minipage} \\
ERROR OUTPUT: & \begin{minipage}[t]{\linewidth}\raggedright
(CC)~=~C~bit~set.\\
(B)~=~Appropriate~error~code.
\end{minipage} \\
FUNCTION: & LOAD THE BOOT FILE INTO MEMORY FROM MASS-STORAGE \\
\end{longtable}
}

NOTE: The BOOT module is \emph{not} part of the disk driver. It is a
separate module which is normally co-resident with the ``OS9P2'' module
in the system firmware.

The bootstrap module contains one subroutine that loads the bootstrap
file and some related information into memory, it uses the standard
executable module format with a module type of ``system'' (code \$C).
The execution offset in the module header contains the offset to the
entry point of this subroutine.

It obtains the starting sector number and size of the \texttt{OS9Boot}
file from the identification sector (LSN 0). OS-9 is called to allocate
a memory area large enough for the boot file, and then it loads the boot
file into this memory area.

1. Read the identification sector (sector zero) from the disk. BOOT must
pick its own buffer area. The identification sector contains the values
for DD.BT (the 24 bit logical sector number of the bootstrap file), and
DD.BSZ (the size of the bootstrap file in bytes). For a full description
of the identification sector. See \hyperref[lsn1]{Disk Allocation Map
Sector}.

2. After reading the identification sector into the buffer, get the 24
bit logical sector number of the bootstrap file from DD.BT.

3. Get the size (in bytes) of the bootstrap file from DD.BSZ. The boot
is contained in one logically contiguous block beginning at the logical
sector specified in DD.BT and extending for (DD.BSZ/256+1) sectors.

4. Use the OS9 \hyperref[f.srqmem]{F\$SRqMem} service request to request
the memory area where the boot file will be loaded into.

5. Read the boot file into this memory area.

6. Return the size of the boot file and its location.

\section{Sequential Character File Manager}

The Sequential Character File Manager (SCFMAN) is the OS-9 file manager
module that supports devices that operate on a character-by-character
basis, such as terminals, printers, modems, etc. SCFMAN can handle any
number or type of such devices. It is a reentrant subroutine package
called by IOMAN for I/O service requests to sequential
character-oriented devices. It includes the extensive input and output
editing functions typical of line-oriented operation such as: backspace,
line delete, repeat line, auto line feed. Screen pause, return delay
padding, etc.

Standard OS-9 systems are supplied with SCFMAN and two SCF-type device
driver modules: ACIA, which run 6850 serial interfaces, and PIA, which
drives a 6821-type parallel interface for printers.

\subsection{SCFMAN Line Editing Functions}

\hyperref[i.read]{I\$Read} and \hyperref[i.write]{I\$Write} service
requests (which correspond to Basic09 GET and PUT statements) to
SCFMAN-type devices pass data to/from the device without any
modification, except that keyboard interrupt, keyboard abort, and pause
character are filtered out of the input (editing is disabled if the
corresponding character in the path descriptor contains a zero). In
particular, carriage returns are not automatically followed by line
feeds or nulls, and the high order bits are passed as sent/received.

\hyperref[i.readln]{I\$ReadLn} and \hyperref[i.writln]{I\$WritLn}
service requests (which correspond to Basic09 INPUT, PRINT, READ and
WRITE statements) to SCFMAN-type devices perform full line editing of
all functions enabled for the particular device. These functions are
initialized when the device is first used by copying the option table
from the device descriptor table associated with the specific device.
They may be altered anytime afterwards from assembly language programs
using the \hyperref[i.setstt]{I\$SetStt} and
\hyperref[i.getstt]{I\$GetStt} service requests, or from the keyboard
using the \texttt{tmode} command. Also, all bytes transfered in this
mode will have the high order bit cleared.

The following path descriptor values control the line editing functions:

If PD.UPC \textless\textgreater{} 0 bytes input or output in the range
``a..z'' are made ``A..Z''

If PD.EKO \textless\textgreater{} 0, input bytes are echoed, except that
undefined control characters in the range \$0..\$1F print as ``.''

If PD.ALF \textless\textgreater{} 0, carriage returns are automatically
followed by line feeds.

If PD.NUL \textless\textgreater{} 0, After each CR/LF a PD.NUL ``nulls''
(always \$00) are sent.

If PD.PAU \textless\textgreater{} 0, Auto page pause will occur after
every PD.PAU lines since the last input.

If PD.BSP \textless\textgreater{} 0, SCF will recognize PD.BSP as the
``input'' backspace character, and will echo PD.BSE (the backspace echo
character) if PD.BSO = 0, or PD.BSE, space, PD.BSE if PD.BSO
\textless\textgreater{} 0.

If PD.DEL \textless\textgreater{} 0, SCF will recognize PD.DEL the
delete line character (on input), and echo the backspace sequence over
the entire line if PD.DLO = 0, or echo CR/LF if PD.DLO
\textless\textgreater{} 0.

PD.EOR defines the end of record character. This is the last character
an each line entered (\hyperref[i.readln]{I\$ReadLn}), and terminates
the output (\hyperref[i.writln]{I\$WritLn}) when this character is sent.
Normally PD.EOR will be set to \$0D. If it is set to zero,
SCF\textquotesingle s \hyperref[i.readln]{I\$ReadLn} will NEVER
terminate, unless an EOF occurs.

If PD.EOF \textless\textgreater{} 0, it defines the end of file
character. SCFMAN will return an end-of-file error on
\hyperref[i.read]{I\$Read} or \hyperref[i.readln]{I\$ReadLn} if this is
the first (and only) character input. It can be disabled by setting its
value to zero.

If PD.RPR \textless\textgreater{} 0, SCF
(\hyperref[i.readln]{I\$ReadLn}) will, upon receipt of this character,
echo a carriage return {[}optional line feed{]}, and then reprint the
current line.

If PD.DUP \textless\textgreater{} 0, SCF
(\hyperref[i.readln]{I\$ReadLn}) will duplicate whatever is in the input
buffer through the first ``PD.EOR'' character.

If PD.PSC \textless\textgreater{} 0, output is suspended before the next
``PD.EOR'' character when this character is input. This will also delete
any ``type ahead'' input for \hyperref[i.readln]{I\$ReadLn}.

If PD.INT \textless\textgreater{} 0, and is received on input, a
keyboard interrupt signal is sent to the last user of this path. Also it
will terminate the current I/O request (it any) with an error identical
to the keyboard interrupt signal code. This location normally is set to
a { +control+C} character.

If PD.QUT \textless\textgreater{} 0, and is received on input, a
keyboard abort signal is sent to the last user of this path. Also it
will terminate the current I/O request (if any) with an error code
identical to the keyboard interrrupt signal code. This location is
normally set to a { +control+Q} character.

If PD.OVF \textless\textgreater{} 0, It is echoed when
\hyperref[i.readln]{I\$ReadLn} has satisfied its input byte count
without finding a ``PD.EOR'' character.

NOTE: It is possible to disable most of these special editing functions
by setting the corresponding control character in the path descriptor to
zero by using the \hyperref[i.setstt]{I\$SetStt} service request, or by
running the \texttt{tmode} utility. A more permanent solution may be had
by setting the corresponding control character value in the device
descriptor module to zero.

Device descriptors may be inspected to determine the default settings
for these values for specific devices.

\subsection{SCFMAN Definitions of The Path Descriptor}

The table below describes the path descriptors used by SCFMAN and
SCFMAN-type device drivers.

{\def\LTcaptype{none} % do not increment counter
\begin{longtable}[]{@{}
  >{\raggedright\arraybackslash}p{(\linewidth - 6\tabcolsep) * \real{0.1613}}
  >{\raggedright\arraybackslash}p{(\linewidth - 6\tabcolsep) * \real{0.0968}}
  >{\raggedleft\arraybackslash}p{(\linewidth - 6\tabcolsep) * \real{0.0968}}
  >{\raggedright\arraybackslash}p{(\linewidth - 6\tabcolsep) * \real{0.6452}}@{}}
\toprule\noalign{}
\begin{minipage}[b]{\linewidth}\raggedright
Name
\end{minipage} & \begin{minipage}[b]{\linewidth}\raggedright
Offset
\end{minipage} & \begin{minipage}[b]{\linewidth}\raggedleft
Size
\end{minipage} & \begin{minipage}[b]{\linewidth}\raggedright
Description
\end{minipage} \\
\midrule\noalign{}
\endhead
\bottomrule\noalign{}
\endlastfoot
Universal Section (same for all file managers) & & & \\
PD.PD & \$00 & 1 & Path number \\
PD.MOD & \$01 & 1 & Mode (read/write/update) \\
PD.CNT & \$02 & 1 & Number of open images \\
PD.DEV & \$03 & 2 & Address of device table entry \\
PD.CPR & \$05 & 1 & Current process ID \\
PD.RGS & \$06 & 2 & Address of callers register stack \\
PD.BUF & \$08 & 2 & Buffer address \\
SCFMAN Path Descriptor Definitions & & & \\
PD.DV2 & \$0A & 2 & Device table addr of 2nd (echo) device \\
PD.RAW & \$0C & 1 & Edit flag: 0=raw mode, 1=edit mode \\
PD.MAX & \$0D & 2 & Headline maximum character count \\
PD.MIN & \$0F & 1 & Devices are ``mine'' if cleared \\
PD.STS & \$10 & 2 & Status routine module address \\
PD.STM & \$12 & 2 & Reserved for status routine \\
SCFMAN Option Section Definition & & & \\
& \$20 & 1 & Device class 0=SCF 1=RBF 2=PIPE 3=SBF \\
PD.UPC & \$21 & 1 & Case (0=BOTH, 1=UPPER ONLY) \\
PD.BSO & \$22 & 1 & Backsp (0=BSE, 1=BSE SP BSE) \\
PD.DLO & \$23 & 1 & Delete (0 = BSE over line, 1=CR LF) \\
PD.EKO & \$24 & 1 & Echo (0=no echo) \\
PD.ALF & \$25 & 1 & Auto LF (0=no auto LF) \\
PD.NUL & \$26 & 1 & End of line null count \\
PD.PAU & \$27 & 1 & Pause (0= no end of page pause) \\
PD.PAG & \$28 & 1 & Lines per page \\
PD.BSP & \$29 & 1 & Backspace character \\
PD.DEL & \$2A & 1 & Delete line character \\
PD.EOR & \$25 & 1 & End of record character (read only) \\
PD.EOF & \$2C & 1 & End of file character (read only) \\
PD.RPR & \$2D & 1 & Reprint line character \\
PD.DUP & \$25 & 1 & Duplicate last line character \\
PD.PSC & \$2F & 1 & Pause character \\
PD.INT & \$30 & 1 & Keyboard interrupt character (CTL C) \\
PD.QUT & \$31 & 1 & Keyboard abort character (CTL Q) \\
PD.BSE & \$32 & 1 & Backspace echo character (BSE) \\
PD.OVF & \$33 & 1 & Line overflow character (bell) \\
PD.PAR & \$34 & 1 & Device initialization value (parity) \\
PD.BAU & \$35 & 1 & Software settable baud rate \\
PD.D2P & \$36 & 2 & Offset to 2nd device name string \\
PD.STN & \$38 & 2 & Offset of status routine name \\
PD.ERR & \$3A & 1 & Most recent I/O error status \\
\end{longtable}
}

The first section is universal for all file managers, the second and
third section are specific for SCFMAN and SCFMAN-type device drivers.
The option section of the path descriptor contains many device operating
parameters whicb may be read or written by the OS9
\hyperref[i.getstt]{I\$GetStt} or \hyperref[i.setstt]{I\$SetStt} service
requests. IOMAN initializes this section when a path is opened to a
device by copying the corresponding device descriptor initialization
table. Any values not determined by this table will default to zero.

Special editing functions may be disabled by setting the corresponding
control character value to zero.

\subsection{SCF Device Descriptor Modules}

Device descriptor modules for SCF-type devices contain the device
address and an initialization table which defines inital values for the
I/O editing features, as listed below.

{\def\LTcaptype{none} % do not increment counter
\begin{longtable}[]{@{}
  >{\raggedright\arraybackslash}p{(\linewidth - 6\tabcolsep) * \real{0.1169}}
  >{\raggedright\arraybackslash}p{(\linewidth - 6\tabcolsep) * \real{0.1429}}
  >{\raggedright\arraybackslash}p{(\linewidth - 6\tabcolsep) * \real{0.0909}}
  >{\raggedright\arraybackslash}p{(\linewidth - 6\tabcolsep) * \real{0.6494}}@{}}
\toprule\noalign{}
\begin{minipage}[b]{\linewidth}\raggedright
MODULE OFFSET
\end{minipage} & \begin{minipage}[b]{\linewidth}\raggedright
\end{minipage} & \begin{minipage}[b]{\linewidth}\raggedright
ORG \$12
\end{minipage} & \begin{minipage}[b]{\linewidth}\raggedright
\end{minipage} \\
\midrule\noalign{}
\endhead
\bottomrule\noalign{}
\endlastfoot
& TABLE & EQU . & BEGINING OF OPTION TABLE \\
\$12 & IT.DVC & RMB 1 & DEVICE CLASS (0=SCF 1=RBF 2=PIPE 3=SBF) \\
\$13 & IT.UPC & RMB 1 & CASE (0=BOTH, 1=UPPER ONLY) \\
\$14 & IT.BSO & RMB 1 & BACK SPACE (0=BSE, 1=BSE,SP,BSE) \\
\$15 & IT.DLO & RMB 1 & DELETE (0=BSE OVER LINE, 1=CR) \\
\$16 & IT.EKO & RMB 1 & ECHO (0=NO ECHO) \\
\$17 & IT.ALF & RMB 1 & AUTO LINE FEED (0= NO AUTO LF) \\
\$18 & IT.NUL & RMB 1 & END OF LINE NULL COUNT \\
\$19 & IT.PAU & RMB 1 & PAUSE (0= NO END OF PAGE PAUSE) \\
\$1A & IT.PAG & RMB 1 & LINES PER PAGE \\
\$1B & IT.BSP & RMB 1 & BACKSPACE CHARACTER \\
\$1C & IT.DEL & RMB 1 & DELETE LINE CHARACTER \\
\$1D & IT.EOR & RMB 1 & END OF RECORD CHARACTER \\
\$1E & IT.EOF & RMB 1 & END OF FILE CHARACTER \\
\$1F & IT.RPR & RMB 1 & REPRINT LINE CHARACTER \\
\$20 & IT.DUP & RMB 1 & DUP LAST LINE CHARACTER \\
\$21 & IT.PSC & RMB 1 & PAUSE CHARACTER \\
\$22 & IT.INT & RMB 1 & INTERRUPT CHARACTER \\
\$23 & IT.QUT & RMB 1 & QUIT CHARACTER \\
\$24 & IT.BSE & RMB 1 & BACKSPACE ECHO CHARACTER \\
\$25 & IT.OVF & RMB 1 & LINE OVERFLOW CHARACTER (BELL) \\
\$26 & IT.PAR & RMB 1 & INITIALIZATION VALUE (PARITY) \\
\$27 & IT.BAU & RMB 1 & BAUD RATE \\
\$28 & IT.D2P & RMB 2 & ATTACHED DEVICE NAME STRING OFFSET \\
\$2A & IT.STN & RMB 2 & OFFSET TO STATUS ROUTINE \\
\$2C & IT.ERR & RMB 1 & INITIAL ERROR STATUS \\
\end{longtable}
}

NOTES:

SCF editing functions will be ``turned off'' if the corresponding
special character is a zero. For example, it IT.EOF was a zero, there
would be no end of file character.

IT.PAR is typically used to intitialize the device\textquotesingle s
control register when a path is opened to it.

\subsection{SCF Device Driver Storage Definitions}

An SCFMAN-type device driver module contains a package of subroutines
that perform raw I/O transfers to or from a specific hardware
controller. These modules are usually reentrant so that one copy of the
module can simultaneously run several different devices that use
identical I/O controllers. For each ``incarnation'' of the driver, IOMAN
will allocate a static storage area for that device. The size of the
storage area is given in the device driver module header. Some of this
storage area will be used by IOMAN and SCFMAN, the device driver is free
to use the remainder in any way (typically as variables and butters).
This static storage is defined as:

{\def\LTcaptype{none} % do not increment counter
\begin{longtable}[]{@{}
  >{\raggedright\arraybackslash}p{(\linewidth - 6\tabcolsep) * \real{0.1200}}
  >{\raggedright\arraybackslash}p{(\linewidth - 6\tabcolsep) * \real{0.1800}}
  >{\raggedright\arraybackslash}p{(\linewidth - 6\tabcolsep) * \real{0.1400}}
  >{\raggedright\arraybackslash}p{(\linewidth - 6\tabcolsep) * \real{0.5600}}@{}}
\toprule\noalign{}
\begin{minipage}[b]{\linewidth}\raggedright
OFFSET
\end{minipage} & \begin{minipage}[b]{\linewidth}\raggedright
\end{minipage} & \begin{minipage}[b]{\linewidth}\raggedright
ORG 0
\end{minipage} & \begin{minipage}[b]{\linewidth}\raggedright
\end{minipage} \\
\midrule\noalign{}
\endhead
\bottomrule\noalign{}
\endlastfoot
\$0 & V.PAGE & RMB 1 & PORT EXTENDED ADDRESS \\
\$1 & V.PORT & RMB 2 & DEVICE BASE ADDRESS \\
\$3 & V.LPRC & RMB 1 & LAST ACTIVE PROCESS ID \\
\$4 & V.BUSY & RMB 1 & ACTIVE PROCESS ID (0 NOT BUSY) \\
\$5 & V. WAKE & RMB 1 & PROCESS ID TO REAWAKEN \\
& V. USER & EQU . & END OF OS9 DEFINITIONS \\
\$6 & V.TYPE & RMB 1 & DEVICE TYPE OR PARITY \\
\$7 & V.LINE & RMB 1 & LINES LEFT TILL END OF PAGE \\
\$8 & V.PAUS & RMB 1 & PAUSE REQUEST (0 = NO PAUSE) \\
\$9 & V.DEV2 & RMB 2 & ATTACHED DEVICE STATIC STORAGE \\
\$B & V. INTR & RMB 1 & INTERRUPT CHARACTER \\
\$C & V.QUIT & RMB 1 & QUIT CHARACTER \\
\$D & V.PCHR & RMB 1 & PAUSE CHARACTER \\
\$E & V. ERR & RMB 1 & ERROR ACCUMULATOR \\
\$F & V.SCF & EQU . & END OF SCFMAN DEFINITIONS \\
& FREE & EQU . & FREE FOR DEVICE DRIVER TO USE \\
\end{longtable}
}

V.PAGE, V.PORT These three bytes are defined by IOMAN to be the 24 bit
device address.

V.LPRC This location contains the process ID of the last process to use
the device. The IRQ service routine is responsible for sending this
process the proper signal in case a ``QUIT'' character or an
``INTERRUPT'' character is recieved. Defined by SCFMAN.

V. BUSY This location contains the process ID of the process currently
using the device (zero if it is not being used). This is used by SCFMAN
to prevent more than one process from using the device at the same
moment. Defined by SCFMAN.

V.WAKE This location contains the process ID of any process that is
waiting for the device to complete I/O (or zero if there is none
waiting). The interrupt service routine should check this location to
see if a process is waiting and if so, send it a wake up signal. Defined
by the device driver.

V.TYPE This location contains any special characteristics of a device.
It is typically used as a value to initialize the device control
register, for parity etc. It is defined by SCFMAN which copies its value
from PD.PAR in the path descriptor.

V.LINE This location contains the number of lines left till end of page.
Paging is handled by SCFMAN and not by the device driver.

V.PAUS This location is a flag used by SCFMAN to indicate that a pause
character has been recieved. Setting its value to anything other than
zero will cause SCFMAN to stop transmitting characters at the end of the
next line. Device driver input routines must set V.PAUS in the ECHO
device\textquotesingle s static storage area. SCFMAN will check this
value in the ECHO device\textquotesingle s static storage when output is
sent.

V.DEV2 This location contains the address of the ECHO (attached)
device\textquotesingle s static storage area. Typically the device and
the attached device are one and the same. However they may be different
as in the case of a keyboard and a memory mapped video display. Defined
by SCFMAN.

V.INTR Keyboard interrupt character. It is defined by SCFMAN, which
copies its value from PD.INT in the path descriptor.

V.QUIT Keyboard abort character. It is defined by SCFMAN which copies
its value from PD.QUT in the path descriptor.

V.PCHR Pause character. It is defined by SCFMAN which copies its value
from PD.PSC in the path descriptor.

V.ERR This location is used to accumulate I/O errors. Typically it is
used by the IRQ service routine to record errors so that they may be
reported later when SCFMAN calls one of the device driver routines.

\subsection{SCFMAN Device Driver Subroutines}

As with all device drivers. SCFMAN device drivers use a standard
executable memory module format with a module type of ``device driver''
(CODE \$5). The execution offset address in the module header points to
a branch table that has six three byte entries. Each entry is typically
a LBRA to the corresponding subroutine. The branch table is as follows:

{\def\LTcaptype{none} % do not increment counter
\begin{longtable}[]{@{}
  >{\raggedright\arraybackslash}p{(\linewidth - 6\tabcolsep) * \real{0.2000}}
  >{\raggedright\arraybackslash}p{(\linewidth - 6\tabcolsep) * \real{0.2000}}
  >{\raggedright\arraybackslash}p{(\linewidth - 6\tabcolsep) * \real{0.2000}}
  >{\raggedright\arraybackslash}p{(\linewidth - 6\tabcolsep) * \real{0.4000}}@{}}
\toprule\noalign{}
\endhead
\bottomrule\noalign{}
\endlastfoot
\multirow{6}{=}{ENTRY} & LBRA & INIT & INITIALIZE DEVICE \\
& LBRA & READ & READ CHARACTER \\
& LBRA & WRITE & WRITE CHARACTER \\
& LBRA & GETSTA & GET DEVICE STATUS \\
& LBRA & SETSTA & SET DEVICE STATUS \\
& LBRA & TERM & TERMINATE DEVICE \\
\end{longtable}
}

Each subroutine should exit with the condition code register C bit
cleared it no error occured. Otherwise the C bit should be set and an
appropriate error code returned in the B register. Below is a
description of each subroutine, its input parameters and its output
parameters.

\subsubsection{NAME: INIT}

{\def\LTcaptype{none} % do not increment counter
\begin{longtable}[]{@{}
  >{\raggedright\arraybackslash}p{(\linewidth - 2\tabcolsep) * \real{0.3182}}
  >{\raggedright\arraybackslash}p{(\linewidth - 2\tabcolsep) * \real{0.6818}}@{}}
\toprule\noalign{}
\endhead
\bottomrule\noalign{}
\endlastfoot
NAME: & INIT \\
INPUT: & \begin{minipage}[t]{\linewidth}\raggedright
(U)~=~ADDRESS~OF~DEVICE~STATIC~STORAGE\\
(Y)~=~ADDRESS~OF~DEVICE~DESCRIPTOR~MODULE
\end{minipage} \\
OUTPUT: & NONE \\
ERROR OUTPUT: & \begin{minipage}[t]{\linewidth}\raggedright
(CC)~=~C~bit~set.\\
(B)~=~Appropriate~error~code.
\end{minipage} \\
FUNCTION: & INITIALIZE DEVICE AND ITS STATIC STORAGE \\
\end{longtable}
}

1. Initialize the device static storage.

2. Place the IRQ service routine on the IRQ polling list by using the
OS9 \hyperref[f.irq]{F\$IRQ} service request.

3. Initialize the device control registers (enable interrupts if
necessary).

NOTE: Prior to being called, the device static storage will be cleared
(set to zero) except for V.PAGE and V.PORT which will contain the 24 bit
device address. There is no need to initialize the portion of static
storage used by IOMAN and SCFMAN.

\subsubsection{NAME: READ}

{\def\LTcaptype{none} % do not increment counter
\begin{longtable}[]{@{}
  >{\raggedright\arraybackslash}p{(\linewidth - 2\tabcolsep) * \real{0.3182}}
  >{\raggedright\arraybackslash}p{(\linewidth - 2\tabcolsep) * \real{0.6818}}@{}}
\toprule\noalign{}
\endhead
\bottomrule\noalign{}
\endlastfoot
NAME: & READ \\
INPUT: & \begin{minipage}[t]{\linewidth}\raggedright
(U)~=~ADDRESS~OF~DEVICE~STATIC~STORAGE\\
(Y)~=~ADDRESS~OF~PATH~DESCRIPTOR
\end{minipage} \\
OUTPUT: & (A) = CHARACTER READ \\
ERROR OUTPUT: & \begin{minipage}[t]{\linewidth}\raggedright
(CC)~=~C~bit~set.\\
(B)~=~Appropriate~error~code.
\end{minipage} \\
FUNCTION: & GET NEXT CHARACTER \\
\end{longtable}
}

This routine should get the next character from the input buffer. If
there is no data ready, this routine should copy its process ID from
V.BUSY into V.WAKE and then use the \hyperref[f.sleep]{F\$Sleep} service
request to put itself to sleep.

Later when data is recieved, the IRQ service routine will leave the data
in a buffer, then check V.WAKE to see if any process is waiting for the
device to complete I/O. If so, the IRQ service routine should send a
wakeup signal to it.

NOTE: Data buffers are \emph{not} automatically allocated. It any are
used, they should be defined in the device\textquotesingle s static
storage area.

\subsubsection{NAME: WRITE}

{\def\LTcaptype{none} % do not increment counter
\begin{longtable}[]{@{}
  >{\raggedright\arraybackslash}p{(\linewidth - 2\tabcolsep) * \real{0.3182}}
  >{\raggedright\arraybackslash}p{(\linewidth - 2\tabcolsep) * \real{0.6818}}@{}}
\toprule\noalign{}
\endhead
\bottomrule\noalign{}
\endlastfoot
NAME: & WRITE \\
INPUT: & \begin{minipage}[t]{\linewidth}\raggedright
(U)~=~ADDRESS~OF~DEVICE~STATIC~STORAGE\\
(Y)~=~ADDRESS~OF~THE~PATH~DESCRIPTOR\\
(A)~=~CHAR~TO~WRITE
\end{minipage} \\
OUTPUT: & NONE \\
ERROR OUTPUT: & \begin{minipage}[t]{\linewidth}\raggedright
(CC)~=~C~bit~set.\\
(B)~=~Appropriate~error~code.
\end{minipage} \\
FUNCTION: & OUTPUT A CHARACTER \\
\end{longtable}
}

This routine places a data byte into an output buffer and enables the
device output interrupts. It the data buffer is already full, this
routine should copy its process ID from V.BUSY into V.WAKE and then put
itself to sleep.

Later when the IRQ service routine transmits a character and makes room
for more data in the buffer, it will check V.WAKE to see if there is a
process waiting for the device to complete I/O. It there is, it will
send a wake up signal to that process.

NOTE: This routine must ensure that the IRQ service routine will start
up when data is placed into the buffer. After an interrupt is generated
the IRQ service routine will continue to transmit data until the data
butter is empty, and then it will disable the device\textquotesingle s
``ready to transmit'' interrupts.

NOTE: Data buffers are \emph{not} automatically allocated. If any are
used, they should be defined in the device\textquotesingle s static
storage.

\subsubsection{NAME: GETSTA/SETSTA}

{\def\LTcaptype{none} % do not increment counter
\begin{longtable}[]{@{}
  >{\raggedright\arraybackslash}p{(\linewidth - 2\tabcolsep) * \real{0.3182}}
  >{\raggedright\arraybackslash}p{(\linewidth - 2\tabcolsep) * \real{0.6818}}@{}}
\toprule\noalign{}
\endhead
\bottomrule\noalign{}
\endlastfoot
NAME: & GETSTA/SETSTA \\
INPUT: & \begin{minipage}[t]{\linewidth}\raggedright
(U)~=~ADDRESS~OF~DEVICE~STATIC~STORAGE\\
(Y)~=~ADDRESS~OF~PATH~DESCRIPTOR\\
(A)~=~STATUS~CODE
\end{minipage} \\
OUTPUT: & DEPENDS UPON FUNCTION CODE \\
FUNCTION: & GET/SET DEVICE STATUS \\
\end{longtable}
}

This routine is a wild card call used to get (set) the device parameters
specified in the \hyperref[i.getstt]{I\$GetStt} and
\hyperref[i.setstt]{I\$SetStt} service requests. Currently all of the
function codes defined by Microware for SCF-type devices are handled by
IOMAN or SCFMAN. Any codes not defined by Microware will be passed to
the device driver.

It may be necessary to examine or change the register packet which
contains the values of the 6809 registers at the time the OS9 service
request was issued. The address of the register packet may be found in
PD.RGS, which is located in the path descriptor. The following offsets
may be used to access any particular value in the register packet:

{\def\LTcaptype{none} % do not increment counter
\begin{longtable}[]{@{}
  >{\raggedright\arraybackslash}p{(\linewidth - 8\tabcolsep) * \real{0.1333}}
  >{\raggedright\arraybackslash}p{(\linewidth - 8\tabcolsep) * \real{0.1333}}
  >{\raggedright\arraybackslash}p{(\linewidth - 8\tabcolsep) * \real{0.1111}}
  >{\raggedright\arraybackslash}p{(\linewidth - 8\tabcolsep) * \real{0.0667}}
  >{\raggedright\arraybackslash}p{(\linewidth - 8\tabcolsep) * \real{0.5556}}@{}}
\toprule\noalign{}
\begin{minipage}[b]{\linewidth}\raggedright
OFFSET
\end{minipage} & \begin{minipage}[b]{\linewidth}\raggedright
MNEMONIC
\end{minipage} & \begin{minipage}[b]{\linewidth}\raggedright
MPU REGISTER
\end{minipage} & \begin{minipage}[b]{\linewidth}\raggedright
\end{minipage} & \begin{minipage}[b]{\linewidth}\raggedright
\end{minipage} \\
\midrule\noalign{}
\endhead
\bottomrule\noalign{}
\endlastfoot
\$0 & R\$CC & RMB & 1 & CONDITIONS CODE REGISTER \\
\$1 & R\$D & EQU & . & D REGISTER \\
\$1 & R\$A & RMB & 1 & A REGISTER \\
\$2 & R\$B & RMB & 1 & B REGISTER \\
\$3 & R\$DP & RMB & 1 & DP REGISTER \\
\$4 & R\$X & RMB & 2 & X REGISTER \\
\$6 & R\$Y & RMB & 2 & Y REGISTER \\
\$8 & R\$U & RMB & 2 & U REGISTER \\
\$A & R\$PC & RMB & 2 & PROGRAM COUNTER \\
\end{longtable}
}

\subsubsection{NAME. TERM}

{\def\LTcaptype{none} % do not increment counter
\begin{longtable}[]{@{}
  >{\raggedright\arraybackslash}p{(\linewidth - 2\tabcolsep) * \real{0.3182}}
  >{\raggedright\arraybackslash}p{(\linewidth - 2\tabcolsep) * \real{0.6818}}@{}}
\toprule\noalign{}
\endhead
\bottomrule\noalign{}
\endlastfoot
NAME: & TERM \\
INPUT: & (U) = PTR TO DEVICE STATIC STORAGE \\
OUTPUT: & NONE \\
ERROR OUTPUT: & \begin{minipage}[t]{\linewidth}\raggedright
(CC)~=~C~bit~set.\\
(B)~=~Appropriate~error~code.
\end{minipage} \\
FUNCTION: & TERMINATE DEVICE \\
\end{longtable}
}

This routine is called when a device is no longer in use, defined as
when its device descriptor module\textquotesingle s link count becomes
zero). It must perform the following:

1. Wait until the output buffer has been emptied (by the IRQ service
routine)

2. Disable device interrupts.

3. Remove device from the IRQ polling list.

NOTE: Static storage used by device drivers is never returned to the
free memory pool. Therefore, it is desirable to NEVER terminate any
device that might be used again. Modules contained in the \texttt{Boot}
file will NEVER be terminated.

\subsubsection{NAME: IRQ SERVICE ROUTINE}

{\def\LTcaptype{none} % do not increment counter
\begin{longtable}[]{@{}
  >{\raggedright\arraybackslash}p{(\linewidth - 2\tabcolsep) * \real{0.3182}}
  >{\raggedright\arraybackslash}p{(\linewidth - 2\tabcolsep) * \real{0.6818}}@{}}
\toprule\noalign{}
\endhead
\bottomrule\noalign{}
\endlastfoot
NAME: & IRQ SERVICE ROUTINE \\
FUNCTION: & SERVICE DEVICE INTERRUPTS \\
\end{longtable}
}

Although this routine is not included in the device drivers branch table
and not called directly from SCFMAN, it is an important routine in
device drivers. The main things that it does are:

\begin{enumerate}
\def\labelenumi{\arabic{enumi}.}
\item
  Service the device interrupts (recieve data from device or send data
  to it). This routine should put its data into and get its data from
  buffers which are defined in the device static storage.
\item
  Wake up any process waiting for I/O to complete by checking to see if
  there is a process ID in V.WAKE (non-zero) and it so send a wakeup
  signal to that process.
\item
  If the device is ready to send more data and the output buffer is
  emoty, disable the device\textquotesingle s ``ready to transmit''
  interrupts.
\item
  If a pause character is recieved, set V.PAUS in the attached device
  static storage to a non-zero value. The address of the attached device
  static storage is in V.DEV2.
\end{enumerate}

When the IRQ service routine finishes servicing an interrupt, it must
clear the carry and exit with an RTS instruction.

\section{Assembly Language Programming Techniques}

There are four key rules for programmers writing OS-9 assembly language
programs:

\begin{enumerate}
\def\labelenumi{\arabic{enumi}.}
\item
  All programs \emph{must} use position-independent-code (PIC). OS9
  selects load addresses based on available memory at run-time. There is
  no way to force a program to be loaded at a specific address.
\item
  All programs must use the standard OS-9 memory module formats or they
  cannot be loaded and run. Programs must not use self-modifying code.
  Programs must not change anything in a memory module or use any part
  of it for variables.
\item
  Storage for all variables and data structures must be within a data
  area which is assigned by OS-9 at run-time, and is separate from the
  program memory module.
\item
  All input and output operations should be made using OS-9 service
  request calls.
\end{enumerate}

Fortunately, the 6809\textquotesingle s versatile addressing modes make
the rules above easy to follow,. The OS-9 Assembler also helps because
it has special capabilities to assist the programmer in creating
programs and memory modules for the OS-9 execution environment.

\subsection{How to Write Position-Independent Code}

The 6809 instruction set was optimized to allow efficient use of
Position Independent Code (PIC). The basic technique is to always use
PC-relative addressing; for example BRA, LBRA, BSR and LBSR. Get
addresses of constants and tables using LEA instructions instead of load
immediate instructions. If you use dispatch tables, use tables of
RELATIVE, not absolute, addresses.

{\def\LTcaptype{none} % do not increment counter
\begin{longtable}[]{@{}ll@{}}
\toprule\noalign{}
INCORRECT & CORRECT \\
\midrule\noalign{}
\endhead
\bottomrule\noalign{}
\endlastfoot
LDX \#CONSTANT & LEAX CONSTANT,PCR \\
JSR SUBR & BSR SUBR or LBSR SUBR \\
JMP LABEL & BRA LABEL or LBRA LABEL \\
\end{longtable}
}

\subsection{Addressing Variables and Data Structures}

Programs executed as processes (by \hyperref[f.fork]{F\$Fork} and
\hyperref[f.chain]{F\$Chain} system calls or by the Shell) are assigned
a RAM memory area for variables, stacks, and data structures at
execution-time. The addresses cannot be determined or specified ahead of
time. However, a minimum size for this area is specified in the
program\textquotesingle s module header. Again, thanks to the
6809\textquotesingle s full compliment of addressing modes this presents
no problem to the OS-9 programmer.

When the program is first entered, the Y register will have the address
of the top of the process\textquotesingle{} data memory area. If the
creating process passed a parameter area, it will be located from the
value of the SP to the top of memory (Y), and the D register will
contain the parameter area size in bytes. If the new process was called
by the shell, the parameter area will contain the part of the shell
command line that includes the argument (parameter) text. The U register
will have the lower bound of the data memory area, and the DP register
will contain its page number.

The most important rule is to \emph{not use extended addressing!}
Indexed and direct page addressing should be used exclusively to access
data area values and structures. Do not use program-counter relative
addressing to find addresses in the data area, but do use it to refer to
addresses within the program area.

The most efficient way to handle tables, buffers, stacks, etc., is to
have the program\textquotesingle s initialization routine compute their
absolute addresses using the data area bounds passed by OS-9 in the
registers. These addresses can then be saved in the direct page where
they can be loaded into registers quickly, using short instructions.
This technique has advantages: it is faster than extended addressing,
and the program is inherently reentrant.

\subsection{Stack Requirements}

Because OS-9 uses interrupts extensively, and also because many
reentrant 6809 programs use the MPU stack for local variable storage, a
generous stack should be maintained at all times. The recommended
minimum is approximately 200 bytes.

\subsection{Interrupt Masks}

User programs should keep the condition codes register F (FIRQ mask) and
I (IRQ mask) bits off. They can be set during critical program sequences
to avoid task-switching or interrupts, but this time should be kept to a
minimum. If they are set for longer than a tick period, system
timekeeping accuracy may be affected. Also, some Level Two systems will
abort programs having a set IRQ mask.

\subsection{Writing Interrupt-driven Device Drivers}

OS-9 programs do not use interrupts directly. Any interrupt-driven
function should be implemented as a device driver module which should
handle all interrupt-related functions. When it is necessary for a
program to be synchronized to an interrupt-causing event, a driver can
send a semaphore to a program (or the reverse) using
OS-9\textquotesingle s \emph{signal} facilities.

It is important to understand that interrupt service routines are
asynchronous and somewhat nebulous in that they are not distinct
processes. They are in effect subroutines called by OS-9 when an
interrupt occurs.

Therefore, all interrupt-driven device drivers have two basic parts: the
``mainline'' subroutines that execute as part of the calling process,
and a separate interrupt service routine.

\emph{The two routines are asynchronous and therefore must use signals
for communications and coordination.}

The INIT initialization subroutine within the driver package should
allocate static storage for the service routine, get the service routine
address, and execute the \hyperref[f.irq]{F\$IRQ} system call to add it
to the IRQ polling table.

When a device driver routine does something that will result in an
interrupt, it should immediately execute a \hyperref[f.sleep]{F\$Sleep}
service request. This results in the process\textquotesingle{}
deactivation. When the interrupt in question occurs, its service routine
is executed after some random interval. It should then do the minimal
amount of processing required, and send a ``wakeup'' signal to its
associated process using the \hyperref[f.send]{F\$Send} service request.
It may also put some data in its static storage (I/O data and status)
which is shared with its associated ``sleeping'' process.

Some time later, the device driver ``mainline'' routine is awakened by
the signal, and can process the data or status returned by the interrupt
service routine.

\subsection{Using Standard I/O Paths}

Programs should be written to use standard I/O paths wherever practical.
Usually, this involves I/O calls that are intended to communicate to the
user\textquotesingle s terminal, or any other case where the OS-9
redirected I/O capability is desirable.

All three standard I/O paths will already be open when the program is
entered (they are inherited from the parent process). Programs should
\emph{not} close these paths except under very special circumstances.

Standard I/O paths are always assigned path numbers zero, one, and two,
as down below:

Path 0 - Standard Input. Analogous to the keyboard or other main data
input source.

Path 1 - Standard Output. Analogous to the terminal display or other
main data output destination.

Path 2 - Standard Error/Status. This path is provided so output messages
which are not part of the actual program output can be kept separate.
Many times paths 1 and 2 will be directed to the same device.

\subsection{A Sample Program}

The OS-9 \texttt{list} utility command program is shown on this and the
next page as an example of assembly language programming.

\begin{verbatim}
Microware OS-9 Assembler 2.1    01/04/82 23:39:37         Page 001
LIST - File List Utility


     *****
     *  LIST UTILITY COMMAND
     *  Syntax: list <pathname> 
     *  COPIES INPUT FROM SPECIFIED FILE TO STANDARD OUTPUT 
0000 87CD0048           mod LSTEND,LSTNAM,PRGRM+OBJCT,
                            REENT+1,LSTENT,LSTMEM
000D 4C6973F4   LSTNAM   fcs   "List"

     * STATIC STORAGE OFFSETS 
     *
00C8            BUFSIZ   equ   200        size of input buffer
0000                     ORG   0
0000            IPATH    rmb   1          input path number
0001            PRMPTR   rmb   2          parameter pointer
0003            BUFFER   rmb   BUFSIZ     allocate line buffer
00CB                     rmb   200        allocate stack
0193                     rmb   200        room for parameter list
025B            LSTMEM   EQU   .

0011 9F01       LSTENT   stx   PRMPTR     save parameter ptr
0013 8601                lda   #READ.     select read access mode
0015 103F84              os9   I$Open     open input file
0018 252E                bcs   LIST50     exit if error
001A 9700                sta   IPATH      save input path number
001C 9F01                stx   PRMPTR     save updated param ptr

001E 9600       LIST20   lda   IPATH      load input path number
0020 3043                leax  BUFFER,U   load buffer pointer
0022 10BE0C88            ldy   #BUFSIZ    maximum bytes to read
0026 103F8B              os9   I$ReadLn   read line of input
0029 2509                bcs   LIST30     exit if error
002B 8601                lda   #1         load std. out. path #
002D 103F8C              os9   I$WritLn   output line
0030 24EC                bcc   LIST20     Repeat if no error
0032 2014                bra   LIST50     exit if error

0034 C1D3       LIST30   cmpb  #E$EOF     at end of file?
0036 2610                bne   LIST50     branch if not
0038 9600                lda   IPATH      load input path number
003A 103F8F              os9   I$Close    close input path
003D 2509                bcs   LIST50     ..exit if error
003F 9E01                ldx   PRMPTR     restore parameter ptr
0041 A684                lda   0,X
0043 810D                cmpa  #$0D       End of parameter line?
0045 26CA                bne   LSTENT     ..no; list next file
0047 5F                  clrb
0048 103F06     LIST50   os9   F$Exit     ... terminate

004B 95BB58              emod             Module CRC

004E            LSTEND   EQU   *
\end{verbatim}

\section{Adapting OS-9 to a New System}

Thanks to OS-9\textquotesingle s modular structure, it is easily
portable to almost any 6809-based computer, and in fact it has been
installed on an incredible variety of hardware. Usually only device
driver and device descriptor modules need by rewritten or modified for
the target system\textquotesingle s specific hardware devices. The
larger and more complex kernel and file manager modules almost never
need adaptation.

One essential point is that you will need a functional OS-9 development
system to use during installation of OS-9 on a new target system.
Although it is possible to use a non-OS-9 system, or if you are truly
masochistic, the target system itself, lack of facilities to generate
and test memory modules and create system disks can make an otherwise
straightforward job a time-consuming headache that is seldom less costly
than a commercial OS-9-equipped computer. Over a dozen manufacturers
offer OS-9 based development systems in all price ranges with an
excellent selection of time-saving options such as hard disks, line
printers. PROM programmers, etc.

A group of OS-9 aficionados has over the years first written an
alternative to OS-9 called NitrOS9, then also disassembled all the
original Microware code and commented the source code. The software is
available on the Internet. You can find source code to the Kernel,
Shell, INIT, SYSGO, device driver and descriptor modules, and a
selection of utility commands which can be useful when moving OS-9 to a
new target system.

\subsection{Adapting OS-9 to Disk-based Systems}

Usually, most of the work in moving OS-9 to a disk-based target system
is writing a device driver module for the target
system\textquotesingle s disk controller. Part of this task involves
producing a subset of the driver (mostly disk read functions) for use as
a bootstrap module.

If terminal and/or parallel I/O for terminals, printers, etc., will use
ACIA and/or PIA-type devices, the standard ACIA and PIA device driver
modules may be used, or device drivers of your own design may be used in
place of or in addition to these standard modules Device descriptor
modules may also require adaptation to match device addresses and
initialization required by the target system.

A CLOCK module may be adapted from a standard version, or a new one may
be created. All other component modules, such as IOMAN, RBFMAN, SCFMAN,
Shell, and utilities seldom require modification.

\subsection{Using OS-9 in ROM-based Systems}

One of OS-9\textquotesingle s major features is its ability to reside in
ROM memory and work effectively with ROMed applications programs written
in assembler or high-level languages such as Basic09, Pascal, and C.

All the component modules of OS-9 (including all commands and utilities)
are directly ROMable without modification. In some cases, particularly
when the target system is to automatically execute an application
program upon system start-up, it may be necessary to reassemble the two
modules used during system startup, INIT and SYSGO.

The first step in designing a ROM-based system is to select which OS-9
modules to include in ROM. The following checklist is designed to help
you do so:

\begin{enumerate}
\def\labelenumi{\alph{enumi}.}
\item
  Include OS9P1, OS9P2, SYSGO, and INIT. These modules are required in
  any OS-9 system.
\item
  If the target system is perform any I/O or interrupt functions include
  IOMAN.
\item
  If the target system is to perform I/O to character-oriented I/O
  devices using ACIAs, PIAs, etc., include SCFMAN, required device
  drivers (such as ACIA and PIA, and/or your own), and device
  descriptors as needed (such as TERM, T1, P, and/or your own). If
  device addresses and/or initialization functions need to be changed,
  the device descriptor modules must be modified before being ROMed.
\item
  If the target system is to perform disk I/O, include RBFMAN, and
  appropriate disk driver and device descriptor modules. As in (c)
  above, change device addresses and initialization if needed. If RBFMAN
  \emph{will not} be included, the INIT and SYSGO modules \emph{must} be
  altered to remove references to disk files.
\item
  If the target system requires multiprogramming, time-of-day, or other
  time-related functions, include a CLOCK module for the target
  system\textquotesingle s real-time clock. Also consider how the clock
  is to be started,. You may want to ROM the \texttt{setime} command, or
  have SYSGO start the clock.
\item
  It the target system will receive commands manually, or if any
  application program uses Shell functions, include the SHELL and SYSGO
  modules, otherwise include a modified SYSGO module which calls your
  application program instead of Shell.
\end{enumerate}

\subsection{Adapting the Initialization Module}

INIT is a module that contains system startup parameters. It \emph{must}
be in ROM in any OS-9 system (it usually resides in the same ROM as the
kernel). It is a non-executable module named ``INIT'' and has type
``system'' (code \$C). It is scanned once during the system startup. It
begins with the standard header followed by:

{\def\LTcaptype{none} % do not increment counter
\begin{longtable}[]{@{}
  >{\raggedright\arraybackslash}p{(\linewidth - 2\tabcolsep) * \real{0.2200}}
  >{\raggedright\arraybackslash}p{(\linewidth - 2\tabcolsep) * \real{0.7800}}@{}}
\toprule\noalign{}
\begin{minipage}[b]{\linewidth}\raggedright
MODULE OFFSET
\end{minipage} & \begin{minipage}[b]{\linewidth}\raggedright
\end{minipage} \\
\midrule\noalign{}
\endhead
\bottomrule\noalign{}
\endlastfoot
\$9,\$A,\$B & This location contains an upper limit RAM memory address
used to override OS-9\textquotesingle s automatic end-of-RAM search so
that memory may be reserved for I/O device addresses or other special
purposes. \\
\$C & Number of entries to create in the IRQ polling table. One entry is
required for each interrupt- generating device control register. \\
\$D & Number of entries to create in the system device table. One entry
is required for each device in the system. \\
\$E,\$F & Offset to a string which is the name of the first module to be
executed after startup, usually ``SYSGO''. There must always be a
startup module. \\
\$10,\$11 & Offset to the default directory name string (normally /D0).
This device is assumed when device names are omitted from pathlists. If
the system will not use disks (e.g., RBFMAN will not be used) this
offset \emph{must}be zero. \\
\$12,\$13 & Offset to the initial standard path string (typically
/TERM). This path is opened as the standard paths for the initial
startup module. This offset \emph{must} contain zero if there is
none. \\
\$14,\$15 & Offset to bootstrap module name string. If OS-9 does not
find IOMAN in ROM during the start-up module search, it will execute the
bootstrap module named to load additional modules from a file on a
mass-storage device. \\
\$16 to N & All name strings referred to above go here. Each must have
the sign bit (bit 7) of the last character set. \\
\end{longtable}
}

\subsection{Adapting the SYSGO Module}

SYSGO is a program which is the first process started after the system
start-up sequence. Its function is threefold:

\begin{itemize}
\item
  It does additional high-level system initialization, for example, disk
  system SYSGO call the shell to process the \texttt{Startup} shell
  procedure file.
\item
  It starts the first ``user'' process.
\item
  It thereafter remains in a ``wait'' state as insurance against all
  user processes terminating, thus leaving the system halted. If this
  happens. SYSGO can restart the first user program.
\end{itemize}

The standard SYSGO module for disk systems cannot be used on non-disk
based systems unless it is modified to:

\begin{enumerate}
\def\labelenumi{\arabic{enumi}.}
\item
  Remove initialization of the working execution directory.
\item
  Remove processing of the \texttt{Startup} procedure file.
\item
  Possibly change the name of the first user program from Shell to the
  name of a applications program. Here are some example name strings:

  {\def\LTcaptype{none} % do not increment counter
  \begin{longtable}[]{@{}ll@{}}
  \toprule\noalign{}
  \endhead
  \bottomrule\noalign{}
  \endlastfoot
  fcs /userpqm/ & (object code module ``userpgm'') \\
  fcs /RunB userpgm/ & (compiled Basie09 program using RunB
  run-time-only system) \\
  fcs /Basic09 userpgm/ & (compiled Basic09 program using Basic09) \\
  \end{longtable}
  }
\end{enumerate}

\section{OS-9 Service Request Descriptions}

System calls are used to communicate between the OS-9 operating system
and assembly-language-level programs. There are three general
categories:

1.~User~mode~function~requests\\
2.~System~mode~function~requests\\
3.~I/O~requests\\

System mode function requests are privileged and may be executed only
while OS-9 is in the system state (when it is processinq another service
request, executing a file manager, device drivers, etc.). They are
included in this manual primarily for the benefit of those programmers
who will be writing device drivers and other system-level applications.

The system calls are performed by loading the MPU registers with the
appropriate parameters (if any), and executing a SWI2 instruction
immediately followed by a constant byte which is the request code.
Parameters (if any) will be returned in the MPU registers after OS-9 has
processed the service request. A standard convention for reporting
errors is used in all system calls; if an error occurred, the ``C bit''
of the condition code register will be set and accumulator B will
contain the appropriate error code. This permits a BCS or BCC
instruction immediately following the system call to branch on error/no
error.

Here is an example system call for the \hyperref[i.close]{I\$Close}
service request:

\begin{verbatim}
LDA PATHNUM 
SWI2 
FCB $8B 
BCS ERROR 
\end{verbatim}

Using the assembler\textquotesingle s ``OS9'' directive simplifies the
call:

\begin{verbatim}
LDA PATHNUM 
OS9 I$Close 
BCS ERROR 
\end{verbatim}

The I/O service requests are simpler to use than in many other operating
systems because the calling program does not have to allocate and set up
``file control blocks'', ``sector buffers'', etc. Instead OS-9 will
return a one byte path number when a path to a file/device is opened or
created; then this path number may be used in subsequent I/O requests to
identify the file/device until the path is closed. OS-9 internally
allocates and maintains its own data structures and users never have to
deal with them: in fact attempts to do so are memory violations.

All system calls have a mnemonic name that starts with ``F\$'' for
system functions, or ``I\$'' for I/O related requests. These are defined
in the assembler-input equate file called \texttt{OS9Defs}.

In the service request descriptions which follow, registers not
explicitly specified as input or output parameters are not altered.
Strings passed as parameters are normally terminated by having bit seven
of the last character set, a space character, or an end of line
character.

\subsection{F\$AllBit - Set bits in an allocation bit map}

{\def\LTcaptype{none} % do not increment counter
\begin{longtable}[]{@{}
  >{\raggedright\arraybackslash}p{(\linewidth - 2\tabcolsep) * \real{0.2800}}
  >{\raggedright\arraybackslash}p{(\linewidth - 2\tabcolsep) * \real{0.7200}}@{}}
\toprule\noalign{}
\endhead
\bottomrule\noalign{}
\endlastfoot
ASSEMBLER CALL: & OS9 F\$AllBit \\
MACHINE CODE: & 103F 13 \\
INPUT: & \begin{minipage}[t]{\linewidth}\raggedright
(X)~=~Base~address~of~allocation~bit~map.\\
(D)~=~Bit~number~of~first~bit~to~set.\\
(Y)~=~Bit~count~(number~of~bits~to~set)
\end{minipage} \\
OUTPUT: & None. \\
ERROR OUTPUT: & \begin{minipage}[t]{\linewidth}\raggedright
(CC)~=~C~bit~set.\\
(B)~=~Appropriate~error~code.
\end{minipage} \\
\end{longtable}
}

This system mode service request sets bits in the allocation bit map
specified by the X register.

Bit numbers range from 0..N-1, where N is the number of bits in the
allocation bit map.

\subsection{F\$Chain - Load and execute a new primary
module}\label{f.chain}

{\def\LTcaptype{none} % do not increment counter
\begin{longtable}[]{@{}
  >{\raggedright\arraybackslash}p{(\linewidth - 2\tabcolsep) * \real{0.2800}}
  >{\raggedright\arraybackslash}p{(\linewidth - 2\tabcolsep) * \real{0.7200}}@{}}
\toprule\noalign{}
\endhead
\bottomrule\noalign{}
\endlastfoot
ASSEMBLER CALL: & OS9 F\$Chain \\
MACHINE CODE: & 103F 05 \\
INPUT: & \begin{minipage}[t]{\linewidth}\raggedright
(X)~=~Address~of~module~name~or~file~name\\
(Y)~=~Parameter~area~size~(256~byte~pages)\\
(U)~=~Beginning~address~of~parameter~area\\
(A)~=~Lanquage~/~type~code\\
(B)~=~Optional~data~area~size~(256~byte~pages)
\end{minipage} \\
ERROR OUTPUT: & \begin{minipage}[t]{\linewidth}\raggedright
(CC)~=~C~bit~set.\\
(B)~=~Appropriate~error~code.
\end{minipage} \\
\end{longtable}
}

This system call is similar to \hyperref[f.fork]{F\$Fork}, but it does
not create a new process. It effectively ``resets'' the calling
process\textquotesingle{} program and data memory areas and begins
execution of a new primary module. Open paths are not closed or
otherwise affected.

This system call is used when it is necessary to execute an entirely new
program, but without the overhead of creating a new process. It is
functionally similar to a \hyperref[f.fork]{F\$Fork} followed by an
\hyperref[f.exit]{F\$Exit}, but with less processing overhead.

The sequence of operations taken by \hyperref[f.chain]{F\$Chain} is as
follows:

1. The system parses the name string ot the new proces\textquotesingle{}
``primary module'' - the program that will initially be executed. Then
the system module directory is searched to see if a module with the same
name and type / language is already in memory. If so it is linked to. If
not, the name string is used as the pathlist of a file which is to be
loaded into memory. Then the first module in this file is linked to
(several modules may have been loaded from a single file).

2. The process\textquotesingle{} old primary module is \emph{unlinked}.

3. The data memory area is reconfigured to the size specified in the new
primary module\textquotesingle s header.

The diagram below shows how \hyperref[f.chain]{F\$Chain} sets up the
data memory area and registers for the new module.

\begin{figure}
\centering
\begin{verbatim}
   +-----------------+  <--  Y          (highest address)
   !                 !
   !   Parameter     !
   !     Area        !
   !                 !
   +-----------------+  <-- X, SP
   !                 !
   !                 !
   !   Data Area     !
   !                 !
   !                 !
   +-----------------+
   !   Direct Page   !
   +-----------------+  <-- U, DP       (lowest address)

   D = parameter area size
  PC = module entry point abs. address
  CC = F=0, I=0, others undefined
\end{verbatim}
\caption{}
\end{figure}

Y (top of memory pointer) and U (bottom of memory pointer) will always
have a values at 256-byte page boundaries. If the parent does not
specify a parameter area, Y, X, and SP will be the same, and D will
equal zero. The minimum overall data area size is one page (256 bytes).

The hardware stack pointer (SP) should be located somewhere in the
direct page before the \hyperref[f.chain]{F\$Chain} service request is
executed to prevent a ``suicide attempt'' error or an actual suicide
(system crash). This will prevent a suicide from occurring in case the
new module requires a smaller data area than what is currently being
used. You should allow approximately 200 bytes of stack space for
execution of the \hyperref[f.chain]{F\$Chain} service request and other
system ``overhead''.

For more information, please see the \hyperref[f.fork]{F\$Fork} service
request description.

\subsection{F\$CmpNam - Compare two names}

{\def\LTcaptype{none} % do not increment counter
\begin{longtable}[]{@{}
  >{\raggedright\arraybackslash}p{(\linewidth - 2\tabcolsep) * \real{0.2800}}
  >{\raggedright\arraybackslash}p{(\linewidth - 2\tabcolsep) * \real{0.7200}}@{}}
\toprule\noalign{}
\endhead
\bottomrule\noalign{}
\endlastfoot
ASSEMBLER CALL: & OS9 F\$CmpNam \\
MACHINE CODE: & 103F 11 \\
INPUT: & \begin{minipage}[t]{\linewidth}\raggedright
(X)~=~Address~of~first~name.\\
(B)~=~Length~of~first~name.\\
(Y)~=~Address~of~second~name.
\end{minipage} \\
OUTPUT: & (CC) = C bit clear if the strings match. \\
\end{longtable}
}

Given the address and length of a string, and the address of a second
string, compares them and indicates whether they match. Typically used
in conjunction with ``parsename''.

The second name must bave the sign bit (bit 7) of the last character
set.

\subsection{F\$CRC - Compute CRC}\label{f.crc}

{\def\LTcaptype{none} % do not increment counter
\begin{longtable}[]{@{}
  >{\raggedright\arraybackslash}p{(\linewidth - 2\tabcolsep) * \real{0.2800}}
  >{\raggedright\arraybackslash}p{(\linewidth - 2\tabcolsep) * \real{0.7200}}@{}}
\toprule\noalign{}
\endhead
\bottomrule\noalign{}
\endlastfoot
ASSEMBLER CALL: & OS9 F\$CRC \\
MACHINE CODE: & 103F 17 \\
INPUT: & \begin{minipage}[t]{\linewidth}\raggedright
(X)~=~Starting~byte~address.\\
(Y)~=~Byte~count.\\
(U)~=~Address~of~3~byte~CRC~accumulator.
\end{minipage} \\
OUTPUT: & CRC accumulator is updated. \\
ERROR OUTPUT: & None. \\
\end{longtable}
}

This service request calculates the CRC (cyclic redundancy count) for
use by compilers, assemblers, or other module generators. The CRC is
calculated starting at the source address over ``byte count'' bytes, it
is not necessary to cover an entire module in one call, since the CRC
may be ``accumulated'' over several calls. The CRC accumulator can be
any three byte memory location and must be initialized to \$FFFFFF
before the first F\$CRC call.

The last three bytes in the module (where the three CRC bytes will be
stored) are not included in the CRC generation.

The polynomial is \$800063. If you perform the CRC over the module minus
the stored CRC then you must exclusive-or the result with \$FFFFFF to
compare directly with the stored CRC. If you perform CRC over the module
including the last three bytes then the result must be \$800FE3.

Example C code of CRC algorithm (with 32-bit longs):

\begin{verbatim}
unsigned long compute_crc()
    unsigned long crc;
    unsigned char *octets;
    int len;
{
    int i;

    while (len--) {
        crc ^= (*octets++) << 16;
        for (i = 0; i < 8; i++) {
            crc <<= 1;
            if (crc & 0x1000000L)
                crc ^= 0x800063L;
        }
    }
    return crc & 0xffffffL;
}
\end{verbatim}

\subsection{F\$DelBit - Deallocate in a bit map}

{\def\LTcaptype{none} % do not increment counter
\begin{longtable}[]{@{}
  >{\raggedright\arraybackslash}p{(\linewidth - 2\tabcolsep) * \real{0.2800}}
  >{\raggedright\arraybackslash}p{(\linewidth - 2\tabcolsep) * \real{0.7200}}@{}}
\toprule\noalign{}
\endhead
\bottomrule\noalign{}
\endlastfoot
ASSEMBLER CALL: & OS9 F\$DelBit \\
MACHINE CODE: & 103F 14 \\
INPUT: & \begin{minipage}[t]{\linewidth}\raggedright
(X)~=~Base~address~of~an~allocation~bit~map.\\
(D)~=~Bit~number~of~first~bit~to~clear.\\
(Y)~=~Bit~count~(number~of~bits~to~clear).
\end{minipage} \\
OUTPUT: & None. \\
ERROR OUTPUT: & \begin{minipage}[t]{\linewidth}\raggedright
(CC)~=~C~bit~set.\\
(B)~=~Appropriate~error~code.
\end{minipage} \\
\end{longtable}
}

This system mode service request is used to clear bits in the allocation
bit map pointed to by X.

Bit numbers range from 0..N-1, where N is the number of bits in the
allocation bit map.

\subsection{F\$Exit - Terminate the calling process}\label{f.exit}

{\def\LTcaptype{none} % do not increment counter
\begin{longtable}[]{@{}
  >{\raggedright\arraybackslash}p{(\linewidth - 2\tabcolsep) * \real{0.2800}}
  >{\raggedright\arraybackslash}p{(\linewidth - 2\tabcolsep) * \real{0.7200}}@{}}
\toprule\noalign{}
\endhead
\bottomrule\noalign{}
\endlastfoot
ASSEMBLER CALL: & OS9 F\$Exit \\
MACHINE CODE: & 103F 06 \\
INPUT: & (B) = Status code to be returned to the parent process \\
OUTPUT: & Process is terminated. \\
\end{longtable}
}

This call kills the calling process and is the only means by which a
process can terminate itself. Its data memory area is deallocated, and
its primary module is UNLINKed. All open paths are automatically closed.

The death of the process can be detected by the parent executing a
\hyperref[f.wait]{F\$Wait} call, which returns to the parent the status
byte passed by the child in its \hyperref[f.exit]{F\$Exit} call. The
status byte can be an OS-9 error code the terminating process wishes to
pass back to its parent process (the shell assumes this), or can be used
to pass a user-defined status value. Processes to be called directly by
the shell should only return an OS-9 error code or zero if no error
occurred.

\subsection{F\$Fork - Create a new process}\label{f.fork}

{\def\LTcaptype{none} % do not increment counter
\begin{longtable}[]{@{}
  >{\raggedright\arraybackslash}p{(\linewidth - 2\tabcolsep) * \real{0.2800}}
  >{\raggedright\arraybackslash}p{(\linewidth - 2\tabcolsep) * \real{0.7200}}@{}}
\toprule\noalign{}
\endhead
\bottomrule\noalign{}
\endlastfoot
ASSEMBLER CALL: & OS9 F\$Fork \\
MACHINE CODE: & 103F 03 \\
INPUT: & \begin{minipage}[t]{\linewidth}\raggedright
(X)~=~Address~of~module~name~or~file~name.\\
(Y)~=~Parameter~area~size.\\
(U)~=~Beginning~address~of~the~parameter~area.\\
(A)~=~Language~/~Type~code.\\
(B)~=~Optional~data~area~size~(pages).
\end{minipage} \\
OUTPUT: & \begin{minipage}[t]{\linewidth}\raggedright
(X)~=~Updated~past~the~name~string.\\
(A)~=~New~process~ID~number.
\end{minipage} \\
ERROR OUTPUT: & \begin{minipage}[t]{\linewidth}\raggedright
(CC)~=~C~bit~set.\\
(B)~=~Appropriate~error~code.
\end{minipage} \\
\end{longtable}
}

This system call creates a new process which becomes a ``child'' of the
caller, and sets up the new process\textquotesingle{} memory and MPU
registers.

The system parses the name string of the new process\textquotesingle{}
``primary module'' - the program that will initially be executed. Then
the system module directory is searched to see if the program is already
in memory. If so, the module is linked to and executed. If not, the name
string is used as the pathlist of the file which is to loaded into
memory. Then the first module in this file is linked to and executed
(several modules may have been loaded from a single file).

The primary module\textquotesingle s module header is used to determine
the process\textquotesingle{} initial data area size. OS-9 then attempts
to allocate a contiguous RAM area equal to the required data storage
size, (includes the parameter passing area, which is copied from the
parent process\textquotesingle{} data area). The new
process\textquotesingle{} registers are set up as shown in the diagram
on the next page. The execution offset given in the module header is
used to set the PC to the module\textquotesingle s entry point.

When the shell processes a command line it passes a string in the
parameter area which is a copy of the parameter part (if any) of the
command line. It also inserts an end-of-line character at the end of the
parameter string to simplify string-oriented processing. The X register
will point to the beginning of the parameter string. If the command line
included the optional memory size specification (\#n or \#nK), the shell
will pass that size as the requested memory size when executing the
\hyperref[f.fork]{F\$Fork}.

If any of the above operations are unsuccessful, the
\hyperref[f.fork]{F\$Fork} is aborted and the caller is returned an
error.

The diagram below shows how \hyperref[f.fork]{F\$Fork} sets up the data
memory area and registers for a newly-created process.

\begin{figure}
\centering
\begin{verbatim}
   +-----------------+  <--  Y          (highest address)
   !                 !
   !   Parameter     !
   !     Area        !
   !                 !
   +-----------------+  <-- X, SP
   !                 !
   !                 !
   !   Data Area     !
   !                 !
   !                 !
   +-----------------+
   !   Direct Page   !
   +-----------------+  <-- U, DP       (lowest address)

   D = parameter area size
  PC = module entry point abs. address
  CC = F=0, I=0, others undefined
\end{verbatim}
\caption{}
\end{figure}

Y (top of memory pointer) and U (bottom of memory pointer) will always
have a values at 256-byte page boundaries. If the parent does not
specify a parameter area, Y, X, and SP will be the same, and D will
equal zero. The minimum overall data area size is one page (256 bytes).
Shell will always pass at least an end of line character in the
parameter area.

NOTE: Both the child and parent process will execute concurrently. If
the parent executes a \hyperref[f.wait]{F\$Wait} call immediately after
the fork, it will wait until the child dies before it resumes execution.
Caution should be exercised when recursively calling a program that uses
the \hyperref[f.fork]{F\$Fork} service request since another child may
be created with each ``incarnation''. This will continue until the
process table becomes full.

\subsection{F\$ICPT - Set up a signal intercept trap}\label{f.icpt}

{\def\LTcaptype{none} % do not increment counter
\begin{longtable}[]{@{}
  >{\raggedright\arraybackslash}p{(\linewidth - 2\tabcolsep) * \real{0.2800}}
  >{\raggedright\arraybackslash}p{(\linewidth - 2\tabcolsep) * \real{0.7200}}@{}}
\toprule\noalign{}
\endhead
\bottomrule\noalign{}
\endlastfoot
ASSEMBLER CALL: & OS9 F\$ICPT \\
MACHINE CODE: & 103F 09 \\
INPUT: & \begin{minipage}[t]{\linewidth}\raggedright
(X)~=~Address~of~the~intercept~routine.\\
(U)~=~Address~of~the~intercept~routine~local~storage.
\end{minipage} \\
OUTPUT: & None. \\
ERROR OUTPUT: & \begin{minipage}[t]{\linewidth}\raggedright
(CC)~=~C~bit~set.\\
(B)~=~Appropriate~error~code.
\end{minipage} \\
\end{longtable}
}

This system call tells OS-9 to set a signal intercept trap, where X
contains the address of the signal handler routine, and U contains the
base address of the routine\textquotesingle s storage area. After a
signal trap has been set, whenever the process receives a signal, its
intercept routine will be executed. A signal will abort any process
which has not used the \hyperref[f.icpt]{F\$ICPT} service request to set
a signal trap, and its termination status (B register) will be the
signal code. Many interactive programs will set up an intercept routine
to handle keyboard abort ({ +controlQ}), and keyboard interrupt ({
+controlC}).

The intercept routine is entered asynchronously because a signal may be
sent at any time (it is like an interrupt) and is passed the following:

U = Address of intercept routine local storage.

B = Signal code.

NOTE: The value of DP may not be the same as it was when the
\hyperref[f.icpt]{F\$ICPT} call was made.

Whenever a signal is received. OS-9 will pass the signal code and the
base address of its data area (which was defined by a
\hyperref[f.icpt]{F\$ICPT} service request) to the signal intercept
routine. The base address of the data area is selected by the user and
is typically a pointer to the process\textquotesingle{} data area.

The intercept routine is activated when a signal is received, then it
takes some action based upon the value of the signal code such as
setting a flag in the process\textquotesingle{} data area. After the
signal has been processed, the handler routine should terminate with an
RTI instruction.

\subsection{F\$ID - Get process ID / user ID}

{\def\LTcaptype{none} % do not increment counter
\begin{longtable}[]{@{}
  >{\raggedright\arraybackslash}p{(\linewidth - 2\tabcolsep) * \real{0.2800}}
  >{\raggedright\arraybackslash}p{(\linewidth - 2\tabcolsep) * \real{0.7200}}@{}}
\toprule\noalign{}
\endhead
\bottomrule\noalign{}
\endlastfoot
ASSEMBLER CALL: & OS9 F\$ID \\
MACHINE CODE: & 103F 0C \\
INPUT: & None \\
OUTPUT: & \begin{minipage}[t]{\linewidth}\raggedright
(A)~=~Process~ID.\\
(Y)~=~User~ID.
\end{minipage} \\
ERROR OUTPUT: & \begin{minipage}[t]{\linewidth}\raggedright
(CC)~=~C~bit~set.\\
(B)~=~Appropriate~error~code.
\end{minipage} \\
\end{longtable}
}

Returns the caller\textquotesingle s process ID number, which is a byte
value in the range of 1 to 255, and the user ID which is a integer in
the range 0 to 65535. The process ID is assigned by OS-9 and is unique
to the process. The user ID is defined in the system password file, and
is used by the file security system and a few other functions. Several
processes can have the same user ID.

\subsection{F\$LINK - Link to memory module}\label{f.link}

{\def\LTcaptype{none} % do not increment counter
\begin{longtable}[]{@{}
  >{\raggedright\arraybackslash}p{(\linewidth - 2\tabcolsep) * \real{0.2800}}
  >{\raggedright\arraybackslash}p{(\linewidth - 2\tabcolsep) * \real{0.7200}}@{}}
\toprule\noalign{}
\endhead
\bottomrule\noalign{}
\endlastfoot
ASSEMBLER CALL: & OS9 F\$LINK \\
MACHINE CODE: & 103F 00 \\
INPUT: & \begin{minipage}[t]{\linewidth}\raggedright
(X)~=~Address~of~the~module~name~string.\\
(A)~=~Module~type~/~language~byte.
\end{minipage} \\
OUTPUT: & \begin{minipage}[t]{\linewidth}\raggedright
(X)~=~Advanced~past~the~module~name.\\
(Y)~=~Module~entry~point~absolute~address.\\
(U)~=~Module~header~absolute~address.\\
(A)~=~Module~type~/~language.\\
(B)~=~Module~attributes~/~revision~level.
\end{minipage} \\
ERROR OUTPUT: & \begin{minipage}[t]{\linewidth}\raggedright
(CC)~=~C~bit~set.\\
(B)~=~Appropriate~error~code.
\end{minipage} \\
\end{longtable}
}

This system call causes OS-9 to search the module directory for a module
having a name, language and type as given in the parameters. If found,
the address of the module\textquotesingle s header is returned in U, and
the absolute address of the module\textquotesingle s execution entry
point is returned in Y (as a convenience: this and other information can
be obtained from the module header). The module\textquotesingle s link
count\textquotesingle{} is incremented whenever a
\hyperref[f.link]{F\$LINK} references its name, thus keeping track of
how many processes are using the module. If the module requested has an
attribute byte indicating it is not sharable (meaning it is not
reentrant) only one process may link to it at a time.

Possible errors:

(A) Module not found.

(B) Module busy (not sharable and in use).

(C) Incorrect or defective module header.

\subsection{F\$LOAD - Load module(s) from a file}

{\def\LTcaptype{none} % do not increment counter
\begin{longtable}[]{@{}
  >{\raggedright\arraybackslash}p{(\linewidth - 2\tabcolsep) * \real{0.2800}}
  >{\raggedright\arraybackslash}p{(\linewidth - 2\tabcolsep) * \real{0.7200}}@{}}
\toprule\noalign{}
\endhead
\bottomrule\noalign{}
\endlastfoot
ASSEMBLER CALL: & OS9 F\$LOAD \\
MACHINE CODE: & 103F 01 \\
INPUT: & \begin{minipage}[t]{\linewidth}\raggedright
(X)~=~Address~of~pathlist~(file~name)\\
(A)~=~Language~/~type~(0~=~any~language~/~type)
\end{minipage} \\
OUTPUT: & \begin{minipage}[t]{\linewidth}\raggedright
(X)~=~Advanced~past~pathlist\\
(Y)~=~Primary~module~entry~point~address\\
(U)~=~Address~of~module~header\\
(A;~-~Language~/~type\\
(B)~=~Attributes~/~revision~level
\end{minipage} \\
ERROR OUTPUT: & \begin{minipage}[t]{\linewidth}\raggedright
(CC)~=~C~bit~set.\\
(B)~=~Appropriate~error~code.
\end{minipage} \\
\end{longtable}
}

Opens a file specified by the pathlist, reads one or more memory modules
from the file into memory, then closes the file. All modules loaded are
added to the system module directory, and the first module read is
LINKed. The parameters returned are the same as the
\hyperref[f.link]{F\$LINK} call and apply only to the first module
loaded.

In order to be loaded, the file must have the ``execute'' permission and
contain a module or modules that have a proper module header. The file
will be loaded from the working execution directory unless a complete
pathlist is given.

Possible errors: module directory full; memory full; plus errors that
occur on \hyperref[i.open]{I\$Open}, \hyperref[i.read]{I\$Read},
\hyperref[i.close]{I\$Close} and \hyperref[f.link]{F\$LINK} system
calls.

\subsection{F\$Mem - Resize data memory area}\label{f.mem}

{\def\LTcaptype{none} % do not increment counter
\begin{longtable}[]{@{}
  >{\raggedright\arraybackslash}p{(\linewidth - 2\tabcolsep) * \real{0.2800}}
  >{\raggedright\arraybackslash}p{(\linewidth - 2\tabcolsep) * \real{0.7200}}@{}}
\toprule\noalign{}
\endhead
\bottomrule\noalign{}
\endlastfoot
ASSEMBLER CALL: & OS9 F\$Mem \\
MACHINE CODE: & 103F 07 \\
INPUT: & (D) = Desired new memory area size in bytes \\
OUTPUT: & \begin{minipage}[t]{\linewidth}\raggedright
(Y)~=~Address~of~new~memory~area~upper~bound\\
(D)~=~Actual~new~memory~area~size~in~bytes
\end{minipage} \\
ERROR OUTPUT: & \begin{minipage}[t]{\linewidth}\raggedright
(CC)~=~C~bit~set.\\
(B)~=~Appropriate~error~code.
\end{minipage} \\
\end{longtable}
}

Used to expand or contract the process\textquotesingle{} data memory
area. The new size requested is rounded up to the next 256-byte page
boundary. Additional memory is allocated contiguously upward (towards
higher addresses), or deallocated downward from the old highest address.
If D = 0, then the current upper bound and size will be returned.

This request can never return all of a process\textquotesingle{} memory,
or the page in which its SP register points to.

In Level One systems, the request may return an error upon an expansion
request even though adequate free memory exists. This is because the
data area is always made contiguous, and memory requests by other
processes may fragment free memory into smaller, scattered blocks that
are not adjacent to the caller\textquotesingle s present data area.
Level Two systems do not have this restriction because of the
availability of hardware for memory relocation, and because each process
has its own ``address space''.

\subsection{F\$PErr - Print error message}

{\def\LTcaptype{none} % do not increment counter
\begin{longtable}[]{@{}
  >{\raggedright\arraybackslash}p{(\linewidth - 2\tabcolsep) * \real{0.2800}}
  >{\raggedright\arraybackslash}p{(\linewidth - 2\tabcolsep) * \real{0.7200}}@{}}
\toprule\noalign{}
\endhead
\bottomrule\noalign{}
\endlastfoot
ASSEMBLER CALL: & OS9 F\$PErr \\
MACHINE CODE: & 103F 0F \\
INPUT: & \begin{minipage}[t]{\linewidth}\raggedright
(A)~=~Output~path~number.\\
(B)~=~Error~code.
\end{minipage} \\
OUTPUT: & None. \\
ERROR OUTPUT: & \begin{minipage}[t]{\linewidth}\raggedright
(CC)~=~C~bit~set.\\
(B)~=~Appropriate~error~code.
\end{minipage} \\
\end{longtable}
}

This is the system\textquotesingle s error reporting utility. It writes
an error message to the output path specified. Most OS-9 systems will
display:

ERROR \#\textless decimal number\textgreater{}

by default. The error reporting routine is vectored and can be replaced
with a more elaborate reporting module. To replace this routine use the
\hyperref[f.ssvc]{F\$SSVC} service request.

\subsection{F\$PrsNam - Parse a path name}

{\def\LTcaptype{none} % do not increment counter
\begin{longtable}[]{@{}
  >{\raggedright\arraybackslash}p{(\linewidth - 2\tabcolsep) * \real{0.2800}}
  >{\raggedright\arraybackslash}p{(\linewidth - 2\tabcolsep) * \real{0.7200}}@{}}
\toprule\noalign{}
\endhead
\bottomrule\noalign{}
\endlastfoot
ASSEMBLER CALL: & OS9 F\$PrsNam \\
MACHINE CODE: & 103F 10 \\
INPUT: & (X) = Address of the pathlist \\
OUTPUT: & \begin{minipage}[t]{\linewidth}\raggedright
(X)~=~Updated~past~the~optional~``/''\\
(Y)~=~Address~of~the~last~character~of~the~name~+~1.\\
(B)~=~Length~of~the~name\\
\end{minipage} \\
ERROR OUTPUT: & \begin{minipage}[t]{\linewidth}\raggedright
(CC)~=~C~bit~set.\\
(B)~=~Appropriate~error~code.\\
(X)~=~Updated~past~space~characters.
\end{minipage} \\
\end{longtable}
}

Parses the input text string for a legal OS-9 name. The name is
terminated by any character that is not a legal component character.
This system call is useful for processing pathlist arguments passed to
new processes. Also if X was at the end of a pathlist, a bad name error
will be returned and X will be moved past any space characters so that
the next pathlist in a command line may be parsed.

Note that this system call processes only one name, so several calls may
be needed to process a pathlist that has more than one name.

BEFORE F\$PrsNam CALL:

\begin{verbatim}
+---+---+---+---+---+---+---+---+---+---+---+---+---
! / ! D ! 0 ! / ! F ! I ! L ! E !   !   !   !   !
+---+---+---+---+---+---+---+---+---+---+---+---+---
  ^
  X
\end{verbatim}

AFTER THE F\$PrsNam CALL:

\begin{verbatim}
+---+---+---+---+---+---+---+---+---+---+---+---+---
! / ! D ! 0 ! / ! F ! I ! L ! E !   !   !   !   !
+---+---+---+---+---+---+---+---+---+---+---+---+---
      ^       ^
      X       Y       (B) = 2
\end{verbatim}

\subsection{F\$SchBit - Search bit map for a free area}

{\def\LTcaptype{none} % do not increment counter
\begin{longtable}[]{@{}
  >{\raggedright\arraybackslash}p{(\linewidth - 2\tabcolsep) * \real{0.2800}}
  >{\raggedright\arraybackslash}p{(\linewidth - 2\tabcolsep) * \real{0.7200}}@{}}
\toprule\noalign{}
\endhead
\bottomrule\noalign{}
\endlastfoot
ASSEMBLER CALL: & OS9 F\$SchBit \\
MACHINE CODE: & 103F 12 \\
INPUT: & \begin{minipage}[t]{\linewidth}\raggedright
(X)~=~Beginning~address~of~a~bit~map.\\
(D)~=~Beginning~bit~number.\\
(Y)~=~Bit~count~(free~bit~block~size).\\
(U)~=~End~of~bit~map~address.
\end{minipage} \\
OUTPUT: & \begin{minipage}[t]{\linewidth}\raggedright
(D)~=~Beginning~bit~number.\\
(Y)~=~Bit~count.
\end{minipage} \\
\end{longtable}
}

This system mode service request searches the specified allocation bit
map starting at the ``beginning bit number'' for a free block (cleared
bits) of the required length.

If no block of the specified size exists, it returns with the carry set,
beginning bit number and size of the largest block.

\subsection{F\$Send - Send a signal to another process}\label{f.send}

{\def\LTcaptype{none} % do not increment counter
\begin{longtable}[]{@{}
  >{\raggedright\arraybackslash}p{(\linewidth - 2\tabcolsep) * \real{0.2800}}
  >{\raggedright\arraybackslash}p{(\linewidth - 2\tabcolsep) * \real{0.7200}}@{}}
\toprule\noalign{}
\endhead
\bottomrule\noalign{}
\endlastfoot
ASSEMBLER CALL: & OS9 F\$Send \\
MACHINE CODE: & 103F 08 \\
INPUT: & \begin{minipage}[t]{\linewidth}\raggedright
(A)~=~Receiver\textquotesingle s~process~ID~number.\\
\strut ~(B)~=~Signal~code.
\end{minipage} \\
OUTPUT: & None. \\
ERROR OUTPUT: & \begin{minipage}[t]{\linewidth}\raggedright
(CC)~=~C~bit~set.\\
(B)~=~Appropriate~error~code.
\end{minipage} \\
\end{longtable}
}

This system call sends a ``signal'' to the process specified. The signal
code is a single byte value of 1 - 255.

If the signal\textquotesingle s destination process is sleeping or
waiting, it will be activated so that it may process the signal. The
signal processing routine (intercept) will be executed if a signal trap
was set up (see \hyperref[f.icpt]{F\$ICPT}), otherwise the signal will
abort the destination process, and the signal code becomes the exit
status (see \hyperref[f.wait]{F\$Wait}). An exception is the WAKEUP
signal, which activates a sleeping process but does not cause the signal
intercept routine to be executed.

Some of the signal codes have meanings defined by convention:

0~=~System~Abort~(cannot~be~intercepted)\\
1~=~Wake~Up~Process\\
2~=~Keyboard~Abort\\
3~=~Keyboard~Interrupt\\
4-255~=~user~defined\\

If an attempt is made to send a signal to a process that has an
unprocessed, previous signal pending, the current ``send'' request will
be cancelled and an error will be returned. An attempt can be made to
re-send the signal later. It is good practice to issue a ``sleep'' call
for a few ticks before a retry to avoid wasting MPU time.

For related information see the \hyperref[f.icpt]{F\$ICPT},
\hyperref[f.wait]{F\$Wait} and \hyperref[f.sleep]{F\$Sleep} service
request descriptions.

\subsection{F\$Sleep - Put calling process to sleep}\label{f.sleep}

{\def\LTcaptype{none} % do not increment counter
\begin{longtable}[]{@{}
  >{\raggedright\arraybackslash}p{(\linewidth - 2\tabcolsep) * \real{0.2800}}
  >{\raggedright\arraybackslash}p{(\linewidth - 2\tabcolsep) * \real{0.7200}}@{}}
\toprule\noalign{}
\endhead
\bottomrule\noalign{}
\endlastfoot
ASSEMBLER CALL: & OS9 F\$Sleep \\
MACHINE CODE: & 103F 0A \\
INPUT: & (X) = Sleep time in ticks (0 = indefinitely) \\
OUTPUT: & (X) = Decremented by the number of ticks that the process was
asleep. \\
ERROR OUTPUT: & \begin{minipage}[t]{\linewidth}\raggedright
(CC)~=~C~bit~set.\\
(B)~=~Appropriate~error~code.
\end{minipage} \\
\end{longtable}
}

This call deactivates the calling process for a specified time, or
indefinitely if X = 0. If X = 1, the effect is to have the caller give
up its current time slice. The process will be activated before the full
time interval if a signal is received, therefore sleeping indefinitely
is a good way to wait for a signal or interrupt without wasting CPU
time.

The duration of a ``tick'' is system dependent but is most commonly 100
milliseconds.

Due to the fact that it is not known when the
\hyperref[f.sleep]{F\$Sleep} request was made during the current tick,
\hyperref[f.sleep]{F\$Sleep} can not be used for precise timing. A sleep
of one tick is effectively a ``give up remaining time slice'' request;
the process is immediately inserted into the active process queue and
will resume execution when it reaches the front of the queue. A sleep of
two or more ticks causes the process to be inserted into the active
process queue after N-1 ticks occur and will resume execution when it
reaches the front of the queue.

\subsection{F\$SPrior - Set process priority}

{\def\LTcaptype{none} % do not increment counter
\begin{longtable}[]{@{}
  >{\raggedright\arraybackslash}p{(\linewidth - 2\tabcolsep) * \real{0.2800}}
  >{\raggedright\arraybackslash}p{(\linewidth - 2\tabcolsep) * \real{0.7200}}@{}}
\toprule\noalign{}
\endhead
\bottomrule\noalign{}
\endlastfoot
ASSEMBLER CALL: & OS9 F\$SPrior \\
MACHINE CODE: & 103F 0D \\
INPUT: & \begin{minipage}[t]{\linewidth}\raggedright
(A)~=~Process~ID~number.\\
(B)~=~Priority:\\
\strut ~~~~~~~0~=~lowest\\
\strut ~~~~~~~255~-~highest
\end{minipage} \\
OUTPUT: & None. \\
ERROR OUTPUT: & \begin{minipage}[t]{\linewidth}\raggedright
(CC)~=~C~bit~set.\\
(B)~=~Appropriate~error~code.
\end{minipage} \\
\end{longtable}
}

Changes the process\textquotesingle{} priority to the new value given.
\$FF is the highest possible priority, \$00 is the lowest. A process can
change another process\textquotesingle{} priority only if it has the
same user ID.

\subsection{F\$SSVC - Install function request}\label{f.ssvc}

{\def\LTcaptype{none} % do not increment counter
\begin{longtable}[]{@{}
  >{\raggedright\arraybackslash}p{(\linewidth - 2\tabcolsep) * \real{0.2800}}
  >{\raggedright\arraybackslash}p{(\linewidth - 2\tabcolsep) * \real{0.7200}}@{}}
\toprule\noalign{}
\endhead
\bottomrule\noalign{}
\endlastfoot
ASSEMBLER CALL: & OS9 F\$SSVC \\
MACHINE CODE: & 103F 32 \\
INPUT: & (Y) = Address of service request initialization table. \\
OUTPUT: & None. \\
ERROR OUTPUT: & \begin{minipage}[t]{\linewidth}\raggedright
(CC)~=~C~bit~set.\\
(B)~=~Appropriate~error~code.
\end{minipage} \\
\end{longtable}
}

This system mode service request is used to add a new function request
to OS-9\textquotesingle s user and privileged system service request
tables, or to replace an old one. The Y register passes the address of a
table which contains the function codes and offsets to the corresponding
service request handler routines. This table has the following format:

\begin{verbatim}
OFFSET

          +----------------------+
 $00      !     Function Code    !  <--- First entry
          +----------------------+
 $01      ! Offset From Byte 3   !
          +--                  --+
 $02      !  To Function Handler !
          +----------------------+
 $03      !     Function Code    !  <--- Second entry
          +----------------------+
 $04      ! Offset From Byte 6   !
          +--                  --+
 $05      !  To Function Handler !
          +----------------------+
          !                      !   <--- Third entry etc.
          !     MORE ENTRIES     !
          !                      !
          !                      !
          +----------------------+
          !         $80          !   <--- End of table mark
          +----------------------+
\end{verbatim}

NOTE: If the sign bit of the function code is set, only the system table
will be updated. Otherwise both the system and user tables will be
updated. Privileged system service requests may be called only while
executing a system routine.

The service request handler routine should process the service request
and return from subroutine with an RTS instruction. They may alter all
MPU registers (except for SP). The U register will pass the address of
the register stack to the service request handler as shown in the
following diagram:

\begin{figure}
\centering
\begin{verbatim}
                         OFFSET   OS9Defs
                                  MNEMONIC
        +------+
U --->  !  CC  !            $0      R$CC
        +------+            $1      R$D
        !  A   !            $1      R$A
        +------+
        !  B   !            $2      R$B
        +------+
        !  DP  !            $3      R$DP
        +------+------+
        !      X      !     $4      R$X
        +-------------+
        !      Y      !     $6      R$Y
        +-------------+
        !      U      !     $8      R$U
        +-------------+
        !      PC     !     $A      R$PC
        +-------------+
\end{verbatim}
\caption{}
\end{figure}

Function request codes are broken into the two categories as shown
below:

{\def\LTcaptype{none} % do not increment counter
\begin{longtable}[]{@{}
  >{\raggedright\arraybackslash}p{(\linewidth - 2\tabcolsep) * \real{0.3000}}
  >{\raggedright\arraybackslash}p{(\linewidth - 2\tabcolsep) * \real{0.7000}}@{}}
\toprule\noalign{}
\endhead
\bottomrule\noalign{}
\endlastfoot
\$00 - \$27 & User mode service request codes. \\
\$29 - \$34 & Privileged system mode service request codes. When
installing these service request, the sign bit should be set if it is to
be placed into the system table only. \\
\end{longtable}
}

NOTE: These categories are defined by convention and not enforced by
OS9.

Codes \$25..\$27, and \$70..\$7F will not be used by the operating
system and are free for user definition.

\subsection{F\$SSWI - Set SWI vector}\label{f.sswi}

{\def\LTcaptype{none} % do not increment counter
\begin{longtable}[]{@{}
  >{\raggedright\arraybackslash}p{(\linewidth - 2\tabcolsep) * \real{0.2800}}
  >{\raggedright\arraybackslash}p{(\linewidth - 2\tabcolsep) * \real{0.7200}}@{}}
\toprule\noalign{}
\endhead
\bottomrule\noalign{}
\endlastfoot
ASSEMBLER CALL: & OS9 F\$SSWI \\
MACHINE CODE: & 103F 0E \\
INPUT: & \begin{minipage}[t]{\linewidth}\raggedright
(A)~=~SWI~type~code.\\
(X)~=~Address~of~user~SWI~service~routine.
\end{minipage} \\
OUTPUT: & None. \\
ERROR OUTPUT: & \begin{minipage}[t]{\linewidth}\raggedright
(CC)~=~C~bit~set.\\
(B)~=~Appropriate~error~code.
\end{minipage} \\
\end{longtable}
}

Sets up the interrupt vectors for SWI, SWI2 and SWI3 instructions. Each
process has its own local vectors. Each F\$SSWI call sets up one type of
vector according to the code number passed in A.

1~=~SWI~\\
2~=~SWI2~\\
3~=~SWI3~\\

When a process is created, all three vectors are initialized with the
address of the OS-9 service call processor.

Software built for OS-9 uses SWI2 to call OS-9. If you reset this vector
these programs will not work. If you change all three vectors, you will
not be able to call OS-9 at all.

\subsection{F\$STime - Set system date and time}

{\def\LTcaptype{none} % do not increment counter
\begin{longtable}[]{@{}
  >{\raggedright\arraybackslash}p{(\linewidth - 2\tabcolsep) * \real{0.2800}}
  >{\raggedright\arraybackslash}p{(\linewidth - 2\tabcolsep) * \real{0.7200}}@{}}
\toprule\noalign{}
\endhead
\bottomrule\noalign{}
\endlastfoot
ASSEMBLER CALL: & OS9 F\$STime \\
MACHINE CODE: & 103F 16 \\
INPUT: & (X) = Address of time packet (see below) \\
OUTPUT: & Time/date is set. \\
ERROR OUTPUT: & \begin{minipage}[t]{\linewidth}\raggedright
(CC)~=~C~bit~set.\\
(B)~=~Appropriate~error~code.
\end{minipage} \\
\end{longtable}
}

This service request is used to set the current system date/time and
start the system real-time clock. The date and time are passed in a time
packet as follows:

{\def\LTcaptype{none} % do not increment counter
\begin{longtable}[]{@{}
  >{\raggedright\arraybackslash}p{(\linewidth - 2\tabcolsep) * \real{0.4000}}
  >{\raggedright\arraybackslash}p{(\linewidth - 2\tabcolsep) * \real{0.6000}}@{}}
\toprule\noalign{}
\begin{minipage}[b]{\linewidth}\raggedright
OFFSET
\end{minipage} & \begin{minipage}[b]{\linewidth}\raggedright
VALUE
\end{minipage} \\
\midrule\noalign{}
\endhead
\bottomrule\noalign{}
\endlastfoot
0 & year \\
1 & month \\
2 & day \\
3 & hours \\
4 & minutes \\
5 & seconds \\
\end{longtable}
}

\subsection{F\$Time - Get system date and time}

{\def\LTcaptype{none} % do not increment counter
\begin{longtable}[]{@{}
  >{\raggedright\arraybackslash}p{(\linewidth - 2\tabcolsep) * \real{0.2800}}
  >{\raggedright\arraybackslash}p{(\linewidth - 2\tabcolsep) * \real{0.7200}}@{}}
\toprule\noalign{}
\endhead
\bottomrule\noalign{}
\endlastfoot
ASSEMBLER CALL: & OS9 F\$Time \\
MACHINE CODE: & 103F 15 \\
INPUT: & (X) = Address of place to store the time packet. \\
OUTPUT: & Time packet (see below). \\
ERROR OUTPUT: & \begin{minipage}[t]{\linewidth}\raggedright
(CC)~=~C~bit~set.\\
(B)~=~Appropriate~error~code.
\end{minipage} \\
\end{longtable}
}

This returns the current system date and time in the form of a six byte
packet (in binary). The packet is copied to the address passed in X. The
packet looks like:

{\def\LTcaptype{none} % do not increment counter
\begin{longtable}[]{@{}
  >{\raggedright\arraybackslash}p{(\linewidth - 2\tabcolsep) * \real{0.4000}}
  >{\raggedright\arraybackslash}p{(\linewidth - 2\tabcolsep) * \real{0.6000}}@{}}
\toprule\noalign{}
\begin{minipage}[b]{\linewidth}\raggedright
OFFSET
\end{minipage} & \begin{minipage}[b]{\linewidth}\raggedright
VALUE
\end{minipage} \\
\midrule\noalign{}
\endhead
\bottomrule\noalign{}
\endlastfoot
0 & year \\
1 & month \\
2 & day \\
3 & hours \\
4 & minutes \\
5 & seconds \\
\end{longtable}
}

\subsection{F\$Unlink - Unlink a module}

{\def\LTcaptype{none} % do not increment counter
\begin{longtable}[]{@{}
  >{\raggedright\arraybackslash}p{(\linewidth - 2\tabcolsep) * \real{0.2800}}
  >{\raggedright\arraybackslash}p{(\linewidth - 2\tabcolsep) * \real{0.7200}}@{}}
\toprule\noalign{}
\endhead
\bottomrule\noalign{}
\endlastfoot
ASSEMBLER CALL: & OS9 F\$Unlink \\
MACHINE CODE: & 103F 02 \\
INPUT: & (U) = Address of the module header. \\
OUTPUT: & None \\
ERROR OUTPUT: & \begin{minipage}[t]{\linewidth}\raggedright
(CC)~=~C~bit~set.\\
(B)~=~Appropriate~error~code.
\end{minipage} \\
\end{longtable}
}

Tells OS-9 that the module is no longer needed by the calling process.
The module\textquotesingle s link count is decremented, and the module
is destroyed and its memory deallocated when the link count equals zero.
The module will not be destroyed if in use by any other process(es)
because its link count will be non-zero. In Level Two systems, the
module is usually switched out of the process\textquotesingle{} address
space.

Device driver modules in use or certain system modules cannot he
unlinked. ROMed modules can be unlinked but cannot be deleted from the
module directory.

\subsection{F\$Wait - Wait for child process to die}\label{f.wait}

{\def\LTcaptype{none} % do not increment counter
\begin{longtable}[]{@{}
  >{\raggedright\arraybackslash}p{(\linewidth - 2\tabcolsep) * \real{0.2800}}
  >{\raggedright\arraybackslash}p{(\linewidth - 2\tabcolsep) * \real{0.7200}}@{}}
\toprule\noalign{}
\endhead
\bottomrule\noalign{}
\endlastfoot
ASSEMBLER CALL: & OS9 F\$Wait \\
MACHINE CODE: & 103F 04 \\
INPUT: & None \\
OUTPUT: & \begin{minipage}[t]{\linewidth}\raggedright
(A)~=~Deceased~child~process\textquotesingle~process~ID\\
(B)~=~Child~process\textquotesingle~exit~status~code
\end{minipage} \\
ERROR OUTPUT: & \begin{minipage}[t]{\linewidth}\raggedright
(CC)~=~C~bit~set.\\
(B)~=~Appropriate~error~code.
\end{minipage} \\
\end{longtable}
}

The calling process is deactivated until a child process terminates by
executing an \hyperref[f.exit]{F\$Exit} system call, or by receiving a
signal. The child\textquotesingle s ID number and exit status is
returned to the parent. If the child died due to a signal, the exit
status byte (B register) is the signal code.

If the caller has several children, the caller is activated when the
first one dies, so one \hyperref[f.wait]{F\$Wait} system call is
required to detect termination of each child.

If a child died before the \hyperref[f.wait]{F\$Wait} call, the caller
is reactivated almost immediately. \hyperref[f.wait]{F\$Wait} will
return an error if the caller has no children.

See the \hyperref[f.exit]{F\$Exit} description for more related
information.

\subsection{F\$All64 - Allocate a 64 byte memory block}\label{f.all64}

{\def\LTcaptype{none} % do not increment counter
\begin{longtable}[]{@{}
  >{\raggedright\arraybackslash}p{(\linewidth - 2\tabcolsep) * \real{0.2800}}
  >{\raggedright\arraybackslash}p{(\linewidth - 2\tabcolsep) * \real{0.7200}}@{}}
\toprule\noalign{}
\endhead
\bottomrule\noalign{}
\endlastfoot
ASSEMBLER CALL: & OS9 F\$All64 \\
MACHINE CODE: & 103F 30 \\
INPUT: & (X) = Base address of page table (zero if the page table has
not yet been allocated). \\
OUTPUT: & \begin{minipage}[t]{\linewidth}\raggedright
(A)~=~Block~number\\
(X)~=~Base~address~of~page~table\\
(Y)~=~Address~of~block.
\end{minipage} \\
ERROR OUTPUT: & \begin{minipage}[t]{\linewidth}\raggedright
(CC)~=~C~bit~set.\\
(B)~=~Appropriate~error~code.
\end{minipage} \\
\end{longtable}
}

This system mode service request is used to dynamically allocate 64 byte
blocks of memory by splitting whole pages (256 byte) into tour sections.
The first 64 bytes of the base page are used as a ``page table'', which
contains the MSB of all pages in the memory structure. Passing a value
of zero in the X register will cause the \hyperref[f.all64]{F\$All64}
service request to allocate a new base page and the first 64 byte memory
block. Whenever a new page is needed, an \hyperref[f.srqmem]{F\$SRqMem}
service request will automatically be executed. The first byte of each
block contains the block number; routines using this service request
should not alter it. Below is a diagram to show how 7 blocks might be
allocated:

\begin{figure}
\centering
\begin{verbatim}
                   ANY 256 BYTE            ANY 256 BYTE
                   MEMORY PAGE             MEMORY PAGE
BASE PAGE  --->   +-------------+         +-------------+
                  !             !         !X            !
                  !  PAGE TABLE !         !   BLOCK 4   !
                  !  (64 bytes) !         !  (64 bytes) !
                  +-------------+         +-------------+
                  !X            !         !X            !
                  !   BLOCK 1   !         !   BLOCK 5   !
                  !  (64 bytes) !         !  (64 bytes) !
                  +-------------+         +-------------+
                  !X            !         !X            !
                  !   BLOCK 2   !         !   BLOCK 6   !
                  !  (64 bytes) !         !  (64 bytes) !
                  +-------------+         +-------------+
                  !X            !         !X            !
                  !   BLOCK 3   !         !   BLOCK 7   !
                  !  (64 bytes) !         !  (64 bytes) !
                  +-------------+         +-------------+
\end{verbatim}
\caption{}
\end{figure}

This is a privileged system mode service request.

\subsection{F\$AProc - Insert process in active process queue}

{\def\LTcaptype{none} % do not increment counter
\begin{longtable}[]{@{}
  >{\raggedright\arraybackslash}p{(\linewidth - 2\tabcolsep) * \real{0.2800}}
  >{\raggedright\arraybackslash}p{(\linewidth - 2\tabcolsep) * \real{0.7200}}@{}}
\toprule\noalign{}
\endhead
\bottomrule\noalign{}
\endlastfoot
ASSEMBLER CALL: & OS9 F\$AProc \\
MACHINE CODE: & 103F 2C \\
INPUT: & (X) = Address of process descriptor. \\
OUTPUT: & None. \\
ERROR OUTPUT: & \begin{minipage}[t]{\linewidth}\raggedright
(CC)~=~C~bit~set.\\
(B)~=~Appropriate~error~code.
\end{minipage} \\
\end{longtable}
}

This system mode service request inserts a process into the active
process queue so that it may be scheduled for execution.

All processes already in the active process queue are aged, and the age
of the specified process is set to its priority. If the process is in
system state, it is inserted after any other process\textquotesingle s
also in system state, but before any process in user state. If the
process is in user state, it is inserted according to its age.

This is a privileged system mode service request.

\subsection{F\$Find64 - Find a 64 byte memory block}

{\def\LTcaptype{none} % do not increment counter
\begin{longtable}[]{@{}
  >{\raggedright\arraybackslash}p{(\linewidth - 2\tabcolsep) * \real{0.2800}}
  >{\raggedright\arraybackslash}p{(\linewidth - 2\tabcolsep) * \real{0.7200}}@{}}
\toprule\noalign{}
\endhead
\bottomrule\noalign{}
\endlastfoot
ASSEMBLER CALL: & OS9 F\$Find64 \\
MACHINE CODE: & 103F 2F \\
INPUT: & \begin{minipage}[t]{\linewidth}\raggedright
(X)~=~Address~of~base~page.\\
(A)~=~Block~number.
\end{minipage} \\
OUTPUT: & (Y) = Address of block. \\
ERROR OUTPUT: & \begin{minipage}[t]{\linewidth}\raggedright
(CC)~=~C~bit~set.\\
(B)~=~Appropriate~error~code.
\end{minipage} \\
\end{longtable}
}

This system mode service request will return the address of a 64 byte
memory block as described in the \hyperref[f.all64]{F\$All64} service
request. OS-9 used this service request to find process descriptors and
path descriptors when given their number.

Block numbers range from 1..N

This is a privileged system mode service request.

\subsection{F\$IODel - Delete I/O device from system}

{\def\LTcaptype{none} % do not increment counter
\begin{longtable}[]{@{}
  >{\raggedright\arraybackslash}p{(\linewidth - 2\tabcolsep) * \real{0.2800}}
  >{\raggedright\arraybackslash}p{(\linewidth - 2\tabcolsep) * \real{0.7200}}@{}}
\toprule\noalign{}
\endhead
\bottomrule\noalign{}
\endlastfoot
ASSEMBLER CALL: & OS9 F\$IODel \\
MACHINE CODE: & 103F 33 \\
INPUT: & (X) = Address of an I/O module, (see description). \\
OUTPUT: & None. \\
ERROR OUTPUT: & \begin{minipage}[t]{\linewidth}\raggedright
(CC)~=~C~bit~set.\\
(B)~=~Appropriate~error~code.
\end{minipage} \\
\end{longtable}
}

This system mode service request is used to determine whether or not an
I/O module is being used. The X register passes the address of a device
descriptor module, device driver module, or file manager module. The
address is used to search the device table, and if found the use count
is checked to see if it is zero. If it is not zero, an error condition
is returned.

This service request is used primarily by IOMAN and may be of limited or
no use for other applications.

This is a privileged system mode service request.

\subsection{F\$IOQU - Enter I/O queue}

{\def\LTcaptype{none} % do not increment counter
\begin{longtable}[]{@{}
  >{\raggedright\arraybackslash}p{(\linewidth - 2\tabcolsep) * \real{0.2800}}
  >{\raggedright\arraybackslash}p{(\linewidth - 2\tabcolsep) * \real{0.7200}}@{}}
\toprule\noalign{}
\endhead
\bottomrule\noalign{}
\endlastfoot
ASSEMBLER CALL: & OS9 F\$IOQU \\
MACHINE CODE: & 103F 2B \\
INPUT: & (A) = Process Number. \\
OUTPUT: & None. \\
ERROR OUTPUT: & \begin{minipage}[t]{\linewidth}\raggedright
(CC)~=~C~bit~set.\\
(B)~=~Appropriate~error~code.
\end{minipage} \\
\end{longtable}
}

This system mode service request links the calling process into the I/O
queue of the specified process and performs an untimed sleep. It is
assumed that routines associated with the specified process will send a
wakeup signal to the calling process.

This is a privileged system mode service request.

\subsection{F\$IRQ - Add or remove device from IRQ table}\label{f.irq}

{\def\LTcaptype{none} % do not increment counter
\begin{longtable}[]{@{}
  >{\raggedright\arraybackslash}p{(\linewidth - 2\tabcolsep) * \real{0.2800}}
  >{\raggedright\arraybackslash}p{(\linewidth - 2\tabcolsep) * \real{0.7200}}@{}}
\toprule\noalign{}
\endhead
\bottomrule\noalign{}
\endlastfoot
ASSEMBLER CALL: & OS9 F\$IRQ \\
MACHINE CODE: & 103F 2A \\
INPUT: & \begin{minipage}[t]{\linewidth}\raggedright
(X)~=~Zero~to~remove~device~from~table,~or~the~address~of~a\\
packet~as~defined~below~to~add~a~device~to~the~IRQ~polling~table:\\
\strut ~\\
\strut ~~~~~~~~{[}x{]}~=~flip~byte\\
\strut ~~~~~~~~{[}X+1{]}~=~mask~byte\\
\strut ~~~~~~~~{[}X+2{]}~=~priority\\
(U)~=~Address~of~service~routine\textquotesingle s~static~storage~area.\\
(Y)~=~Device~IRQ~service~routine~address.\\
(D)~=~Address~of~the~device~status~register.
\end{minipage} \\
OUTPUT: & None. \\
ERROR OUTPUT: & \begin{minipage}[t]{\linewidth}\raggedright
(CC)~=~C~bit~set.\\
(B)~=~Appropriate~error~code.
\end{minipage} \\
\end{longtable}
}

This service request is used to add a device to or remove a device from
the IRQ polling table. To remove a device from the table the input
should be (X)=0, (U)= Addr of service routine\textquotesingle s static
storage. This service request is primarily used by device driver
routines. See the text of this manual for a complete discussion of the
interrupt polling system.

PACKET DEFINITIONS:

{\def\LTcaptype{none} % do not increment counter
\begin{longtable}[]{@{}
  >{\raggedright\arraybackslash}p{(\linewidth - 2\tabcolsep) * \real{0.3000}}
  >{\raggedright\arraybackslash}p{(\linewidth - 2\tabcolsep) * \real{0.7000}}@{}}
\toprule\noalign{}
\endhead
\bottomrule\noalign{}
\endlastfoot
Flip Byte & This byte selects whether the bits in the device status
register are active when set or active when cleared. A set bit(s)
identifies the active bit(s). \\
Mask Byte & This byte selects one or more bits within the device status
register that are interrupt request flag(s). A set bit identifies an
active bit(s) \\
Priority & \begin{minipage}[t]{\linewidth}\raggedright
The~device~priority~number:\\
\strut ~~~~~~~~~~0~=~lowest\\
\strut ~~~~~~~~255~=~highest
\end{minipage} \\
\end{longtable}
}

This is a privileged system mode service request.

\subsection{F\$NProc - Start next process}

{\def\LTcaptype{none} % do not increment counter
\begin{longtable}[]{@{}
  >{\raggedright\arraybackslash}p{(\linewidth - 2\tabcolsep) * \real{0.2800}}
  >{\raggedright\arraybackslash}p{(\linewidth - 2\tabcolsep) * \real{0.7200}}@{}}
\toprule\noalign{}
\endhead
\bottomrule\noalign{}
\endlastfoot
ASSEMBLER CALL: & OS9 F\$NProc \\
MACHINE CODE: & 103F 2D \\
INPUT: & None. \\
OUTPUT: & Control does not return to caller. \\
\end{longtable}
}

This system mode service request takes the next process out of the
Active Process Queue and initiates its execution. If there is no process
in the queue, OS-9 waits for an interrupt, and then checks the active
process queue again.

This is a privileged system mode service request.

\subsection{F\$Ret64 - Deallocate a 64 byte memory block}

{\def\LTcaptype{none} % do not increment counter
\begin{longtable}[]{@{}
  >{\raggedright\arraybackslash}p{(\linewidth - 2\tabcolsep) * \real{0.2800}}
  >{\raggedright\arraybackslash}p{(\linewidth - 2\tabcolsep) * \real{0.7200}}@{}}
\toprule\noalign{}
\endhead
\bottomrule\noalign{}
\endlastfoot
ASSEMBLER CALL: & OS9 F\$Ret64 \\
MACHINE CODE: & 103F 31 \\
INPUT: & \begin{minipage}[t]{\linewidth}\raggedright
(X)~=~Address~of~the~base~page.\\
(A)~=~Block~number.
\end{minipage} \\
OUTPUT: & None. \\
ERROR OUTPUT: & \begin{minipage}[t]{\linewidth}\raggedright
(CC)~=~C~bit~set.\\
(B)~=~Appropriate~error~code.
\end{minipage} \\
\end{longtable}
}

This system mode service request deallocates a 64 byte block of memory
as described in the \hyperref[f.all64]{F\$All64} service request.

This is a privileged system mode service request.

\subsection{F\$SRqMem - System memory request}\label{f.srqmem}

{\def\LTcaptype{none} % do not increment counter
\begin{longtable}[]{@{}
  >{\raggedright\arraybackslash}p{(\linewidth - 2\tabcolsep) * \real{0.2800}}
  >{\raggedright\arraybackslash}p{(\linewidth - 2\tabcolsep) * \real{0.7200}}@{}}
\toprule\noalign{}
\endhead
\bottomrule\noalign{}
\endlastfoot
ASSEMBLER CALL: & OS9 F\$SRqMem \\
MACHINE CODE: & 103F 28 \\
INPUT: & (D) = Byte count. \\
OUTPUT: & (U) = Beginning address of memory area. \\
ERROR OUTPUT: & \begin{minipage}[t]{\linewidth}\raggedright
(CC)~=~C~bit~set.\\
(B)~=~Appropriate~error~code.
\end{minipage} \\
\end{longtable}
}

This system mode service request allocates a block of memory from the
top of available RAM of the specified size. The size requested is
rounded to the next 256 byte page boundary.

This is a privileged system mode service request.

\subsection{F\$SRtMem - Return System Memory}

{\def\LTcaptype{none} % do not increment counter
\begin{longtable}[]{@{}
  >{\raggedright\arraybackslash}p{(\linewidth - 2\tabcolsep) * \real{0.2800}}
  >{\raggedright\arraybackslash}p{(\linewidth - 2\tabcolsep) * \real{0.7200}}@{}}
\toprule\noalign{}
\endhead
\bottomrule\noalign{}
\endlastfoot
ASSEMBLER CALL: & OS9 F\$SRtMem \\
MACHINE CODE: & 103F 29 \\
INPUT: & \begin{minipage}[t]{\linewidth}\raggedright
(U)~=~Beginning~address~of~memory~to~return.\\
(D)~=~Number~of~bytes~to~return.
\end{minipage} \\
OUTPUT: & None. \\
ERROR OUTPUT: & \begin{minipage}[t]{\linewidth}\raggedright
(CC)~=~C~bit~set.\\
(B)~=~Appropriate~error~code.
\end{minipage} \\
\end{longtable}
}

This system mode service request is used to deallocate a block of
contigous 256 byte pages. The U register must point to an even page
boundary.

This is a privileged system mode service request.

\subsection{F\$VModul - Verify module}

{\def\LTcaptype{none} % do not increment counter
\begin{longtable}[]{@{}
  >{\raggedright\arraybackslash}p{(\linewidth - 2\tabcolsep) * \real{0.2800}}
  >{\raggedright\arraybackslash}p{(\linewidth - 2\tabcolsep) * \real{0.7200}}@{}}
\toprule\noalign{}
\endhead
\bottomrule\noalign{}
\endlastfoot
ASSEMBLER CALL: & OS9 F\$VModul \\
MACHINE CODE: & 103F 2E \\
INPUT: & (X) = Address of module to verify \\
OUTPUT: & (U) = Address of module directory entry \\
ERROR OUTPUT: & \begin{minipage}[t]{\linewidth}\raggedright
(CC)~=~C~bit~set.\\
(B)~=~Appropriate~error~code.
\end{minipage} \\
\end{longtable}
}

This system mode service request checks the module header parity and CRC
bytes of an OS-9 module. If these values are valid, then the module
directory is searched for a module with the same name. If a module with
the same name exists, the one with the highest revision level is
retained in the module directory. Ties are broken in favor of the
established module.

This is a privileged system mode service request.

\subsection{I\$Attach - Attach a new device to the
system}\label{i.attach}

{\def\LTcaptype{none} % do not increment counter
\begin{longtable}[]{@{}
  >{\raggedright\arraybackslash}p{(\linewidth - 2\tabcolsep) * \real{0.2800}}
  >{\raggedright\arraybackslash}p{(\linewidth - 2\tabcolsep) * \real{0.7200}}@{}}
\toprule\noalign{}
\endhead
\bottomrule\noalign{}
\endlastfoot
ASSEMBLER CALL: & OS9 I\$Attach \\
MACHINE CODE: & 103F 80 \\
INPUT: & \begin{minipage}[t]{\linewidth}\raggedright
(X)~=~Address~of~device~name~string\\
(A)~=~Access~mode.
\end{minipage} \\
OUTPUT: & (U) = Address of device table entry \\
ERROR OUTPUT: & \begin{minipage}[t]{\linewidth}\raggedright
(CC)~=~C~bit~set.\\
(B)~=~Appropriate~error~code.
\end{minipage} \\
\end{longtable}
}

This service request is used to attach a new device to the system, or
verify that it is already attached. The device\textquotesingle s name
string is used to search the system module directory to see if a device
descriptor module with the same name is in memory (this is the name the
device will be known by). The descriptor module will contain the name of
the device\textquotesingle s file manager, device driver and other
related information. It it is found and the device is not already
attached, OS-9 will link to its file manager and device driver, and then
place their address\textquotesingle{} in a new device table entry. Any
permanent storage needed by the device driver is allocated, and the
driver\textquotesingle s initialization routine is called (which usually
initializes the hardware).

If the device has already been attached, it will not be reinitialized.

An I\$Attach system call is not required to perform routine I/O. It does
\emph{not} ``reserve'' the device in question - it just prepares it for
subsequent use by any process. Most devices are automatically installed,
so it is used mostly when devices are dynamically installed or to verify
the existence of a device.

The access mode parameter specifies which subsequent read and/or write
operations will be permitted as follows:

0 = Use device capabilities.

1 = Read only.

2 = Write only.

3 = Both read and write.

\subsection{I\$ChgDir - Change working directory}

{\def\LTcaptype{none} % do not increment counter
\begin{longtable}[]{@{}
  >{\raggedright\arraybackslash}p{(\linewidth - 2\tabcolsep) * \real{0.2800}}
  >{\raggedright\arraybackslash}p{(\linewidth - 2\tabcolsep) * \real{0.7200}}@{}}
\toprule\noalign{}
\endhead
\bottomrule\noalign{}
\endlastfoot
ASSEMBLER CALL: & OS9 I\$ChgDir \\
MACHINE CODE: & 103F 86 \\
INPUT: & \begin{minipage}[t]{\linewidth}\raggedright
(X)~=~Address~of~the~pathlist.\\
(A)~=~Access~mode.
\end{minipage} \\
OUTPUT: & None. \\
ERROR OUTPUT: & \begin{minipage}[t]{\linewidth}\raggedright
(CC)~=~C~bit~set.\\
(B)~=~Appropriate~error~code.
\end{minipage} \\
\end{longtable}
}

Changes a process\textquotesingle{} working directory to another
directory file specified by the pathlist. Depending on the access mode
given, the current execution or the current data directory may be
changed (but only one may be changed per call). The file specified must
be a directory file, and the caller must have read permission for it
(public read if not owned by the calling process).

ACCESS MODES

1 = Read

2 = Write

3 = Update (read or write)

4 = Execute

If the access mode is read, write, or update the current data directory
is changed. If the access mode is execute, the current execution
directory is changed.

\subsection{I\$Close - Close a path to a file/device}\label{i.close}

{\def\LTcaptype{none} % do not increment counter
\begin{longtable}[]{@{}
  >{\raggedright\arraybackslash}p{(\linewidth - 2\tabcolsep) * \real{0.2800}}
  >{\raggedright\arraybackslash}p{(\linewidth - 2\tabcolsep) * \real{0.7200}}@{}}
\toprule\noalign{}
\endhead
\bottomrule\noalign{}
\endlastfoot
ASSEMBLER CALL: & OS9 I\$Close \\
MACHINE CODE: & 103F 8F \\
INPUT: & (A) = Path number. \\
OUTPUT: & None. \\
ERROR OUTPUT: & \begin{minipage}[t]{\linewidth}\raggedright
(CC)~=~C~bit~set.\\
(B)~=~Appropriate~error~code.
\end{minipage} \\
\end{longtable}
}

Terminates the I/O path specified by the path number. I/O can no longer
be performed to the file/device, unless another
\hyperref[i.open]{I\$Open} or \hyperref[i.create]{I\$Create} call is
used. Devices that are non-sharable become available to other requesting
processes. All OS-9 internally managed buffers and descriptors are
deallocated.

Note: Because the OS9 F\$Exit service request automatically closes all
open paths (except the standard I/O paths), it may not he necessary to
close them individually with the OS9 \hyperref[i.close]{I\$Close}
service request.

Standard I/O paths are not typically closed except when it is desired to
change the files/devices they correspond to.

\subsection{I\$Create - Create a path to a new file}\label{i.create}

{\def\LTcaptype{none} % do not increment counter
\begin{longtable}[]{@{}
  >{\raggedright\arraybackslash}p{(\linewidth - 2\tabcolsep) * \real{0.2800}}
  >{\raggedright\arraybackslash}p{(\linewidth - 2\tabcolsep) * \real{0.7200}}@{}}
\toprule\noalign{}
\endhead
\bottomrule\noalign{}
\endlastfoot
ASSEMBLER CALL: & OS9 I\$Create \\
MACHINE CODE: & 103F 83 \\
INPUT: & \begin{minipage}[t]{\linewidth}\raggedright
(X)~=~Address~of~the~pathlist.\\
(A)~=~Access~mode.\\
(B)~=~File~attributes.
\end{minipage} \\
OUTPUT: & \begin{minipage}[t]{\linewidth}\raggedright
(X)~=~Updated~past~the~pathlist~(trailing~blanks~skipped)\\
(A)~=~Path~number.
\end{minipage} \\
ERROR OUTPUT: & \begin{minipage}[t]{\linewidth}\raggedright
(CC)~=~C~bit~set.\\
(B)~=~Appropriate~error~code.
\end{minipage} \\
\end{longtable}
}

Used to create a new file on a multifile mass storage device. The
pathlist is parsed, and the new file name is entered in the specified
(or default working) directory. The file is given the attributes passed
in the B register, which has individual bits defined as follows:

bit~0~=~read~permit\\
bit~1~=~write~permit\\
bit~2~=~execute~permit\\
bit~3~=~public~read~permit\\
bit~4~=~public~write~permit\\
bit~5~-~public~execute~permit\\
bit~6~=~sharable~file\\

The access mode parameter passed in register A must be either ``WRITE''
or ``UPDATE''. This only affects the file until it is closed; it can be
reopened later in any access mode allowed by the file attributes (see
\hyperref[i.open]{I\$Open}). Files open for ``WRITE'' may allow faster
data transfer than ``UPDATE'', which sometimes needs to preread setors.
These access codes are defined as given below:

2~=~Write~only\\
3~=~Update~(read~and~write)\\

NOTE: If the execute bit (bit 2) is set, the file will be created in the
working execution directory instead of the working data directory.

The path number returned by OS-9 is used to identify the file in
subsequent I/O service requests until the file is closed.

No data storage is initially allocated for the file at the time it is
created; this is done automatically by \hyperref[i.write]{I\$Write} or
explicitly by the \hyperref[i.setstt]{I\$SetStt} call.

An error will occur if the file name already exists in the directory.
\hyperref[i.create]{I\$Create} calls that specify non-multiple file
devices (such as printers, terminals, etc.) work correctly: the
\hyperref[i.create]{I\$Create} behaves the same as
\hyperref[i.open]{I\$Open}. Create cannot be used to make directory
files (see \hyperref[i.makdir]{I\$MakDir}).

\subsection{I\$Delete - Delete a file}

{\def\LTcaptype{none} % do not increment counter
\begin{longtable}[]{@{}
  >{\raggedright\arraybackslash}p{(\linewidth - 2\tabcolsep) * \real{0.2800}}
  >{\raggedright\arraybackslash}p{(\linewidth - 2\tabcolsep) * \real{0.7200}}@{}}
\toprule\noalign{}
\endhead
\bottomrule\noalign{}
\endlastfoot
ASSEMBLER CALL: & OS9 I\$Delete \\
MACHINE CODE: & 103F 87 \\
INPUT: & (X) = Address of pathlist. \\
OUTPUT: & (X) = Updated past pathlist (trailing spaces skipped). \\
ERROR OUTPUT: & \begin{minipage}[t]{\linewidth}\raggedright
(CC)~=~C~bit~set.\\
(B)~=~Appropriate~error~code.
\end{minipage} \\
\end{longtable}
}

This service request deletes the file specified by the pathlist. The
file must have write permission attributes (public write if not the
owner), and reside on a multifile mass storage device. Attempts to
delete devices will result in an error.

\subsection{I\$Detach - Remove a device from the system}\label{i.detach}

{\def\LTcaptype{none} % do not increment counter
\begin{longtable}[]{@{}
  >{\raggedright\arraybackslash}p{(\linewidth - 2\tabcolsep) * \real{0.2800}}
  >{\raggedright\arraybackslash}p{(\linewidth - 2\tabcolsep) * \real{0.7200}}@{}}
\toprule\noalign{}
\endhead
\bottomrule\noalign{}
\endlastfoot
ASSEMBLER CALL: & OS9 I\$Detach \\
MACHINE CODE: & 103F 81 \\
INPUT: & (U) = Address of the device table entry. \\
OUTPUT: & None. \\
ERROR OUTPUT: & \begin{minipage}[t]{\linewidth}\raggedright
(CC)~=~C~bit~set.\\
(B)~=~Appropriate~error~code.
\end{minipage} \\
\end{longtable}
}

Removes a device from the system device table if not in use by any other
process. The device driver\textquotesingle s termination routine is
called, then any permanent storage assigned to the driver is
deallocated. The device driver and file manager modules associated with
the device are unlinked (and may be destroyed if not in use by another
process.

The \hyperref[i.detach]{I\$Detach} service request must be used to
un-attach devices that were attached with the
\hyperref[i.attach]{I\$Attach} service request. Both of these are used
mainly by IOMAN and are of limited (or no use) to the typical user.
SCFMAN also uses
\hyperref[i.attach]{I\$Attach}/\hyperref[i.detach]{I\$Detach} to setup
its second (echo) device.

\subsection{I\$Dup Duplicate a path}

{\def\LTcaptype{none} % do not increment counter
\begin{longtable}[]{@{}
  >{\raggedright\arraybackslash}p{(\linewidth - 2\tabcolsep) * \real{0.2800}}
  >{\raggedright\arraybackslash}p{(\linewidth - 2\tabcolsep) * \real{0.7200}}@{}}
\toprule\noalign{}
\endhead
\bottomrule\noalign{}
\endlastfoot
ASSEMBLER CALL: & OS9 I\$Dup \\
MACHINE CODE: & 103F 82 \\
INPUT: & (A) = Path number of path to duplicate. \\
OUTPUT: & (A) = New path number. \\
ERROR OUTPUT: & \begin{minipage}[t]{\linewidth}\raggedright
(CC)~=~C~bit~set.\\
(B)~=~Appropriate~error~code.
\end{minipage} \\
\end{longtable}
}

Given the number of an existing path, returns another synonymous path
number for the same file or device. Shell uses this service request when
it redirects I/O. Service requests using either the old or new path
numbers operate on the same file or device.

NOTE: This only increments the ``use count'' of a path descriptor and
returns the synonymous path number. The path descriptor is not copied.

\subsection{I\$GetStt - Get file device status}\label{i.getstt}

{\def\LTcaptype{none} % do not increment counter
\begin{longtable}[]{@{}
  >{\raggedright\arraybackslash}p{(\linewidth - 2\tabcolsep) * \real{0.2800}}
  >{\raggedright\arraybackslash}p{(\linewidth - 2\tabcolsep) * \real{0.7200}}@{}}
\toprule\noalign{}
\endhead
\bottomrule\noalign{}
\endlastfoot
ASSEMBLER CALL: & OS9 I\$GetStt \\
MACHINE CODE: & 103F 8D \\
INPUT: & \begin{minipage}[t]{\linewidth}\raggedright
(A)~=~Path~number.\\
(B)~Status~code.\\
(Other~registers~depend~upon~status~code)
\end{minipage} \\
OUTPUT: & (depends upon status code) \\
ERROR OUTPUT: & \begin{minipage}[t]{\linewidth}\raggedright
(CC)~=~C~bit~set.\\
(B)~=~Appropriate~error~code.
\end{minipage} \\
\end{longtable}
}

This system is a ``wild card'' call used to handle individual device
parameters that:

\begin{enumerate}
\def\labelenumi{\alph{enumi}.}
\item
  are not uniform on all devices
\item
  are highly hardware dependent
\item
  need to be user-changable
\end{enumerate}

The exact operation of this call depends on the device driver an file
manager associated with the path. A typical use is to determine a
terminal\textquotesingle s parameters for backspace character, delete
character, echo on/off, null padding, paging, etc. It is commonly used
in conjunction with the \hyperref[i.setstt]{I\$SetStt} service request
which is used to set the device operating parameters. Below are
presently defined function codes for \hyperref[i.getstt]{I\$GetStt}:

{\def\LTcaptype{none} % do not increment counter
\begin{longtable}[]{@{}
  >{\raggedright\arraybackslash}p{(\linewidth - 4\tabcolsep) * \real{0.2105}}
  >{\raggedright\arraybackslash}p{(\linewidth - 4\tabcolsep) * \real{0.1754}}
  >{\raggedright\arraybackslash}p{(\linewidth - 4\tabcolsep) * \real{0.6140}}@{}}
\toprule\noalign{}
\begin{minipage}[b]{\linewidth}\raggedright
MNEMONIC
\end{minipage} & \begin{minipage}[b]{\linewidth}\raggedright
CODE
\end{minipage} & \begin{minipage}[b]{\linewidth}\raggedright
FUNCTION
\end{minipage} \\
\midrule\noalign{}
\endhead
\bottomrule\noalign{}
\endlastfoot
SS.Opt & 0 & Read the 32 byte option section of the path descriptor. \\
SS.Ready & 1 & Test for data ready on SCFMAN-type device. \\
SS.Size & 2 & Return current file size (on RBFMAN-type devices). \\
SS.Pos & 5 & Get current file position. \\
SS.EOF & 6 & Test for end of file. \\
SS.ScSiz & 38 & Width of screen in characters. \\
\end{longtable}
}

CODES 7-127 Reserved for future use.

CODES 128-255 These getstat codes and their parameter passing
conventions are user definable (see the sections of this manual on
writing device drivers). The function code and register stack are passed
to the device driver.

Parameter Passing Conventions

The parameter passing conventions for each of these function codes are
given below:

{\def\LTcaptype{none} % do not increment counter
\begin{longtable}[]{@{}
  >{\raggedright\arraybackslash}p{(\linewidth - 2\tabcolsep) * \real{0.3000}}
  >{\raggedright\arraybackslash}p{(\linewidth - 2\tabcolsep) * \real{0.7000}}@{}}
\toprule\noalign{}
\endhead
\bottomrule\noalign{}
\endlastfoot
SS.Opt (code 0): & Read option section of the path descriptor. \\
INPUT: & \begin{minipage}[t]{\linewidth}\raggedright
(A)~=~Path~number\\
(B)~=~Function~code~0\\
(X)~=~Address~of~place~to~put~a~32~byte~status~packet.
\end{minipage} \\
OUTPUT: & Status packet. \\
ERROR OUTPUT: & \begin{minipage}[t]{\linewidth}\raggedright
(CC)~=~C~bit~set.\\
(B)~=~Appropriate~error~code.
\end{minipage} \\
\end{longtable}
}

This getstat function reads the option section of the path descriptor
and copies it into the 32 byte area pointed to by the X register. It is
typically used to determine the current settings for echo, auto line
feed, etc. For a complete description of the status packet, please see
the section of this manual on path descriptors.

{\def\LTcaptype{none} % do not increment counter
\begin{longtable}[]{@{}
  >{\raggedright\arraybackslash}p{(\linewidth - 2\tabcolsep) * \real{0.3000}}
  >{\raggedright\arraybackslash}p{(\linewidth - 2\tabcolsep) * \real{0.7000}}@{}}
\toprule\noalign{}
\endhead
\bottomrule\noalign{}
\endlastfoot
SS.Ready (code 1): & Test for data available on SCFMAN supported
devices. \\
INPUT: & \begin{minipage}[t]{\linewidth}\raggedright
(A)~=~Path~number\\
(B)~=~Function~code~1
\end{minipage} \\
OUTPUT: & \\
\end{longtable}
}

{\def\LTcaptype{none} % do not increment counter
\begin{longtable}[]{@{}
  >{\raggedright\arraybackslash}p{(\linewidth - 2\tabcolsep) * \real{0.3000}}
  >{\raggedright\arraybackslash}p{(\linewidth - 2\tabcolsep) * \real{0.7000}}@{}}
\toprule\noalign{}
\endhead
\bottomrule\noalign{}
\endlastfoot
SS.Size (code 2): & Get current file size (RBFMAN supported devices
only) \\
INPUT: & \begin{minipage}[t]{\linewidth}\raggedright
(A)~=~Path~number\\
(B)~=~Function~code~2
\end{minipage} \\
OUTPUT: & \begin{minipage}[t]{\linewidth}\raggedright
(X)~=~M.S.~16~bits~of~current~file~size.\\
(U)~=~L.S.~16~bits~of~current~file~size.
\end{minipage} \\
ERROR OUTPUT: & \begin{minipage}[t]{\linewidth}\raggedright
(CC)~=~C~bit~set.\\
(B)~=~Appropriate~error~code.
\end{minipage} \\
\end{longtable}
}

{\def\LTcaptype{none} % do not increment counter
\begin{longtable}[]{@{}
  >{\raggedright\arraybackslash}p{(\linewidth - 2\tabcolsep) * \real{0.3000}}
  >{\raggedright\arraybackslash}p{(\linewidth - 2\tabcolsep) * \real{0.7000}}@{}}
\toprule\noalign{}
\endhead
\bottomrule\noalign{}
\endlastfoot
SS.Pos (code 5): & Get current file position (RBFMAN supported devices
only) \\
INPUT: & \begin{minipage}[t]{\linewidth}\raggedright
(A)~=~Path~number\\
(B)~=~Function~code~5
\end{minipage} \\
OUTPUT: & \begin{minipage}[t]{\linewidth}\raggedright
(X)~=~M.S.~16~bits~of~current~file~position.\\
(U)~=~L.S.~16~bits~of~current~file~position.
\end{minipage} \\
ERROR OUTPUT: & \begin{minipage}[t]{\linewidth}\raggedright
(CC)~=~C~bit~set.\\
(B)~=~Appropriate~error~code.
\end{minipage} \\
\end{longtable}
}

{\def\LTcaptype{none} % do not increment counter
\begin{longtable}[]{@{}
  >{\raggedright\arraybackslash}p{(\linewidth - 2\tabcolsep) * \real{0.3000}}
  >{\raggedright\arraybackslash}p{(\linewidth - 2\tabcolsep) * \real{0.7000}}@{}}
\toprule\noalign{}
\endhead
\bottomrule\noalign{}
\endlastfoot
SS.EOF (code 6): & Test for end of file. \\
INPUT: & \begin{minipage}[t]{\linewidth}\raggedright
(A)~=~Path~number\\
(B)~=~Function~code~6
\end{minipage} \\
OUTPUT: & \\
\end{longtable}
}

{\def\LTcaptype{none} % do not increment counter
\begin{longtable}[]{@{}
  >{\raggedright\arraybackslash}p{(\linewidth - 2\tabcolsep) * \real{0.3000}}
  >{\raggedright\arraybackslash}p{(\linewidth - 2\tabcolsep) * \real{0.7000}}@{}}
\toprule\noalign{}
\endhead
\bottomrule\noalign{}
\endlastfoot
SS.ScSiz (code 38): & Return screen size for COCO (SCFMAN supported
devices only) \\
INPUT: & \begin{minipage}[t]{\linewidth}\raggedright
(A)~=~Path~number\\
(B)~=~Function~code~38
\end{minipage} \\
OUTPUT: & (X) = Width of screen in characters. Typically 32 or 80. \\
ERROR OUTPUT: & \begin{minipage}[t]{\linewidth}\raggedright
(CC)~=~C~bit~set.\\
(B)~=~Appropriate~error~code.
\end{minipage} \\
\end{longtable}
}

\subsection{I\$MakDir - Make a new directory}\label{i.makdir}

{\def\LTcaptype{none} % do not increment counter
\begin{longtable}[]{@{}
  >{\raggedright\arraybackslash}p{(\linewidth - 2\tabcolsep) * \real{0.2800}}
  >{\raggedright\arraybackslash}p{(\linewidth - 2\tabcolsep) * \real{0.7200}}@{}}
\toprule\noalign{}
\endhead
\bottomrule\noalign{}
\endlastfoot
ASSEMBLER CALL: & OS9 I\$MakDir \\
MACHINE CODE: & 103F 85 \\
INPUT: & \begin{minipage}[t]{\linewidth}\raggedright
(X)~=~Address~of~pathlist.\\
(B)~=~Directory~attributes.
\end{minipage} \\
OUTPUT: & (X) = Updated path pathlist (trailing spaces skipped). \\
ERROR OUTPUT: & \begin{minipage}[t]{\linewidth}\raggedright
(CC)~=~C~bit~set.\\
(B)~=~Appropriate~error~code.
\end{minipage} \\
\end{longtable}
}

\hyperref[i.makdir]{I\$MakDir} is the only way a new directory file can
be created. It will create and initialize a new directory as specified
by the pathlist. The new directory file contains no entries, except for
an entry for itself (".") and its parent directory ("..")

The caller is made the owner of the directory.
\hyperref[i.makdir]{I\$MakDir} does not return a path number because
directory files are not ``opened'' by this request (use
\hyperref[i.open]{I\$Open} to do so). The new directory will
automatically have its ``directory'' bit set in the access permission
attributes. The remaining attributes are specified by the byte passed in
the B register, which has individual bits defined as follows:

bit~0~=~read~permit\\
bit~1~=~write~permit\\
bit~2~=~execute~permit\\
bit~3~=~public~read~permit\\
bit~4~=~public~write~permit\\
bit~5~-~public~execute~permit\\
bit~6~=~sharable~file\\
bit~7~=~(don\textquotesingle t~care)\\

\subsection{I\$Open - Open a path to a file or device}\label{i.open}

{\def\LTcaptype{none} % do not increment counter
\begin{longtable}[]{@{}
  >{\raggedright\arraybackslash}p{(\linewidth - 2\tabcolsep) * \real{0.2800}}
  >{\raggedright\arraybackslash}p{(\linewidth - 2\tabcolsep) * \real{0.7200}}@{}}
\toprule\noalign{}
\endhead
\bottomrule\noalign{}
\endlastfoot
ASSEMBLER CALL: & OS9 I\$Open \\
MACHINE CODE: & 103F 84 \\
INPUT: & \begin{minipage}[t]{\linewidth}\raggedright
(X)~=~Address~of~pathlist.\\
(A)~=~Access~mode~(D~S~PE~PW~PR~E~W~R)
\end{minipage} \\
OUTPUT: & \begin{minipage}[t]{\linewidth}\raggedright
(X)~=~Updated~past~pathlist~(trailing~spaces~skipped).\\
(A)~=~Path~number.
\end{minipage} \\
ERROR OUTPUT: & \begin{minipage}[t]{\linewidth}\raggedright
(CC)~=~C~bit~set.\\
(B)~=~Appropriate~error~code.
\end{minipage} \\
\end{longtable}
}

Opens a path to an existing file or device as specified by the pathlist.
A path number is returned which is used in subsequent service requests
to identify the file.

The access mode parameter specifies which subsequent read and/or write
operation are permitted as follows:

1~=~read~mode\\
2~=~write~mode\\
3~=~update~mode~(both~read~and~write)\\

Update mode can be slightly slower because pre-reading of sectors may be
required for random access of bytes within sectors. The access mode must
conform to the access permission attributes associated with the file or
device (see \hyperref[i.create]{I\$Create}). Only the owner may access a
file unless the appropiate ``public permit'' bits are set.

Files can be opened by several processes (users) simultaneously. Devices
have an attribute that specifies whether or not they are sharable on an
individual basis.

NOTES:

If the execution bit is set in the access mode, OS-9 will begin
searching for the file in the working execution directory (unless the
pathlist begins with a slash).

The sharable bit (bit 6) in the access mode can not lock other users out
of file in OS-9 Level I. It is present only for upward compatibility
with OS-9 Level II.

Directory files may be read or written if the D bit (bit 7) is set in
the access mode.

\subsection{I\$Read - Read data from a file or device}\label{i.read}

{\def\LTcaptype{none} % do not increment counter
\begin{longtable}[]{@{}
  >{\raggedright\arraybackslash}p{(\linewidth - 2\tabcolsep) * \real{0.2800}}
  >{\raggedright\arraybackslash}p{(\linewidth - 2\tabcolsep) * \real{0.7200}}@{}}
\toprule\noalign{}
\endhead
\bottomrule\noalign{}
\endlastfoot
ASSEMBLER CALL: & OS9 I\$Read \\
MACHINE CODE: & 103F 89 \\
INPUT: & \begin{minipage}[t]{\linewidth}\raggedright
(X)~=~Address~to~store~data.\\
(Y)~=~Number~of~bytes~to~read.\\
(A)~=~Path~number.
\end{minipage} \\
OUTPUT: & (Y) = Number of bytes actually read. \\
ERROR OUTPUT: & \begin{minipage}[t]{\linewidth}\raggedright
(CC)~=~C~bit~set.\\
(B)~=~Appropriate~error~code.
\end{minipage} \\
\end{longtable}
}

Reads a specified number of bytes from the path number given. The path
must previously have been opened in READ or UPDATE mode. The data is
returned exactly as read from the file/device without additional
processing or editing such as backspace, line delete, end-of-file, etc.

After all data in a file has been read, the next
\hyperref[i.read]{I\$Read} service request will return an end of file
error.

NOTES:

The keyboard abort, keyboard interrupt, and end-of-file characters may
be filtered out of the input data on SCFMAN-type devices unless the
coresponding entries in the path descriptor have been set to zero. It
may be desirable to modify the device descriptor so that these values in
the path descriptor are initialized to zero when the path is opened.

The number of bytes requested will be read unless:

A.~An~end-of-file~occurs\\
B.~An~end-of-record~occurs~(SCFMAN~only)\\
C.~An~error~condition~occurs.\\

\subsection{I\$ReadLn - Read a text line with editing}\label{i.readln}

{\def\LTcaptype{none} % do not increment counter
\begin{longtable}[]{@{}
  >{\raggedright\arraybackslash}p{(\linewidth - 2\tabcolsep) * \real{0.2800}}
  >{\raggedright\arraybackslash}p{(\linewidth - 2\tabcolsep) * \real{0.7200}}@{}}
\toprule\noalign{}
\endhead
\bottomrule\noalign{}
\endlastfoot
ASSEMBLER CALL: & OS9 I\$ReadLn \\
MACHINE CODE: & 103F 8B \\
INPUT: & \begin{minipage}[t]{\linewidth}\raggedright
(X)~=~Address~to~store~data.\\
(Y)~=~Number~of~bytes~to~read.\\
(A)~=~Path~number.
\end{minipage} \\
OUTPUT: & (Y) = Actual number of bytes read. \\
ERROR OUTPUT: & \begin{minipage}[t]{\linewidth}\raggedright
(CC)~=~C~bit~set.\\
(B)~=~Appropriate~error~code.
\end{minipage} \\
\end{longtable}
}

This system call is the same as \hyperref[i.read]{I\$Read} except it
reads data from the input file or device until a carriage return
character is encountered or until the maximum byte count specified is
reached, and that line editing will occur on SCFMAN-type devices. Line
editing refers to backspace, line delete, echo automatic line feed, etc.

SCFMAN requires that the last byte entered be an end-of-record character
(normally carriage return). If more data is entered than the maximum
specified, it will not be accepted and a PD.OVF character (normally
bell) will be echoed.

After all data in the file has been read, the next
\hyperref[i.readln]{I\$ReadLn} service request will return an end of
file error.

NOTE: For more information on line editing, see 7.1.

\subsection{I\$Seek - Reposition the logical file pointer}

{\def\LTcaptype{none} % do not increment counter
\begin{longtable}[]{@{}
  >{\raggedright\arraybackslash}p{(\linewidth - 2\tabcolsep) * \real{0.2800}}
  >{\raggedright\arraybackslash}p{(\linewidth - 2\tabcolsep) * \real{0.7200}}@{}}
\toprule\noalign{}
\endhead
\bottomrule\noalign{}
\endlastfoot
ASSEMBLER CALL: & OS9 I\$Seek \\
MACHINE CODE: & 103F 88 \\
INPUT: & \begin{minipage}[t]{\linewidth}\raggedright
(A)~=~Path~number.\\
(X)~=~M.S.~16~bits~of~desired~file~position.\\
(U)~=~L.S.~16~bits~of~desired~file~position.
\end{minipage} \\
OUTPUT: & None. \\
ERROR OUTPUT: & \begin{minipage}[t]{\linewidth}\raggedright
(CC)~=~C~bit~set.\\
(B)~=~Appropriate~error~code.
\end{minipage} \\
\end{longtable}
}

This system call repositions the path\textquotesingle s ``file
pointer''; which is the 32-bit addres of the next byte in the file to be
read from or written to.

A seek may be performed to any value even if the file is not large
enough. Subsequent WRITEs will automatically expand the file to the
required size (if possible), but READs will return an end-of-file
condition. Note that a SEEK to address zero is the same as a ``rewind''
operation.

Seeks to non-random access devices are usually ignored and return
without error.

\subsection{I\$SetStt - Set file/device status}\label{i.setstt}

{\def\LTcaptype{none} % do not increment counter
\begin{longtable}[]{@{}
  >{\raggedright\arraybackslash}p{(\linewidth - 2\tabcolsep) * \real{0.2800}}
  >{\raggedright\arraybackslash}p{(\linewidth - 2\tabcolsep) * \real{0.7200}}@{}}
\toprule\noalign{}
\endhead
\bottomrule\noalign{}
\endlastfoot
ASSEMBLER CALL: & OS9 I\$SetStt \\
MACHINE CODE: & 103F 8E \\
INPUT: & \begin{minipage}[t]{\linewidth}\raggedright
(A)~=~Path~number.\\
(B)~Function~code.\\
(Other~registers~depend~upon~function~code)
\end{minipage} \\
OUTPUT: & (depends upon function code) \\
ERROR OUTPUT: & \begin{minipage}[t]{\linewidth}\raggedright
(CC)~=~C~bit~set.\\
(B)~=~Appropriate~error~code.
\end{minipage} \\
\end{longtable}
}

This system is a ``wild card'' call used to handle individual device
parameters that:

\begin{enumerate}
\def\labelenumi{\alph{enumi}.}
\item
  are not uniform on all devices
\item
  are highly hardware dependent
\item
  need to be user-changeable
\end{enumerate}

The exact operation of this call depends on the device driver and file
manager associated with the path. A typical use is to set a
terminal\textquotesingle s parameters for backspace character, delete
character, echo on/off, null padding, paging, etc. It is commonly used
in conjunction with the \hyperref[i.getstt]{I\$GetStt} service request
which is used to read the device operating parameters. Below are
presently defined function codes:

{\def\LTcaptype{none} % do not increment counter
\begin{longtable}[]{@{}
  >{\raggedright\arraybackslash}p{(\linewidth - 4\tabcolsep) * \real{0.2182}}
  >{\raggedright\arraybackslash}p{(\linewidth - 4\tabcolsep) * \real{0.1455}}
  >{\raggedright\arraybackslash}p{(\linewidth - 4\tabcolsep) * \real{0.6364}}@{}}
\toprule\noalign{}
\begin{minipage}[b]{\linewidth}\raggedright
MNEMONIC
\end{minipage} & \begin{minipage}[b]{\linewidth}\raggedright
CODE
\end{minipage} & \begin{minipage}[b]{\linewidth}\raggedright
FUNCTION
\end{minipage} \\
\midrule\noalign{}
\endhead
\bottomrule\noalign{}
\endlastfoot
SS.Opt & \$0 & Write the 32 byte option section of the path
descriptor \\
SS.Size & \$2 & Set the file size (RBF) \\
SS.Reset & \$3 & Restore head to track zero (RBF) \\
SS.WRT & \$4 & Write (format) track (RBF) \\
SS.Feed & \$9 & Issue Form Feed (SCF) \\
SS.FRZ & \$A & Freeze DD. Information (RBF) \\
SS.SPT & \$B & Set Sectors per track (RBF) \\
SS.SQD & \$C & Sequence down disk drive (RBF) \\
SS.Dcmd & \$D & Direct command to hard disk controller (RBF) \\
\end{longtable}
}

Codes 128 through 255 their parameter passing conventions are user
definable (see the sections of this manual on writing device drivers).
The function code and register stack are passed to the device driver.

{\def\LTcaptype{none} % do not increment counter
\begin{longtable}[]{@{}
  >{\raggedright\arraybackslash}p{(\linewidth - 2\tabcolsep) * \real{0.3000}}
  >{\raggedright\arraybackslash}p{(\linewidth - 2\tabcolsep) * \real{0.7000}}@{}}
\toprule\noalign{}
\endhead
\bottomrule\noalign{}
\endlastfoot
SS.Opt (code 0): & Write option section of the path descriptor. \\
INPUT: & \begin{minipage}[t]{\linewidth}\raggedright
(A)~=~Path~number\\
(B)~=~Function~code~0\\
(X)~=~Address~of~a~32~byte~status~packet.
\end{minipage} \\
OUTPUT: & None. \\
\end{longtable}
}

This setstat function writes the option section of the path descriptor
from the 32 byte status packet pointed to by the X register. It is
typically used to set the device operating parameters, such as echo,
auto line feed, etc.

{\def\LTcaptype{none} % do not increment counter
\begin{longtable}[]{@{}
  >{\raggedright\arraybackslash}p{(\linewidth - 2\tabcolsep) * \real{0.3000}}
  >{\raggedright\arraybackslash}p{(\linewidth - 2\tabcolsep) * \real{0.7000}}@{}}
\toprule\noalign{}
\endhead
\bottomrule\noalign{}
\endlastfoot
SS.Size (code 2): & Set the file size (RBFMAN-type devices) \\
INPUT: & \begin{minipage}[t]{\linewidth}\raggedright
(A)~=~Path~number\\
(B)~=~Function~code~2\\
(X)~=~M.S.~16~bits~of~desired~file~size.\\
(U)~=~L.S.~16~bits~of~desired~file~size.
\end{minipage} \\
OUTPUT: & None. \\
\end{longtable}
}

This setstat function is used to change the file\textquotesingle s size.

{\def\LTcaptype{none} % do not increment counter
\begin{longtable}[]{@{}
  >{\raggedright\arraybackslash}p{(\linewidth - 2\tabcolsep) * \real{0.3000}}
  >{\raggedright\arraybackslash}p{(\linewidth - 2\tabcolsep) * \real{0.7000}}@{}}
\toprule\noalign{}
\endhead
\bottomrule\noalign{}
\endlastfoot
SS.Reset (code 3): & Restore head to track zero. \\
INPUT: & \begin{minipage}[t]{\linewidth}\raggedright
(A)~=~Path~number\\
(B)~=~Function~code~3
\end{minipage} \\
OUTPUT: & None. \\
\end{longtable}
}

Home disk head to track zero. Used for formatting and for error
recovery.

{\def\LTcaptype{none} % do not increment counter
\begin{longtable}[]{@{}
  >{\raggedright\arraybackslash}p{(\linewidth - 2\tabcolsep) * \real{0.3000}}
  >{\raggedright\arraybackslash}p{(\linewidth - 2\tabcolsep) * \real{0.7000}}@{}}
\toprule\noalign{}
\endhead
\bottomrule\noalign{}
\endlastfoot
SS.WTrk (code 4): & Write track \\
INPUT: & \begin{minipage}[t]{\linewidth}\raggedright
(A)~=~Path~number\\
(B)~=~Function~code~4\\
(X)~=~Address~of~track~buffer.\\
(U)~=~Track~number~(L.S.~8~bits)\\
(Y)~=~Side/density\\
\strut ~~~~~~~~~Bit~B0~=~SIDE~(0~=~side~zero,~1~=~side~one)\\
\strut ~~~~~Bit~B1~=~DENSITY~(0~=~single,~1~=~double)
\end{minipage} \\
OUTPUT: & None. \\
\end{longtable}
}

This code causes a format track (most floppy disks) operation to occur.
For hard disks or floppy disks with a ``format entire disk'' command,
this command should format the entire media only when the track number
equals zero.

{\def\LTcaptype{none} % do not increment counter
\begin{longtable}[]{@{}
  >{\raggedright\arraybackslash}p{(\linewidth - 2\tabcolsep) * \real{0.3000}}
  >{\raggedright\arraybackslash}p{(\linewidth - 2\tabcolsep) * \real{0.7000}}@{}}
\toprule\noalign{}
\endhead
\bottomrule\noalign{}
\endlastfoot
SS.FRZ (code \$A): & Freeze DD. Information \\
INPUT: & none \\
OUTPUT: & none \\
\end{longtable}
}

Inhibits the reading of identification sector (LSN 0) to DD.xxx
variables (that define disk formats) so non-standard disks may be read.

{\def\LTcaptype{none} % do not increment counter
\begin{longtable}[]{@{}
  >{\raggedright\arraybackslash}p{(\linewidth - 2\tabcolsep) * \real{0.3000}}
  >{\raggedright\arraybackslash}p{(\linewidth - 2\tabcolsep) * \real{0.7000}}@{}}
\toprule\noalign{}
\endhead
\bottomrule\noalign{}
\endlastfoot
SS.SPT (code \$B): & Set Sectors Per Track \\
INPUT: & X = new sectors per track \\
\end{longtable}
}

\subsection{I\$Write - Write data to file or device}\label{i.write}

{\def\LTcaptype{none} % do not increment counter
\begin{longtable}[]{@{}
  >{\raggedright\arraybackslash}p{(\linewidth - 2\tabcolsep) * \real{0.2800}}
  >{\raggedright\arraybackslash}p{(\linewidth - 2\tabcolsep) * \real{0.7200}}@{}}
\toprule\noalign{}
\endhead
\bottomrule\noalign{}
\endlastfoot
ASSEMBLER CALL: & OS9 I\$Write \\
MACHINE CODE: & 103F 8A \\
INPUT: & \begin{minipage}[t]{\linewidth}\raggedright
(X)~=~Address~of~data~to~write.\\
(Y)~=~Number~of~bytes~to~write.\\
(A)~=~Path~number.
\end{minipage} \\
OUTPUT: & (Y) = Number of bytes actually written. \\
ERROR OUTPUT: & \begin{minipage}[t]{\linewidth}\raggedright
(CC)~=~C~bit~set.\\
(B)~=~Appropriate~error~code.
\end{minipage} \\
\end{longtable}
}

\hyperref[i.write]{I\$Write} outputs one or more bytes to a file or
device associated with the path number specified. The path must have
been OPENed or CREATEd in the WRITE or UPDATE access modes.

Data is written to the file or device without processing or editing. If
data is written past the present end-of-file, the file is automatically
expanded.

\subsection{I\$WritLn - Write line of text with editing}\label{i.writln}

{\def\LTcaptype{none} % do not increment counter
\begin{longtable}[]{@{}
  >{\raggedright\arraybackslash}p{(\linewidth - 2\tabcolsep) * \real{0.2800}}
  >{\raggedright\arraybackslash}p{(\linewidth - 2\tabcolsep) * \real{0.7200}}@{}}
\toprule\noalign{}
\endhead
\bottomrule\noalign{}
\endlastfoot
ASSEMBLER CALL: & OS9 I\$WritLn \\
MACHINE CODE: & 103F 8C \\
INPUT: & \begin{minipage}[t]{\linewidth}\raggedright
(X)~=~Address~of~data~to~write.\\
(Y)~=~Maximum~number~of~bytes~to~write.\\
(A)~=~Path~number.
\end{minipage} \\
OUTPUT: & (Y) = Actual number of bytes written. \\
ERROR OUTPUT: & \begin{minipage}[t]{\linewidth}\raggedright
(CC)~=~C~bit~set.\\
(B)~=~Appropriate~error~code.
\end{minipage} \\
\end{longtable}
}

This system call is similar to \hyperref[i.write]{I\$Write} except it
writes data until a carriage return character is encountered. Line
editing is also activated for character-oriented devices such as
terminals, printers, etc. The line editing refers to auto line feed,
null padding at end of line, etc.

For more information about line editing, see section 7.1.

\section{Memory Module Diagrams}

\begin{figure}
\centering
\begin{verbatim}
MODULE          EXECUTABLE MEMORY MODULE FORMAT
OFFSET

        +------------------------------+  ---+--------+---
 $00    |                              |     |        |
        +--    Sync Bytes ($87CD)    --+     |        |
 $01    |                              |     |        |
        +------------------------------+     |        |
 $02    |                              |     |        |
        +--   Module Size (bytes)    --+     |        |
 $03    |                              |     |        |
        +------------------------------+     |        |
 $04    |                              |     |        |
        +--   Module Name Offset     --+   header     |
 $05    |                              |   parity     |
        +------------------------------+     |        |
 $06    |     Type     |   Language    |     |        |
        +------------------------------+     |        |
 $07    |  Attributes  |   Revision    |     |        |
        +------------------------------+  ---+--    module
 $08    |     Header Parity Check      |             CRC
        +------------------------------+              |
 $09    |                              |              |
        +--     Execution Offset     --+              |
 $0A    |                              |              |
        +------------------------------+              |
 $0B    |                              |              |
        +--  Permanent Storage Size  --+              |
 $0C    |                              |              |
        +------------------------------+              |
 $0D    |                              |              |
        |  (Add'l optional header      |              |
        |   extensions located here    |              |
        |                              |              |
        |  .  .  .  .  .  .  .  .  .   |              |
        |                              |              |
        |                              |              |
        |       Module Body            |              |
        | object code, constants, etc. |              |
        |                              |              |
        |                              |              |
        +------------------------------+              |
        |                              |              |
        +--                          --+              |
        |       CRC Check Value        |              |
        +--                          --+              |
        |                              |              |
        +------------------------------+  ------------+---
\end{verbatim}
\caption{}
\end{figure}

\begin{figure}
\centering
\begin{verbatim}
MODULE     DEVICE DESCRIPTOR MODULE FORMAT
OFFSET
           +-----------------------------+  ---+--------+---
  $0       |                             |     |        |
           +--   Sync Bytes ($87CD)    --+     |        |
  $1       |                             |     |        |
           +-----------------------------+     |        |
  $2       |                             |     |        |
           +--   Module Size (bytes)   --+     |        |
  $3       |                             |     |        |
           +-----------------------------+     |        |
  $4       |                             |     |        |
           +-- Offset to Module Name   --+   header     |
  $5       |                             |   parity     |
           +-----------------------------+     |        |
  $6       | $F (TYPE)   |  $1 (LANG)    |     |        |
           +-----------------------------+     |        |
  $7       | Attributes  |   Revision    |     |        |
           +-----------------------------+  ---+--    module
  $8       |  Header Parity Check        |             CRC
           +-----------------------------+              |
  $9       |                             |              |
           +--  Offset to File Manager --+              |
  $A       |         Name String         |              |
           +-----------------------------+              |
  $B       |                             |              |
           +-- Offset to Device Driver --+              |
  $C       |         Name String         |              |
           +-----------------------------+              |
  $D       |        Mode Byte            |              |
           +-----------------------------+              |
  $E       |                             |              |
           +--    Device Controller    --+              |
  $F       | Absolute Physical Address   |              |
           +--     (24 bit)            --+              |
 $10       |                             |              |
           +-----------------------------+              |
 $11       |   Option Table Size         |              |
           +-----------------------------+              |
$12,$12+N  |    (Option Table)           |              |
           |                             |              |
           | . . . . . . . . . . . . . . |              |
           |                             |              |
           |    (Name Strings etc)       |              |
           +-----------------------------+              |
           |                             |              |
           +--                         --+              |
           |    CRC Check Value          |              |
           +--                         --+              |
           |                             |              |
           +-----------------------------+  ------------+---
\end{verbatim}
\caption{}
\end{figure}

\begin{figure}
\centering
\begin{verbatim}
MODULE      CONFIGURATION MODULE FORMAT
OFFSET
        +------------------------------+  ---+--------+---
 $00    |                              |     |        |
        +--    Sync Bytes ($87CD)    --+     |        |
 $01    |                              |     |        |
        +------------------------------+     |        |
 $02    |                              |     |        |
        +--   Module Size (bytes)    --+     |        |
 $03    |                              |     |        |
        +------------------------------+     |        |
 $04    |                              |     |        |
        +--   Module Name Offset     --+   header     |
 $05    |                              |   parity     |
        +------------------------------+     |        |
 $06    |  $C (TYPE)   |  0  (LANG)    |     |        |
        +------------------------------+     |        |
 $07    |  Attributes  |   Revision    |     |        |
        +------------------------------+  ---+--    module
 $08    |     Header Parity Check      |             CRC
        +------------------------------+              |
 $09    |                              |              |
        +--     Forced Limit of Top  --+              |
 $0A    |         of Free RAM          |              |
        +--                          --+              |
 $0B    |                              |              |
        +------------------------------+              |
 $0C    | # IRQ Polling Table Entries  |              |
        +------------------------------+              |
 $0D    |    # Device Table Entries    |              |
        +------------------------------+              |
 $0E    |                              |              |
        +--    Offset to Startup     --+              |
 $0F    |      Module Name String      |              |
        +------------------------------+              |
 $10    |                              |              |
        +-- Offset to Default Mass-  --+              |
 $11    |  Storage Device Name String  |              |
        +------------------------------+              |
 $12    |                              |              |
        +--    Offset to Initial     --+              |
 $13    |        Standard Path         |              |
        +------------------------------+              |
 $14    |                              |              |
        +--   Offset to Bootstrap    --+              |
 $15    |     Module Name String       |              |
        +------------------------------+              |
 $16-n  |        Name Strings          |              |
        +------------------------------+              |
        |                              |              |
        +--                          --+              |
        |      CRC Check Value         |              |
        +--                          --+              |
        |                              |              |
        +------------------------------+  ------------+---
\end{verbatim}
\caption{}
\end{figure}

\section{Standard Floppy Disk Formats}

\begin{longtable}[]{@{}
  >{\raggedright\arraybackslash}p{(\linewidth - 4\tabcolsep) * \real{0.3571}}
  >{\raggedright\arraybackslash}p{(\linewidth - 4\tabcolsep) * \real{0.3214}}
  >{\raggedright\arraybackslash}p{(\linewidth - 4\tabcolsep) * \real{0.3214}}@{}}
\caption{Single Density Floppy Disk Format}\tabularnewline
\toprule\noalign{}
\endfirsthead
\endhead
\bottomrule\noalign{}
\endlastfoot
SIZE & 5" & 8" \\
DENSITY & SINGLE & SINGLE \\
\#TRACKS & 35 & 77 \\
\#SECTORS/TRACK & 10 & 16 \\
BYTES/TRACK (UNFORMATTED) & 3125 & 5208 \\
\end{longtable}

\begin{longtable}[]{@{}
  >{\raggedright\arraybackslash}p{(\linewidth - 4\tabcolsep) * \real{0.3571}}
  >{\raggedright\arraybackslash}p{(\linewidth - 4\tabcolsep) * \real{0.3214}}
  >{\raggedright\arraybackslash}p{(\linewidth - 4\tabcolsep) * \real{0.3214}}@{}}
\caption{Double Density Floppy Disk Format}\tabularnewline
\toprule\noalign{}
\endfirsthead
\endhead
\bottomrule\noalign{}
\endlastfoot
SIZE & 5" & 8" \\
DENSITY & DOUBLE & DOUBLE \\
\#TRACKS & 35 & 77 \\
\#SECTORS/TRACK & 16 & 28 \\
BYTES/TRACK (UNFORMATTED) & 6250 & 10,416 \\
\end{longtable}

\section{Service Request Summary}

\begin{longtable}[]{@{}
  >{\raggedright\arraybackslash}p{(\linewidth - 6\tabcolsep) * \real{0.1875}}
  >{\raggedright\arraybackslash}p{(\linewidth - 6\tabcolsep) * \real{0.1562}}
  >{\raggedright\arraybackslash}p{(\linewidth - 6\tabcolsep) * \real{0.5625}}
  >{\raggedright\arraybackslash}p{(\linewidth - 6\tabcolsep) * \real{0.0937}}@{}}
\caption{User Mode Service Requests}\tabularnewline
\toprule\noalign{}
\begin{minipage}[b]{\linewidth}\raggedright
Code
\end{minipage} & \begin{minipage}[b]{\linewidth}\raggedright
Mnemonic
\end{minipage} & \begin{minipage}[b]{\linewidth}\raggedright
Function
\end{minipage} & \begin{minipage}[b]{\linewidth}\raggedright
Page
\end{minipage} \\
\midrule\noalign{}
\endfirsthead
\toprule\noalign{}
\begin{minipage}[b]{\linewidth}\raggedright
Code
\end{minipage} & \begin{minipage}[b]{\linewidth}\raggedright
Mnemonic
\end{minipage} & \begin{minipage}[b]{\linewidth}\raggedright
Function
\end{minipage} & \begin{minipage}[b]{\linewidth}\raggedright
Page
\end{minipage} \\
\midrule\noalign{}
\endhead
\bottomrule\noalign{}
\endlastfoot
103F 00 & F\$LINK & Link to memory module. & \\
103F 01 & F\$LOAD & Load module(s) from a file. & \\
103F 02 & F\$Unlink & Unlink a module. & \\
103F 03 & F\$Fork & Create a new process. & \\
103F 04 & F\$Wait & Wait for child process to die. & \\
103F 05 & F\$Chain & Load and execute a new primary module & \\
103F 06 & F\$Exit & Terminate the calling process. & \\
103F 07 & F\$Mem & Resize data memory area, & \\
103F 08 & F\$Send & Send a signal to another process, & \\
103F 09 & F\$ICPT & Set up a signal intercept trap. & \\
103F 0A & F\$Sleep & Put calling process to sleep. & \\
103F 0C & F\$ID & Get process ID / user ID & \\
103F 0D & F\$SPrior & Set process priority. & \\
103F 0E & F\$SSWI & Set SWI vector. & \\
103F 0F & F\$PErr & Print error message. & \\
103F 10 & F\$PrsNam & Parse a path name, & \\
103F 11 & F\$CmpNam & Compare two names & \\
103F 12 & F\$SchBit & Search bit map for a free area & \\
103F 13 & F\$AllBit & Set bits in an allocation bit map & \\
103F 14 & F\$DelBit & Deallocate in a bit map & \\
103F 15 & F\$Time & Get system date and time. & \\
103F 16 & F\$STime & Set system date and time. & \\
103F 17 & F\$CRC & Compute CRC & \\
103F 18 & F\$GPrDsc & Get Process Descriptor copy & \\
103F 19 & F\$GBlkMp & Get system Block Map copy & \\
103F 1A & F\$GModDr & Get Module Directory copy & \\
103F 1B & F\$CpyMem & Copy external Memory & \\
103F 1C & F\$SUser & Set User ID number & \\
103F 1D & F\$UnLoad & Unlink module by name & \\
\end{longtable}

\begin{longtable}[]{@{}
  >{\raggedright\arraybackslash}p{(\linewidth - 6\tabcolsep) * \real{0.1875}}
  >{\raggedright\arraybackslash}p{(\linewidth - 6\tabcolsep) * \real{0.1562}}
  >{\raggedright\arraybackslash}p{(\linewidth - 6\tabcolsep) * \real{0.5625}}
  >{\raggedright\arraybackslash}p{(\linewidth - 6\tabcolsep) * \real{0.0937}}@{}}
\caption{System Mode Privileged Service Requests}\tabularnewline
\toprule\noalign{}
\begin{minipage}[b]{\linewidth}\raggedright
Code
\end{minipage} & \begin{minipage}[b]{\linewidth}\raggedright
Mnemonic
\end{minipage} & \begin{minipage}[b]{\linewidth}\raggedright
Function
\end{minipage} & \begin{minipage}[b]{\linewidth}\raggedright
Page
\end{minipage} \\
\midrule\noalign{}
\endfirsthead
\toprule\noalign{}
\begin{minipage}[b]{\linewidth}\raggedright
Code
\end{minipage} & \begin{minipage}[b]{\linewidth}\raggedright
Mnemonic
\end{minipage} & \begin{minipage}[b]{\linewidth}\raggedright
Function
\end{minipage} & \begin{minipage}[b]{\linewidth}\raggedright
Page
\end{minipage} \\
\midrule\noalign{}
\endhead
\bottomrule\noalign{}
\endlastfoot
103F 28 & F\$SRqMem & System memory request & \\
103F 29 & F\$SRtMem & System memory return & \\
103F 2A & F\$IRQ & Add or remove device from IRQ table. & \\
103F 2B & F\$IOQU & Enter I/O queue & \\
103F 2C & F\$AProc & Insert process in active process queue & \\
103F 2D & F\$NProc & Start next process & \\
103F 2E & F\$VModul & Validate module & \\
103F 2F & F\$Find64 & Find a 64 byte memory block & \\
103F 30 & F\$All64 & Allocate a 64 byte memory block & \\
103F 31 & F\$Ret64 & Deallocate a 64 byte memory block & \\
103F 32 & F\$SSVC & Install function request & \\
103F 33 & F\$IODel & Delete I/O device from system & \\
103F 34 & F\$SLink & System Link & \\
103F 35 & F\$Boot & Bootstrap system & \\
103F 36 & F\$BtMem & Bootstrap Memory request & \\
103F 37 & F\$GProcP & Get Process Pointer & \\
103F 38 & F\$Move & Move data (low bound first) & \\
103F 39 & F\$AllRAM & Allocate RAM blocks & \\
103F 3A & F\$AllImg & Allocate Image RAM blocks & \\
103F 3B & F\$DelImg & Deallocate Image RAM blocks & \\
103F 3C & F\$SetImg & Set process DAT Image & \\
103F 3D & F\$FreeLB & get Free Low block & \\
103F 3E & F\$FreeHB & get Free High block & \\
103F 3F & F\$AllTsk & Allocate process Task number & \\
103F 40 & F\$DelTsk & Deallocate process Task number & \\
103F 41 & F\$SetTsk & Set process Task DAT registers & \\
103F 42 & F\$ResTsk & Reserve Task number & \\
103F 43 & F\$RelTsk & Release Task number & \\
103F 44 & F\$DATLog & Convert DAT block/offset to Logical Addr & \\
103F 45 & F\$DATTmp & Make Temporary DAT image & \\
103F 46 & F\$LDAXY & Load A {[}X, {[}Y{]} {]} & \\
103F 47 & F\$LDAXYP & Load A {[}X+, {[}Y{]} {]} & \\
103F 48 & F\$LDDDXY & Load D {[}D+X, {[}Y{]} {]} & \\
103F 49 & F\$LDABX & Load A from 0,1 in task B & \\
103F 4A & F\$STABX & Store A at 0,X in task B & \\
103F 4B & F\$AllPrc & Allocate Process descriptor & \\
103F 4C & F\$DelPrc & Deallocate Process descriptor & \\
103F 4D & F\$ELink & Link using module directory Entry & \\
103F 4E & F\$FModul & Find Module directory entry & \\
103F 4F & F\$MapBlk & Map specific Block & \\
103F 50 & F\$ClrBlk & Clear specific Block & \\
103F 51 & F\$DelRam & Deallocate RAM blocks & \\
\end{longtable}

\begin{longtable}[]{@{}
  >{\raggedright\arraybackslash}p{(\linewidth - 6\tabcolsep) * \real{0.1875}}
  >{\raggedright\arraybackslash}p{(\linewidth - 6\tabcolsep) * \real{0.1562}}
  >{\raggedright\arraybackslash}p{(\linewidth - 6\tabcolsep) * \real{0.5625}}
  >{\raggedright\arraybackslash}p{(\linewidth - 6\tabcolsep) * \real{0.0937}}@{}}
\caption{Input/Output Service Requests}\tabularnewline
\toprule\noalign{}
\begin{minipage}[b]{\linewidth}\raggedright
Code
\end{minipage} & \begin{minipage}[b]{\linewidth}\raggedright
Mnemonic
\end{minipage} & \begin{minipage}[b]{\linewidth}\raggedright
Function
\end{minipage} & \begin{minipage}[b]{\linewidth}\raggedright
Page
\end{minipage} \\
\midrule\noalign{}
\endfirsthead
\toprule\noalign{}
\begin{minipage}[b]{\linewidth}\raggedright
Code
\end{minipage} & \begin{minipage}[b]{\linewidth}\raggedright
Mnemonic
\end{minipage} & \begin{minipage}[b]{\linewidth}\raggedright
Function
\end{minipage} & \begin{minipage}[b]{\linewidth}\raggedright
Page
\end{minipage} \\
\midrule\noalign{}
\endhead
\bottomrule\noalign{}
\endlastfoot
103F 80 & I\$Attach & Attach a new device to the system. & \\
103F 81 & I\$Detach & Remove a device from the system. & \\
103F 82 & I\$Dup & Duplicate a path. & \\
103F 83 & I\$Create & Create a path to a new file. & \\
103F 84 & I\$Open & Open a path to a file or device & \\
103F 85 & I\$MakDir & Make a new directory & \\
103F 86 & I\$ChgDir & Change working directory. & \\
103F 87 & I\$Delete & Delete a file. & \\
103F 88 & I\$Seek & Reposition the logical file pointer & \\
103F 89 & I\$Read & Read data from a file or device & \\
103F 8A & I\$Write & Write Data to File or Device & \\
103F 8B & I\$ReadLn & Read a text line with editing. & \\
103F 8C & I\$WritLn & Write Line of Text with Editing & \\
103F 8D & I\$GetStt & Get file device status. & \\
103F 8E & I\$SetStt & Set file/device status & \\
103F 8F & I\$Close & Close a path to a file/device. & \\
103F 90 & I\$DeletX & Delete a file & \\
\end{longtable}

\begin{longtable}[]{@{}l@{}}
\caption{Standard I/O Paths}\tabularnewline
\toprule\noalign{}
\endfirsthead
\endhead
\bottomrule\noalign{}
\endlastfoot
0 = Standard Input \\
1 = Standard Output \\
2 = Standard Error Output \\
\end{longtable}

\begin{longtable}[]{@{}
  >{\raggedright\arraybackslash}p{(\linewidth - 2\tabcolsep) * \real{0.1667}}
  >{\raggedright\arraybackslash}p{(\linewidth - 2\tabcolsep) * \real{0.8333}}@{}}
\caption{Module Types}\tabularnewline
\toprule\noalign{}
\endfirsthead
\endhead
\bottomrule\noalign{}
\endlastfoot
\$1 & Program \\
\$2 & Subroutine module \\
\$3 & Multi-module \\
\$4 & Data module \\
\$C & System Module \\
\$D & File Manager \\
\$E & Device Driver \\
\$F & Device Descriptor \\
\end{longtable}

\begin{longtable}[]{@{}
  >{\raggedright\arraybackslash}p{(\linewidth - 2\tabcolsep) * \real{0.4167}}
  >{\raggedright\arraybackslash}p{(\linewidth - 2\tabcolsep) * \real{0.5833}}@{}}
\caption{File Access Codes}\tabularnewline
\toprule\noalign{}
\endfirsthead
\endhead
\bottomrule\noalign{}
\endlastfoot
READ & \$01 \\
WRITE & \$02 \\
UPDATE & READ + WRITE \\
EXEC & \$04 \\
PREAD & \$08 \\
PWRIT & \$10 \\
PEXEC & \$20 \\
SHARE & \$40 \\
DIR & \$80 \\
\end{longtable}

\begin{longtable}[]{@{}
  >{\raggedright\arraybackslash}p{(\linewidth - 2\tabcolsep) * \real{0.2500}}
  >{\raggedright\arraybackslash}p{(\linewidth - 2\tabcolsep) * \real{0.7500}}@{}}
\caption{Module Languages}\tabularnewline
\toprule\noalign{}
\endfirsthead
\endhead
\bottomrule\noalign{}
\endlastfoot
\$0 & Data \\
\$1 & 6809 Object code \\
\$2 & BASIC09 I-code \\
\$3 & Pascal P-Code \\
\$4 & C I-code \\
\$5 & Cobol I-code \\
\$6 & Fortan I-code \\
\$7 & 6309 Object code \\
\end{longtable}

\begin{longtable}[]{@{}
  >{\raggedright\arraybackslash}p{(\linewidth - 2\tabcolsep) * \real{0.1667}}
  >{\raggedright\arraybackslash}p{(\linewidth - 2\tabcolsep) * \real{0.8333}}@{}}
\caption{Module Attributes}\tabularnewline
\toprule\noalign{}
\endfirsthead
\endhead
\bottomrule\noalign{}
\endlastfoot
\$8 & Reentrant \\
\end{longtable}

\section{Error Codes}\label{errorcodes}

\subsection{OS-9 Error Codes}

The error codes are shown both in hexadecimal (first column) and decimal
(second column). Error codes other than those listed are generated by
programming languages or user programs.

{\def\LTcaptype{none} % do not increment counter
\begin{longtable}[]{@{}
  >{\raggedright\arraybackslash}p{(\linewidth - 4\tabcolsep) * \real{0.1200}}
  >{\raggedright\arraybackslash}p{(\linewidth - 4\tabcolsep) * \real{0.1200}}
  >{\raggedright\arraybackslash}p{(\linewidth - 4\tabcolsep) * \real{0.7600}}@{}}
\toprule\noalign{}
\begin{minipage}[b]{\linewidth}\raggedright
HEX
\end{minipage} & \begin{minipage}[b]{\linewidth}\raggedright
DEC
\end{minipage} & \begin{minipage}[b]{\linewidth}\raggedright
\end{minipage} \\
\midrule\noalign{}
\endhead
\bottomrule\noalign{}
\endlastfoot
\$C8 & 200 & PATH TABLE FULL - The file cannot be opened because the
system path table is currently full. \\
\$C9 & 201 & ILLEGAL PATH NUMBER - Number too large or for non-existant
path. \\
\$CA & 202 & INTERRUPT POLLING TABLE FULL \\
\$CB & 203 & ILLEGAL MODE - attempt to perform I/O function of which the
device or file is incapable. \\
\$CC & 204 & DEVICE TABLE FULL - Can\textquotesingle t add another
device \\
\$CD & 205 & ILLEGAL MODULE HEADER - module not loaded because its sync
code, header parity, or CRC is incorrect. \\
\$CE & 206 & MODULE DIRECTORY FULL - Can\textquotesingle t add another
module \\
\$CF & 207 & MEMORY FULL - Level One: not enough contiquous RAM free.
Level Two: process address space full \\
\$D0 & 208 & ILLEGAL SERVICE REQUEST - System call had an illegal code
number \\
\$D1 & 209 & MODULE BUSY - non-sharable module is in use by another
process. \\
\$D2 & 210 & BOUNDARY ERROR - Memory allocation or deallocation request
not on a page boundary. \\
\$D3 & 211 & END OF FILE - End of file encountered on read. \\
\$D4 & 212 & RETURNING NON-ALLOCATED MEMORY - (NOT YOUR MEMORY)
attempted to deallocate memory not previously assigned. \\
\$D5 & 213 & NON-EXISTING SEGMENT - device has damaged file
structure. \\
\$D6 & 214 & NO PERMISSION - file attributes do not permit access
requested. \\
\$D7 & 215 & BAD PATH NAME - syntax error in pathlist (illegal
character, etc.). \\
\$D8 & 216 & PATH NAME NOT FOUND - can\textquotesingle t find pathlist
specified. \\
\$D9 & 217 & SEGMENT LIST FULL - file is too fragmented to be expanded
further. \\
\$DA & 218 & FILE ALREADY EXISTS - file name already appears in current
directory. \\
\$DB & 219 & ILLEGAL BLOCK ADDRESS - device\textquotesingle s file
structure has been damaged. \\
\$DC & 220 & ILLEGAL BLOCK SIZE - device\textquotesingle s file
structure has been damaged. \\
\$DD & 221 & MODULE NOT FOUND - request for link to module not found in
directory. \\
\$DE & 222 & SECTOR OUT OF RANGE - device file structure damaged or
incorrectly formatted. \\
\$DF & 223 & SUICIDE ATTEMPT - request to return memory where your stack
is located. \\
\$E0 & 224 & ILLEGAL PROCESS NUMBER - no such process exists. \\
\$E2 & 226 & NO CHILDREN - can\textquotesingle t wait because process
has no children. \\
\$E3 & 227 & ILLEGAL SWI CODE - must be 1 to 3. \\
\$E4 & 228 & PROCESS ABORTED - process aborted by signal code 2. \\
\$E5 & 229 & PROCESS TABLE FULL - can\textquotesingle t fork now. \\
\$E6 & 230 & ILLEGAL PARAMETER AREA - high and low bounds passed in fork
call are incorrect. \\
\$E7 & 231 & KNOWN MODULE - for internal use only. \\
\$E8 & 232 & INCORRECT MODULE CRC - module has bad CRC value. \\
\$E9 & 233 & SIGNAL ERROR - receiving process has previous unprocessed
signal pending. \\
\$EA & 234 & NON-EXISTENT MODULE - unable to locate module. \\
\$EB & 235 & BAD NAME - illegal name syntax. \\
\$EC & 236 & BAD HEADER - module header parity incorrect \\
\$ED & 237 & RAM FULL - no free system RAM available at this time \\
\$EE & 238 & UNKNOWN PROCESS ID - incorrect process ID number \\
\$EF & 239 & NO TASK NUMBER AVAILABLE - all task numbers in use \\
\end{longtable}
}

\subsection{Device Driver/Hardware Errors}

The following error codes are generated by I/O device drivers, and are
somewhat hardware dependent. Consult manufacturer\textquotesingle s
hardware manual for more details.

{\def\LTcaptype{none} % do not increment counter
\begin{longtable}[]{@{}
  >{\raggedright\arraybackslash}p{(\linewidth - 4\tabcolsep) * \real{0.1200}}
  >{\raggedright\arraybackslash}p{(\linewidth - 4\tabcolsep) * \real{0.1200}}
  >{\raggedright\arraybackslash}p{(\linewidth - 4\tabcolsep) * \real{0.7600}}@{}}
\toprule\noalign{}
\begin{minipage}[b]{\linewidth}\raggedright
HEX
\end{minipage} & \begin{minipage}[b]{\linewidth}\raggedright
DEC
\end{minipage} & \begin{minipage}[b]{\linewidth}\raggedright
\end{minipage} \\
\midrule\noalign{}
\endhead
\bottomrule\noalign{}
\endlastfoot
\$F0 & 240 & UNIT ERROR - device unit does not exist \\
\$F1 & 241 & SECTOR ERROR - sector number is out of range. \\
\$F2 & 242 & WRITE PROTECT - device is write protected. \\
\$F3 & 243 & CRC ERROR - CRC error on read or write verify \\
\$F4 & 244 & READ ERROR - Data transfer error during disk read
operation, or SCF (terminal) input buffer overrun. \\
\$F5 & 245 & WRITE ERROR - hardware error during disk write
operation. \\
\$F6 & 246 & NOT READY - device has "not ready" status. \\
\$F7 & 247 & SEEK ERROR - physical seek to non-existant sector. \\
\$F8 & 248 & MEDIA FULL - insufficient free space on media. \\
\$F9 & 249 & WRONG TYPE - attempt to read incompatible media (i.e.
attempt to read double-side disk on single-side drive) \\
\$FA & 250 & DEVICE BUSY - non-sharable device is in use \\
\$FB & 251 & DISK ID CHANGE - Media was changed with files open \\
\$FC & 252 & RECORD IS LOCKED-OUT - Another process is accessing the
requested record. \\
\$FD & 253 & NON-SHARABLE FILE BUSY - Another process is accessing the
requested file. \\
\$FE & 254 & I/O DEADLOCK ERROR - Two processes are attempting to use
the same two disk areas simultaneously. \\
\end{longtable}
}

\section{Level Two System Service Requests}

\subsection{F\$AllImg - Allocate Image RAM blocks}

{\def\LTcaptype{none} % do not increment counter
\begin{longtable}[]{@{}
  >{\raggedright\arraybackslash}p{(\linewidth - 2\tabcolsep) * \real{0.2800}}
  >{\raggedright\arraybackslash}p{(\linewidth - 2\tabcolsep) * \real{0.7200}}@{}}
\toprule\noalign{}
\endhead
\bottomrule\noalign{}
\endlastfoot
ASSEMBLER CALL: & OS9 F\$AllImg \\
MACHINE CODE: & 103F 3A \\
INPUT: & \begin{minipage}[t]{\linewidth}\raggedright
(A)~=~Beginning~block~number\\
(B)~=~Number~of~blocks\\
(X)~=~Process~Descriptor~pointer
\end{minipage} \\
OUTPUT: & None. \\
ERROR OUTPUT: & \begin{minipage}[t]{\linewidth}\raggedright
(CC)~=~C~bit~set.\\
(B)~=~Appropriate~error~code.
\end{minipage} \\
\end{longtable}
}

Allocates RAM blocks for process DAT image. The blocks do not need to be
contiguous.

This is a privileged system mode service request.

\subsection{F\$AllPrc - Allocate Process descriptor}

{\def\LTcaptype{none} % do not increment counter
\begin{longtable}[]{@{}
  >{\raggedright\arraybackslash}p{(\linewidth - 2\tabcolsep) * \real{0.2800}}
  >{\raggedright\arraybackslash}p{(\linewidth - 2\tabcolsep) * \real{0.7200}}@{}}
\toprule\noalign{}
\endhead
\bottomrule\noalign{}
\endlastfoot
ASSEMBLER CALL: & OS9 F\$AllPrc \\
MACHINE CODE: & 103F 4B \\
INPUT: & none \\
OUTPUT: & (U) = Process Descriptor pointer \\
ERROR OUTPUT: & \begin{minipage}[t]{\linewidth}\raggedright
(CC)~=~C~bit~set.\\
(B)~=~Appropriate~error~code.
\end{minipage} \\
\end{longtable}
}

Allocates and initializes a 512-byte process descriptor.

This is a privileged system mode service request.

\subsection{F\$AllRAM - Allocate RAM blocks}

{\def\LTcaptype{none} % do not increment counter
\begin{longtable}[]{@{}
  >{\raggedright\arraybackslash}p{(\linewidth - 2\tabcolsep) * \real{0.2800}}
  >{\raggedright\arraybackslash}p{(\linewidth - 2\tabcolsep) * \real{0.7200}}@{}}
\toprule\noalign{}
\endhead
\bottomrule\noalign{}
\endlastfoot
ASSEMBLER CALL: & OS9 F\$AllRAM \\
MACHINE CODE: & 103F 39 \\
INPUT: & (B) = Desired block count \\
OUTPUT: & (D) = Beginning RAM block number \\
ERROR OUTPUT: & \begin{minipage}[t]{\linewidth}\raggedright
(CC)~=~C~bit~set.\\
(B)~=~Appropriate~error~code.
\end{minipage} \\
\end{longtable}
}

Searches the Memory Block map for the desired number of contuguous free
RAM blocks.

NOTE: THIS IS A PRIVILEGED SYSTEM MODE SERVICE REQUEST

\subsection{F\$AllTsk - Allocate process Task number}

{\def\LTcaptype{none} % do not increment counter
\begin{longtable}[]{@{}
  >{\raggedright\arraybackslash}p{(\linewidth - 2\tabcolsep) * \real{0.2800}}
  >{\raggedright\arraybackslash}p{(\linewidth - 2\tabcolsep) * \real{0.7200}}@{}}
\toprule\noalign{}
\endhead
\bottomrule\noalign{}
\endlastfoot
ASSEMBLER CALL: & OS9 F\$AllTsk \\
MACHINE CODE: & 103F 3F \\
INPUT: & (X) = Process Descriptor pointer \\
OUTPUT: & None. \\
ERROR OUTPUT: & \begin{minipage}[t]{\linewidth}\raggedright
(CC)~=~C~bit~set.\\
(B)~=~Appropriate~error~code.
\end{minipage} \\
\end{longtable}
}

Allocates a Task number for the given process.

NOTE: THIS IS A PRIVILEGED SYSTEM MODE SERVICE REQUEST

\subsection{F\$Boot - Bootstrap system}

{\def\LTcaptype{none} % do not increment counter
\begin{longtable}[]{@{}
  >{\raggedright\arraybackslash}p{(\linewidth - 2\tabcolsep) * \real{0.2800}}
  >{\raggedright\arraybackslash}p{(\linewidth - 2\tabcolsep) * \real{0.7200}}@{}}
\toprule\noalign{}
\endhead
\bottomrule\noalign{}
\endlastfoot
ASSEMBLER CALL: & OS9 F\$Boot \\
MACHINE CODE: & 103F 35 \\
INPUT: & none \\
OUTPUT: & none \\
ERROR OUTPUT: & \begin{minipage}[t]{\linewidth}\raggedright
(CC)~=~C~bit~set.\\
(B)~=~Appropriate~error~code.
\end{minipage} \\
\end{longtable}
}

Links to the module named ``Boot'' or as specified in the INIT module;
calls linked module; and expects the return of a pointer and size of an
area which is then searched for new modules.

This is a privileged system mode service request.

\subsection{F\$BtMem - Bootstrap Memory request}

{\def\LTcaptype{none} % do not increment counter
\begin{longtable}[]{@{}
  >{\raggedright\arraybackslash}p{(\linewidth - 2\tabcolsep) * \real{0.2800}}
  >{\raggedright\arraybackslash}p{(\linewidth - 2\tabcolsep) * \real{0.7200}}@{}}
\toprule\noalign{}
\endhead
\bottomrule\noalign{}
\endlastfoot
ASSEMBLER CALL: & OS9 F\$BtMem \\
MACHINE CODE: & 103F 36 \\
INPUT: & (D) = Byte count requested. \\
OUTPUT: & \begin{minipage}[t]{\linewidth}\raggedright
(D)~=~Byte~count~granted.\\
(U)~=~Pointer~to~memory~allocated.
\end{minipage} \\
ERROR OUTPUT: & \begin{minipage}[t]{\linewidth}\raggedright
(CC)~=~C~bit~set.\\
(B)~=~Appropriate~error~code.
\end{minipage} \\
\end{longtable}
}

Allocates requested memory (rounded up to nearest block) as contigous
memory in the system\textquotesingle s address space.

NOTE: THIS IS A PRIVILEGED SYSTEM MODE SERVICE REQUEST

\subsection{F\$ClrBlk - Clear specific Block}

{\def\LTcaptype{none} % do not increment counter
\begin{longtable}[]{@{}
  >{\raggedright\arraybackslash}p{(\linewidth - 2\tabcolsep) * \real{0.2800}}
  >{\raggedright\arraybackslash}p{(\linewidth - 2\tabcolsep) * \real{0.7200}}@{}}
\toprule\noalign{}
\endhead
\bottomrule\noalign{}
\endlastfoot
ASSEMBLER CALL: & OS9 F\$ClrBlk \\
MACHINE CODE: & 103F 50 \\
INPUT: & \begin{minipage}[t]{\linewidth}\raggedright
(B)~=~Number~of~blocks\\
(U)~=~Address~of~first~block
\end{minipage} \\
OUTPUT: & None. \\
ERROR OUTPUT: & None. \\
\end{longtable}
}

Marks blocks in process DAT image as unallocated.

NOTE: THIS IS A PRIVILEGED SYSTEM MODE SERVICE REQUEST

\subsection{F\$CpyMem - Copy external Memory}

{\def\LTcaptype{none} % do not increment counter
\begin{longtable}[]{@{}
  >{\raggedright\arraybackslash}p{(\linewidth - 2\tabcolsep) * \real{0.2800}}
  >{\raggedright\arraybackslash}p{(\linewidth - 2\tabcolsep) * \real{0.7200}}@{}}
\toprule\noalign{}
\endhead
\bottomrule\noalign{}
\endlastfoot
ASSEMBLER CALL: & OS9 F\$CpyMem \\
MACHINE CODE: & 103F 1B \\
INPUT: & \begin{minipage}[t]{\linewidth}\raggedright
(D)~=~Starting~Memory~Block~number\\
(X)~=~Offset~in~block~to~begin~copy\\
(Y)~=~Byte~count\\
(U)~=~Caller\textquotesingle s~destination~buffer
\end{minipage} \\
OUTPUT: & None. \\
ERROR OUTPUT: & \begin{minipage}[t]{\linewidth}\raggedright
(CC)~=~C~bit~set.\\
(B)~=~Appropriate~error~code.
\end{minipage} \\
\end{longtable}
}

Reads external memory into the user\textquotesingle s buffer for
inspection. Any memory in the system may be viewed in this way.

\subsection{F\$DATLog - Convert DAT block/offset to Logical Addr}

{\def\LTcaptype{none} % do not increment counter
\begin{longtable}[]{@{}
  >{\raggedright\arraybackslash}p{(\linewidth - 2\tabcolsep) * \real{0.2800}}
  >{\raggedright\arraybackslash}p{(\linewidth - 2\tabcolsep) * \real{0.7200}}@{}}
\toprule\noalign{}
\endhead
\bottomrule\noalign{}
\endlastfoot
ASSEMBLER CALL: & OS9 F\$DATLog \\
MACHINE CODE: & 103F 44 \\
INPUT: & \begin{minipage}[t]{\linewidth}\raggedright
(B)~=~DAT~image~offset\\
(X)~=~Block~offset
\end{minipage} \\
OUTPUT: & (X) = Logical address. \\
ERROR OUTPUT: & \begin{minipage}[t]{\linewidth}\raggedright
(CC)~=~C~bit~set.\\
(B)~=~Appropriate~error~code.
\end{minipage} \\
\end{longtable}
}

Converts a DAT imaqe block number and block offset to its equivalent
logical address.

This is a privileged system mode service request.

\subsection{F\$DATTmp - Make Temporary DAT image}

{\def\LTcaptype{none} % do not increment counter
\begin{longtable}[]{@{}
  >{\raggedright\arraybackslash}p{(\linewidth - 2\tabcolsep) * \real{0.2800}}
  >{\raggedright\arraybackslash}p{(\linewidth - 2\tabcolsep) * \real{0.7200}}@{}}
\toprule\noalign{}
\endhead
\bottomrule\noalign{}
\endlastfoot
ASSEMBLER CALL: & OS9 F\$DATTmp \\
MACHINE CODE: & 103F 45 \\
INPUT: & (D) = Block number \\
OUTPUT: & (Y) = DAT image pointer \\
ERROR OUTPUT: & \begin{minipage}[t]{\linewidth}\raggedright
(CC)~=~C~bit~set.\\
(B)~=~Appropriate~error~code.
\end{minipage} \\
\end{longtable}
}

Builds a temporary DAT image to access the given memory block.

NOTE: THIS IS A PRIVILEGED SYSTEM MODE SERVICE REQUEST

\subsection{F\$DelImg - Deallocate Image RAM blocks}

{\def\LTcaptype{none} % do not increment counter
\begin{longtable}[]{@{}
  >{\raggedright\arraybackslash}p{(\linewidth - 2\tabcolsep) * \real{0.2800}}
  >{\raggedright\arraybackslash}p{(\linewidth - 2\tabcolsep) * \real{0.7200}}@{}}
\toprule\noalign{}
\endhead
\bottomrule\noalign{}
\endlastfoot
ASSEMBLER CALL: & OS9 F\$DelImg \\
MACHINE CODE: & 103F 3B \\
INPUT: & \begin{minipage}[t]{\linewidth}\raggedright
(A)~=~Beginning~block~number\\
(B)~=~Block~count\\
(X)~=~Process~Descriptor~pointer
\end{minipage} \\
OUTPUT: & None. \\
ERROR OUTPUT: & \begin{minipage}[t]{\linewidth}\raggedright
(CC)~=~C~bit~set.\\
(B)~=~Appropriate~error~code.
\end{minipage} \\
\end{longtable}
}

Deallocated memory from the process\textquotesingle{} address space.

This is a privileged system mode service request.

\subsection{F\$DelPrc - Deallocate Process descriptor}

{\def\LTcaptype{none} % do not increment counter
\begin{longtable}[]{@{}
  >{\raggedright\arraybackslash}p{(\linewidth - 2\tabcolsep) * \real{0.2800}}
  >{\raggedright\arraybackslash}p{(\linewidth - 2\tabcolsep) * \real{0.7200}}@{}}
\toprule\noalign{}
\endhead
\bottomrule\noalign{}
\endlastfoot
ASSEMBLER CALL: & OS9 F\$DelPrc \\
MACHINE CODE: & 103F 4C \\
INPUT: & (A) = Process ID. \\
OUTPUT: & None. \\
ERROR OUTPUT: & \begin{minipage}[t]{\linewidth}\raggedright
(CC)~=~C~bit~set.\\
(B)~=~Appropriate~error~code.
\end{minipage} \\
\end{longtable}
}

Returns process descriptor memory to system free memory pool.

This is a privileged system mode service request.

\subsection{F\$DelRam - Deallocate RAM blocks}

{\def\LTcaptype{none} % do not increment counter
\begin{longtable}[]{@{}
  >{\raggedright\arraybackslash}p{(\linewidth - 2\tabcolsep) * \real{0.2800}}
  >{\raggedright\arraybackslash}p{(\linewidth - 2\tabcolsep) * \real{0.7200}}@{}}
\toprule\noalign{}
\endhead
\bottomrule\noalign{}
\endlastfoot
ASSEMBLER CALL: & OS9 F\$DelRam \\
MACHINE CODE: & 103F 51 \\
INPUT: & \begin{minipage}[t]{\linewidth}\raggedright
(B)~=~Number~of~blocks\\
(X)~=~Beginning~block~number.
\end{minipage} \\
OUTPUT: & None. \\
ERROR OUTPUT: & None. \\
\end{longtable}
}

Marks blocks in system memory block map as unallocated.

This is a privileged system mode service request.

\subsection{F\$DelTsk - Deallocate process Task number}

{\def\LTcaptype{none} % do not increment counter
\begin{longtable}[]{@{}
  >{\raggedright\arraybackslash}p{(\linewidth - 2\tabcolsep) * \real{0.2800}}
  >{\raggedright\arraybackslash}p{(\linewidth - 2\tabcolsep) * \real{0.7200}}@{}}
\toprule\noalign{}
\endhead
\bottomrule\noalign{}
\endlastfoot
ASSEMBLER CALL: & OS9 F\$DelTsk \\
MACHINE CODE: & 103F 40 \\
INPUT: & (X) = Process Descriptor pointer \\
OUTPUT: & None. \\
ERROR OUTPUT: & \begin{minipage}[t]{\linewidth}\raggedright
(CC)~=~C~bit~set.\\
(B)~=~Appropriate~error~code.
\end{minipage} \\
\end{longtable}
}

Releases the Task number in use by the process.

This is a privileged system mode service request.

\subsection{F\$ELink - Link using module directory Entry}

{\def\LTcaptype{none} % do not increment counter
\begin{longtable}[]{@{}
  >{\raggedright\arraybackslash}p{(\linewidth - 2\tabcolsep) * \real{0.2800}}
  >{\raggedright\arraybackslash}p{(\linewidth - 2\tabcolsep) * \real{0.7200}}@{}}
\toprule\noalign{}
\endhead
\bottomrule\noalign{}
\endlastfoot
ASSEMBLER CALL: & OS9 F\$ELink \\
MACHINE CODE: & 103F 4D \\
INPUT: & \begin{minipage}[t]{\linewidth}\raggedright
(B)~=~Module~type.\\
(X)~=~Pointer~to~module~directory~entry.
\end{minipage} \\
OUTPUT: & None. \\
ERROR OUTPUT: & \begin{minipage}[t]{\linewidth}\raggedright
(CC)~=~C~bit~set.\\
(B)~=~Appropriate~error~code.
\end{minipage} \\
\end{longtable}
}

Performs a ``Link'' given a pointer to a module directory entry. Note
that this call differs from F\$Link in that a pointer to the module
directory entry is supplied rather than a pointer to a module name.

This is a privileged system mode service request.

\subsection{F\$FModul - Find Module directory entry}

{\def\LTcaptype{none} % do not increment counter
\begin{longtable}[]{@{}
  >{\raggedright\arraybackslash}p{(\linewidth - 2\tabcolsep) * \real{0.2800}}
  >{\raggedright\arraybackslash}p{(\linewidth - 2\tabcolsep) * \real{0.7200}}@{}}
\toprule\noalign{}
\endhead
\bottomrule\noalign{}
\endlastfoot
ASSEMBLER CALL: & OS9 F\$FModul \\
MACHINE CODE: & 103F 4E \\
INPUT: & \begin{minipage}[t]{\linewidth}\raggedright
(A)~=~Module~type.\\
(X)~=~Module~Name~string~pointer.\\
(Y)~=~Name~string~DAT~image~pointer.
\end{minipage} \\
OUTPUT: & \begin{minipage}[t]{\linewidth}\raggedright
(A)~=~Module~Type.\\
(B)~=~Module~Revision.\\
(X)~=~Updated~past~name~string.\\
(U)~=~Module~directory~entry~pointer.
\end{minipage} \\
ERROR OUTPUT: & \begin{minipage}[t]{\linewidth}\raggedright
(CC)~=~C~bit~set.\\
(B)~=~Appropriate~error~code.
\end{minipage} \\
\end{longtable}
}

This call returns a pointer to the module directory entry given the
module name.

This is a privileged system mode service request.

\subsection{F\$FreeHB - Get Free High block}

{\def\LTcaptype{none} % do not increment counter
\begin{longtable}[]{@{}
  >{\raggedright\arraybackslash}p{(\linewidth - 2\tabcolsep) * \real{0.2800}}
  >{\raggedright\arraybackslash}p{(\linewidth - 2\tabcolsep) * \real{0.7200}}@{}}
\toprule\noalign{}
\endhead
\bottomrule\noalign{}
\endlastfoot
ASSEMBLER CALL: & OS9 F\$FreeHB \\
MACHINE CODE: & 103F 3E \\
INPUT: & \begin{minipage}[t]{\linewidth}\raggedright
(B)~=~Block~count\\
(Y)~=~DAT~image~pointer
\end{minipage} \\
OUTPUT: & (A) High block number \\
ERROR OUTPUT: & \begin{minipage}[t]{\linewidth}\raggedright
(CC)~=~C~bit~set.\\
(B)~=~Appropriate~error~code.
\end{minipage} \\
\end{longtable}
}

Searches the DAT image for the highest free block of given size.

This is a privileged system mode service request.

\subsection{F\$FreeLB - Get Free Low block}

{\def\LTcaptype{none} % do not increment counter
\begin{longtable}[]{@{}
  >{\raggedright\arraybackslash}p{(\linewidth - 2\tabcolsep) * \real{0.2800}}
  >{\raggedright\arraybackslash}p{(\linewidth - 2\tabcolsep) * \real{0.7200}}@{}}
\toprule\noalign{}
\endhead
\bottomrule\noalign{}
\endlastfoot
ASSEMBLER CALL: & OS9 F\$FreeLB \\
MACHINE CODE: & 103F 3D \\
INPUT: & \begin{minipage}[t]{\linewidth}\raggedright
(B)~=~Block~count\\
(Y)~=~DAT~image~pointer
\end{minipage} \\
OUTPUT: & (A) = Low block number \\
ERROR OUTPUT: & \begin{minipage}[t]{\linewidth}\raggedright
(CC)~=~C~bit~set.\\
(B)~=~Appropriate~error~code.
\end{minipage} \\
\end{longtable}
}

Searches the DAT image for the lowest free block of given size.

This is a privileged system mode service request.

\subsection{F\$GBlkMp - Get system Block Map copy}

{\def\LTcaptype{none} % do not increment counter
\begin{longtable}[]{@{}
  >{\raggedright\arraybackslash}p{(\linewidth - 2\tabcolsep) * \real{0.2800}}
  >{\raggedright\arraybackslash}p{(\linewidth - 2\tabcolsep) * \real{0.7200}}@{}}
\toprule\noalign{}
\endhead
\bottomrule\noalign{}
\endlastfoot
ASSEMBLER CALL: & OS9 F\$GBlkMp \\
MACHINE CODE: & 103F 19 \\
INPUT: & (X) = 1024 byte buffer pointer \\
OUTPUT: & \begin{minipage}[t]{\linewidth}\raggedright
(D)~=~Number~of~bytes~per~block~(MMU~block~size~dependent).\\
(Y)~=~Size~of~system\textquotesingle s~memory~block~map.
\end{minipage} \\
ERROR OUTPUT: & \begin{minipage}[t]{\linewidth}\raggedright
(CC)~=~C~bit~set.\\
(B)~=~Appropriate~error~code.
\end{minipage} \\
\end{longtable}
}

Copies the system\textquotesingle s memory block map into the
user\textquotesingle s buffer for inspection.

\subsection{F\$GModDr - Get Module Directory copy}

{\def\LTcaptype{none} % do not increment counter
\begin{longtable}[]{@{}
  >{\raggedright\arraybackslash}p{(\linewidth - 2\tabcolsep) * \real{0.2800}}
  >{\raggedright\arraybackslash}p{(\linewidth - 2\tabcolsep) * \real{0.7200}}@{}}
\toprule\noalign{}
\endhead
\bottomrule\noalign{}
\endlastfoot
ASSEMBLER CALL: & OS9 F\$GModDr \\
MACHINE CODE: & 103F 1A \\
INPUT: & (X) = 2048 byte buffer pointer \\
ERROR OUTPUT: & \begin{minipage}[t]{\linewidth}\raggedright
(CC)~=~C~bit~set.\\
(B)~=~Appropriate~error~code.
\end{minipage} \\
\end{longtable}
}

Copies the system\textquotesingle s module directory into the
user\textquotesingle s buffer for inspection.

\subsection{F\$GPrDsc - Get Process Descriptor copy}

{\def\LTcaptype{none} % do not increment counter
\begin{longtable}[]{@{}
  >{\raggedright\arraybackslash}p{(\linewidth - 2\tabcolsep) * \real{0.2800}}
  >{\raggedright\arraybackslash}p{(\linewidth - 2\tabcolsep) * \real{0.7200}}@{}}
\toprule\noalign{}
\endhead
\bottomrule\noalign{}
\endlastfoot
ASSEMBLER CALL: & OS9 F\$GPrDsc \\
MACHINE CODE: & 103F 18 \\
INPUT: & \begin{minipage}[t]{\linewidth}\raggedright
(A)~=~Requested~process~ID.\\
(X)~=~512~byte~buffer~pointer.
\end{minipage} \\
OUTPUT: & None. \\
ERROR OUTPUT: & \begin{minipage}[t]{\linewidth}\raggedright
(CC)~=~C~bit~set.\\
(B)~=~Appropriate~error~code.
\end{minipage} \\
\end{longtable}
}

Copies a process descriptor into the calling process\textquotesingle{}
buffer for inspection There is no way to change data in a process
descriptor.

\subsection{F\$GProcP - Get Process Pointer}

{\def\LTcaptype{none} % do not increment counter
\begin{longtable}[]{@{}
  >{\raggedright\arraybackslash}p{(\linewidth - 2\tabcolsep) * \real{0.2800}}
  >{\raggedright\arraybackslash}p{(\linewidth - 2\tabcolsep) * \real{0.7200}}@{}}
\toprule\noalign{}
\endhead
\bottomrule\noalign{}
\endlastfoot
ASSEMBLER CALL: & OS9 F\$GProcP \\
MACHINE CODE: & 103F 37 \\
INPUT: & (A) = Process ID \\
OUTPUT: & (Y) = Pointer to Process Descriptor \\
ERROR OUTPUT: & \begin{minipage}[t]{\linewidth}\raggedright
(CC)~=~C~bit~set.\\
(B)~=~Appropriate~error~code.
\end{minipage} \\
\end{longtable}
}

Translates a process ID number to the address of its process descriptor
in the system address space.

This is a privileged system mode service request.

\subsection{F\$LDABX - Load A from 0,1 in task B}

{\def\LTcaptype{none} % do not increment counter
\begin{longtable}[]{@{}
  >{\raggedright\arraybackslash}p{(\linewidth - 2\tabcolsep) * \real{0.2800}}
  >{\raggedright\arraybackslash}p{(\linewidth - 2\tabcolsep) * \real{0.7200}}@{}}
\toprule\noalign{}
\endhead
\bottomrule\noalign{}
\endlastfoot
ASSEMBLER CALL: & OS9 F\$LDABX \\
MACHINE CODE: & 103F 49 \\
INPUT: & \begin{minipage}[t]{\linewidth}\raggedright
(B)~=~Task~number\\
(X)~=~Data~pointer
\end{minipage} \\
OUTPUT: & (A) = Data byte at 0,X in task\textquotesingle s address
space \\
ERROR OUTPUT: & \begin{minipage}[t]{\linewidth}\raggedright
(CC)~=~C~bit~set.\\
(B)~=~Appropriate~error~code.
\end{minipage} \\
\end{longtable}
}

One byte is returned from the logical address in (X) in the given
task\textquotesingle s address space. This is typically used to get one
byte from the current process\textquotesingle s memory in a system state
routine.

This is a privileged system mode service request.

\subsection{F\$LDAXY - Load A {[}X, {[}Y{]} {]}}

{\def\LTcaptype{none} % do not increment counter
\begin{longtable}[]{@{}
  >{\raggedright\arraybackslash}p{(\linewidth - 2\tabcolsep) * \real{0.2800}}
  >{\raggedright\arraybackslash}p{(\linewidth - 2\tabcolsep) * \real{0.7200}}@{}}
\toprule\noalign{}
\endhead
\bottomrule\noalign{}
\endlastfoot
ASSEMBLER CALL: & OS9 F\$LDAXY \\
MACHINE CODE: & 103F 46 \\
INPUT: & \begin{minipage}[t]{\linewidth}\raggedright
(X)~=~Block~offset\\
(Y)~=~DAT~image~pointer
\end{minipage} \\
OUTPUT: & (A) = data byte at (X) offset of (Y) \\
ERROR OUTPUT: & \begin{minipage}[t]{\linewidth}\raggedright
(CC)~=~C~bit~set.\\
(B)~=~Appropriate~error~code.
\end{minipage} \\
\end{longtable}
}

Returns one data byte in the memory block specified by the DAT image in
(Y), offset by (X).

This is a privileged system mode service request.

\subsection{F\$LDAXYP - Load A {[}X+, {[}Y{]} {]}}

{\def\LTcaptype{none} % do not increment counter
\begin{longtable}[]{@{}
  >{\raggedright\arraybackslash}p{(\linewidth - 2\tabcolsep) * \real{0.2800}}
  >{\raggedright\arraybackslash}p{(\linewidth - 2\tabcolsep) * \real{0.7200}}@{}}
\toprule\noalign{}
\endhead
\bottomrule\noalign{}
\endlastfoot
ASSEMBLER CALL: & OS9 F\$LDAXYP \\
MACHINE CODE: & 103F 47 \\
INPUT: & \begin{minipage}[t]{\linewidth}\raggedright
(X)~=~Block~offset\\
(Y)~=~DAT~image~pointer
\end{minipage} \\
OUTPUT: & \begin{minipage}[t]{\linewidth}\raggedright
(A)~=~Data~byte~at~(X)~offset~of~(Y)\\
(X)~=~incremented~by~one
\end{minipage} \\
ERROR OUTPUT: & \begin{minipage}[t]{\linewidth}\raggedright
(CC)~=~C~bit~set.\\
(B)~=~Appropriate~error~code.
\end{minipage} \\
\end{longtable}
}

Similar to the assembly instruction ``LDA ,X+'', except that (X) refers
to an offset in the memory block described by the DAT image at (Y).

This is a privileged system mode service request.

\subsection{F\$LDDDXY - Load D {[}D+X, {[}Y{]} {]}}

{\def\LTcaptype{none} % do not increment counter
\begin{longtable}[]{@{}
  >{\raggedright\arraybackslash}p{(\linewidth - 2\tabcolsep) * \real{0.2800}}
  >{\raggedright\arraybackslash}p{(\linewidth - 2\tabcolsep) * \real{0.7200}}@{}}
\toprule\noalign{}
\endhead
\bottomrule\noalign{}
\endlastfoot
ASSEMBLER CALL: & OS9 F\$LDDDXY \\
MACHINE CODE: & 103F 48 \\
INPUT: & \begin{minipage}[t]{\linewidth}\raggedright
(D)~=~Offset~to~offset\\
(X)~=~Offset\\
(Y)~=~DAT~image~pointer
\end{minipage} \\
OUTPUT: & (D) = bytes address by {[}D+X,Y{]} \\
ERROR OUTPUT: & \begin{minipage}[t]{\linewidth}\raggedright
(CC)~=~C~bit~set.\\
(B)~=~Appropriate~error~code.
\end{minipage} \\
\end{longtable}
}

Loads two bytes from the memory block described by the DAT image pointed
to by (Y). The bytes loaded are at the offset (D+X) in the memory block.

This is a privileged system mode service request.

\subsection{F\$MapBlk - Map specific Block}

{\def\LTcaptype{none} % do not increment counter
\begin{longtable}[]{@{}
  >{\raggedright\arraybackslash}p{(\linewidth - 2\tabcolsep) * \real{0.2800}}
  >{\raggedright\arraybackslash}p{(\linewidth - 2\tabcolsep) * \real{0.7200}}@{}}
\toprule\noalign{}
\endhead
\bottomrule\noalign{}
\endlastfoot
ASSEMBLER CALL: & OS9 F\$MapBlk \\
MACHINE CODE: & 103F 4F \\
INPUT: & \begin{minipage}[t]{\linewidth}\raggedright
(B)~=~Number~of~blocks.\\
(X)~=~Beginning~block~number.
\end{minipage} \\
OUTPUT: & (U) = Address of first block. \\
ERROR OUTPUT: & \begin{minipage}[t]{\linewidth}\raggedright
(CC)~=~C~bit~set.\\
(B)~=~Appropriate~error~code.
\end{minipage} \\
\end{longtable}
}

Maps specified block(s) into unallocated blocks of process address
space.

This is a privileged system mode service request.

\subsection{F\$Move - Move data (low bound first)}

{\def\LTcaptype{none} % do not increment counter
\begin{longtable}[]{@{}
  >{\raggedright\arraybackslash}p{(\linewidth - 2\tabcolsep) * \real{0.2800}}
  >{\raggedright\arraybackslash}p{(\linewidth - 2\tabcolsep) * \real{0.7200}}@{}}
\toprule\noalign{}
\endhead
\bottomrule\noalign{}
\endlastfoot
ASSEMBLER CALL: & OS9 F\$Move \\
MACHINE CODE: & 103F 38 \\
INPUT: & \begin{minipage}[t]{\linewidth}\raggedright
(A)~=~Source~Task~number\\
(B)~=~Destination~Task~number\\
(X)~=~Source~pointer\\
(Y)~=~Byte~count\\
(U)~=~Destination~pointer
\end{minipage} \\
OUTPUT: & None. \\
ERROR OUTPUT: & \begin{minipage}[t]{\linewidth}\raggedright
(CC)~=~C~bit~set.\\
(B)~=~Appropriate~error~code.
\end{minipage} \\
\end{longtable}
}

Moves data bytes from one address space to anotber, usually from
System\textquotesingle s to User\textquotesingle s, or vice-versa.

This is a privileged system mode service request.

\subsection{F\$RelTsk - Release Task number}

{\def\LTcaptype{none} % do not increment counter
\begin{longtable}[]{@{}
  >{\raggedright\arraybackslash}p{(\linewidth - 2\tabcolsep) * \real{0.2800}}
  >{\raggedright\arraybackslash}p{(\linewidth - 2\tabcolsep) * \real{0.7200}}@{}}
\toprule\noalign{}
\endhead
\bottomrule\noalign{}
\endlastfoot
ASSEMBLER CALL: & OS9 F\$RelTsk \\
MACHINE CODE: & 103F 43 \\
INPUT: & (B) = Task number \\
OUTPUT: & None. \\
ERROR OUTPUT: & \begin{minipage}[t]{\linewidth}\raggedright
(CC)~=~C~bit~set.\\
(B)~=~Appropriate~error~code.
\end{minipage} \\
\end{longtable}
}

Releases the specified DAT Task number.

This is a privileged system mode service request.

\subsection{F\$ResTsk - Reserve Task number}

DAT task number.

{\def\LTcaptype{none} % do not increment counter
\begin{longtable}[]{@{}
  >{\raggedright\arraybackslash}p{(\linewidth - 2\tabcolsep) * \real{0.2800}}
  >{\raggedright\arraybackslash}p{(\linewidth - 2\tabcolsep) * \real{0.7200}}@{}}
\toprule\noalign{}
\endhead
\bottomrule\noalign{}
\endlastfoot
ASSEMBLER CALL: & OS9 F\$ResTsk \\
MACHINE CODE: & 103F 42 \\
INPUT: & none \\
OUTPUT: & (B) = Task number \\
ERROR OUTPUT: & \begin{minipage}[t]{\linewidth}\raggedright
(CC)~=~C~bit~set.\\
(B)~=~Appropriate~error~code.
\end{minipage} \\
\end{longtable}
}

Finds a free DAT task number.

This is a privileged system mode service request.

\subsection{F\$SetImg - Set process DAT Image}

{\def\LTcaptype{none} % do not increment counter
\begin{longtable}[]{@{}
  >{\raggedright\arraybackslash}p{(\linewidth - 2\tabcolsep) * \real{0.2800}}
  >{\raggedright\arraybackslash}p{(\linewidth - 2\tabcolsep) * \real{0.7200}}@{}}
\toprule\noalign{}
\endhead
\bottomrule\noalign{}
\endlastfoot
ASSEMBLER CALL: & OS9 F\$SetImg \\
MACHINE CODE: & 103F 3C \\
INPUT: & \begin{minipage}[t]{\linewidth}\raggedright
(A)~=~Beginning~image~block~number\\
(B)~=~Block~count\\
(X)~=~Process~Descriptor~pointer\\
(U)~New~image~pointer
\end{minipage} \\
OUTPUT: & None. \\
ERROR OUTPUT: & \begin{minipage}[t]{\linewidth}\raggedright
(CC)~=~C~bit~set.\\
(B)~=~Appropriate~error~code.
\end{minipage} \\
\end{longtable}
}

Copies a DAT image into the process descriptor.

NOTE: THIS IS A PRIVILEGED SYSTEM MODE SERVICE REQUEST

\subsection{F\$SetTsk - Set process Task DAT registers}

{\def\LTcaptype{none} % do not increment counter
\begin{longtable}[]{@{}
  >{\raggedright\arraybackslash}p{(\linewidth - 2\tabcolsep) * \real{0.2800}}
  >{\raggedright\arraybackslash}p{(\linewidth - 2\tabcolsep) * \real{0.7200}}@{}}
\toprule\noalign{}
\endhead
\bottomrule\noalign{}
\endlastfoot
ASSEMBLER CALL: & OS9 F\$SetTsk \\
MACHINE CODE: & 103F 41 \\
INPUT: & (X) = Process Descriptor pointer \\
OUTPUT: & None. \\
ERROR OUTPUT: & \begin{minipage}[t]{\linewidth}\raggedright
(CC)~=~C~bit~set.\\
(B)~=~Appropriate~error~code.
\end{minipage} \\
\end{longtable}
}

Sets the process Task DAT registers.

This is a privileged system mode service request.

\subsection{F\$SLink - System Link}

{\def\LTcaptype{none} % do not increment counter
\begin{longtable}[]{@{}
  >{\raggedright\arraybackslash}p{(\linewidth - 2\tabcolsep) * \real{0.2800}}
  >{\raggedright\arraybackslash}p{(\linewidth - 2\tabcolsep) * \real{0.7200}}@{}}
\toprule\noalign{}
\endhead
\bottomrule\noalign{}
\endlastfoot
ASSEMBLER CALL: & OS9 F\$SLink \\
MACHINE CODE: & 103F 34 \\
INPUT: & \begin{minipage}[t]{\linewidth}\raggedright
(A)~=~Module~Type.\\
(X)~=~Module~Name~string~pointer.\\
(Y)~=~Name~string~DAT~image~pointer.
\end{minipage} \\
OUTPUT: & \begin{minipage}[t]{\linewidth}\raggedright
(A)~=~Module~Type.\\
(B)~=~Module~Revision.\\
(X)~=~Updated~Name~string~pointer.\\
(Y)~=~Module~Entry~point.\\
(U)~=~Module~pointer.
\end{minipage} \\
ERROR OUTPUT: & \begin{minipage}[t]{\linewidth}\raggedright
(CC)~=~C~bit~set.\\
(B)~=~Appropriate~error~code.
\end{minipage} \\
\end{longtable}
}

Links a module whose name is outside the current (system)
process\textquotesingle{} adress space into the Address space that
contains its name.

This is a privileged system mode service request.

\subsection{F\$SRqMem - System Memory Request}

{\def\LTcaptype{none} % do not increment counter
\begin{longtable}[]{@{}
  >{\raggedright\arraybackslash}p{(\linewidth - 2\tabcolsep) * \real{0.2800}}
  >{\raggedright\arraybackslash}p{(\linewidth - 2\tabcolsep) * \real{0.7200}}@{}}
\toprule\noalign{}
\endhead
\bottomrule\noalign{}
\endlastfoot
ASSEMBLER CALL: & OS9 F\$SRqMem \\
MACHINE CODE: & 103F 28 \\
INPUT: & (D) = byte count of requested memory \\
OUTPUT: & \begin{minipage}[t]{\linewidth}\raggedright
(D)~=~byte~count~of~memory~granted\\
(U)~=~pointer~to~memory~block~allocated
\end{minipage} \\
ERROR OUTPUT: & \begin{minipage}[t]{\linewidth}\raggedright
(CC)~=~C~bit~set.\\
(B)~=~Appropriate~error~code.
\end{minipage} \\
\end{longtable}
}

Allocates the requested memory (rounded up to the nearest page) in the
system\textquotesingle s address space. Useful for allocating I/O
buffers and other semi-permanent system memory.

This is a privileged system mode service request.

\subsection{F\$SRtMem - System Memory Return}\label{f.srtmem}

{\def\LTcaptype{none} % do not increment counter
\begin{longtable}[]{@{}
  >{\raggedright\arraybackslash}p{(\linewidth - 2\tabcolsep) * \real{0.2800}}
  >{\raggedright\arraybackslash}p{(\linewidth - 2\tabcolsep) * \real{0.7200}}@{}}
\toprule\noalign{}
\endhead
\bottomrule\noalign{}
\endlastfoot
ASSEMBLER CALL: & OS9 F\$SRtMem \\
MACHINE CODE: & 103F 29 \\
INPUT: & \begin{minipage}[t]{\linewidth}\raggedright
(D)~=~Byte~count~of~memory~being~returned\\
(U)~=~Address~of~memory~block~being~returned
\end{minipage} \\
ERROR OUTPUT: & \begin{minipage}[t]{\linewidth}\raggedright
(CC)~=~C~bit~set.\\
(B)~=~Appropriate~error~code.
\end{minipage} \\
\end{longtable}
}

Returns system memory (e.g., memory in the system address space) after
it is no longer needed.

This is a privileged system mode service request.

\subsection{F\$STABX - Store A at 0,X in task B}

{\def\LTcaptype{none} % do not increment counter
\begin{longtable}[]{@{}
  >{\raggedright\arraybackslash}p{(\linewidth - 2\tabcolsep) * \real{0.2800}}
  >{\raggedright\arraybackslash}p{(\linewidth - 2\tabcolsep) * \real{0.7200}}@{}}
\toprule\noalign{}
\endhead
\bottomrule\noalign{}
\endlastfoot
ASSEMBLER CALL: & OS9 F\$STABX \\
MACHINE CODE: & 103F 4A \\
INPUT: & \begin{minipage}[t]{\linewidth}\raggedright
(A)~=~Data~byte~to~store~in~Task\textquotesingle s~address~space\\
(B)~=~Task~number\\
(X)~=~Logical~address~in~task\textquotesingle s~address~space~to~store
\end{minipage} \\
OUTPUT: & None. \\
ERROR OUTPUT: & \begin{minipage}[t]{\linewidth}\raggedright
(CC)~=~C~bit~set.\\
(B)~=~Appropriate~error~code.
\end{minipage} \\
\end{longtable}
}

This is similar to the assembly instruction ``STA 0,X'', except that (X)
refers to an address in the given task\textquotesingle s address space
rather than the current address space.

This is a privileged system mode service request.

\subsection{F\$SUser Set User ID number}

{\def\LTcaptype{none} % do not increment counter
\begin{longtable}[]{@{}
  >{\raggedright\arraybackslash}p{(\linewidth - 2\tabcolsep) * \real{0.2800}}
  >{\raggedright\arraybackslash}p{(\linewidth - 2\tabcolsep) * \real{0.7200}}@{}}
\toprule\noalign{}
\endhead
\bottomrule\noalign{}
\endlastfoot
ASSEMBLER CALL: & OS9 F\$SUser \\
MACHINE CODE: & 103F 1C \\
INPUT: & (Y) = desired User ID number \\
OUTPUT: & None \\
ERROR OUTPUT: & \begin{minipage}[t]{\linewidth}\raggedright
(CC)~=~C~bit~set.\\
(B)~=~Appropriate~error~code.
\end{minipage} \\
\end{longtable}
}

Alters the current user ID to that specified, without error checking.

\subsection{F\$UnLoad - Unlink module by name}

{\def\LTcaptype{none} % do not increment counter
\begin{longtable}[]{@{}
  >{\raggedright\arraybackslash}p{(\linewidth - 2\tabcolsep) * \real{0.2800}}
  >{\raggedright\arraybackslash}p{(\linewidth - 2\tabcolsep) * \real{0.7200}}@{}}
\toprule\noalign{}
\endhead
\bottomrule\noalign{}
\endlastfoot
ASSEMBLER CALL: & OS9 F\$UnLoad \\
MACHINE CODE: & 103F 1D \\
INPUT: & \begin{minipage}[t]{\linewidth}\raggedright
(A)~=~Module~Type\\
(X)~=~Module~Name~pointer
\end{minipage} \\
OUTPUT: & None \\
ERROR OUTPUT: & \begin{minipage}[t]{\linewidth}\raggedright
(CC)~=~C~bit~set.\\
(B)~=~Appropriate~error~code.
\end{minipage} \\
\end{longtable}
}

Locates the module to the module directory, decrements its link count
and removes it from the directory if the count reaches zero. Note that
this call differs from F\$UnLink in that the a pointer to module name is
supplied rather than the address of the module header.

\subsection{I\$DeletX - Delete a file}

{\def\LTcaptype{none} % do not increment counter
\begin{longtable}[]{@{}
  >{\raggedright\arraybackslash}p{(\linewidth - 2\tabcolsep) * \real{0.2800}}
  >{\raggedright\arraybackslash}p{(\linewidth - 2\tabcolsep) * \real{0.7200}}@{}}
\toprule\noalign{}
\endhead
\bottomrule\noalign{}
\endlastfoot
ASSEMBLER CALL: & OS9 I\$Deletx \\
MACHINE CODE: & 103F 90 \\
INPUT: & \begin{minipage}[t]{\linewidth}\raggedright
(X)~=~Address~of~pathlist\\
(A)~=~Access~mode.
\end{minipage} \\
OUTPUT: & (X) = Updated past pathlist (trailing spaces skipped). \\
ERROR OUTPUT: & \begin{minipage}[t]{\linewidth}\raggedright
(CC)~=~C~bit~set.\\
(B)~=~Appropriate~error~code.
\end{minipage} \\
\end{longtable}
}

This service request deletes the file specified by the pathlist. The
file must have write permission attributes (public write if not the
owner), and reside on a multi-file mass storage device. Attempts to
delete devices will result in error.

The access mode is used to specify the current working directory or me
current execution directory (but not both) in the absence of a full
pathlist. If the access mode is read, write, or update, the current data
directory is assumed. If the access mode is execute, the current
execution directory is assumed. Note that if a full pathlist (a pathlist
beginning with a ``/'') appears, the access mode is ignored.

ACCESS MODES:

1~=~Read\\
2~=~Write\\
3~=~Update~(read~or~write)\\
4~=~Execute\\

\end{document}
